\documentclass[11pt,letterpaper]{article}

\usepackage[T1]{fontenc}
\usepackage[utf8]{inputenc}
\usepackage{lmodern}
\usepackage[margin=0.82in]{geometry}
\usepackage{microtype}
\usepackage{amsmath,amssymb,amsthm,mathtools}
\usepackage{amscd}
\usepackage{tikz-cd}
\usepackage{booktabs,array}
\usepackage{enumitem}
\usepackage{natbib}
\usepackage{xurl}
\usepackage{xcolor}
\usepackage[colorlinks=true,linkcolor=blue!45!black,
 citecolor=blue!45!black,urlcolor=blue!45!black]{hyperref}
\usepackage[nameinlink,noabbrev]{cleveref}

\setlength{\parindent}{0pt}
\setlength{\parskip}{0.48em}
\numberwithin{equation}{section}
\newtheorem{theorem}{Theorem}[section]
\newtheorem{proposition}[theorem]{Proposition}
\newtheorem{lemma}[theorem]{Lemma}
\newtheorem{corollary}[theorem]{Corollary}
\theoremstyle{definition}
\newtheorem{definition}[theorem]{Definition}
\theoremstyle{remark}
\newtheorem{remark}[theorem]{Remark}

\newcommand{\R}{\mathbb R}
\newcommand{\C}{\mathbb C}
\newcommand{\Z}{\mathbb Z}
\newcommand{\T}{\mathbb T}
\newcommand{\id}{\mathrm{id}}
\newcommand{\sys}{\operatorname{sys}}
\newcommand{\Cone}{\operatorname{Cone}}
\newcommand{\Susp}{\operatorname{Susp}}
\newcommand{\im}{\operatorname{im}}
\newcommand{\rank}{\operatorname{rank}}

\title{Real \(S^{12}\) and \(S^{24}\) Routes from Alp\"oge's Candidate
\((3,4,\infty)\) Two-Torus Fibration\\
\large Exact sphere carriers, four-dimensional cusp limits, and their sharp scope}
\author{A self-contained mathematical note\\
\small based on the candidate $S^6$ construction of Levent Alp\"oge}
\date{5 September 2026}

\begin{document}
\maketitle

\begin{abstract}
Starting from the geometric input of Alp\"oge's candidate
\((3,4,\infty)\) complex two-torus fibration---reconstructed and separately
verified in the accompanying workbench, but imported here as stated
input---we derive the requested
\(S^{12}\)- and \(S^{24}\)-based four-dimensional decompactification
routes directly from its native \(T^2\times\R^2\) cusp geometry.  The
cusp-local rank-two quotient in two copies and the unique
  full-monodromy-invariant rank-one quotient in four copies are optimal among
  native connected torus quotients with unscaled Euclidean quotient metric.
  The two-copy route discards the bounded \(T^2\) in each original fibre and
  combines the two surviving rank-two quotients; it does not retain all four
  directions of one fibre.  After a based choice
  \(X\cong_{\mathrm{diff}}S^6\), they give continuous
surjective quotient maps
\[
 S^{12}\twoheadrightarrow S^6,\qquad
 S^{24}\twoheadrightarrow S^8.
\]
The four-copy map \(S^{24}\to S^8\) is nevertheless based
null-homotopic: more generally, if
\(f=\widehat q_\gamma:S^6\to S^2\), then
\([f^{\wedge r}]=0\in\pi_{6r}(S^{2r})\) for every \(r\geq2\).
The distinct two-copy map \(S^{12}\to S^6\) is also based
null-homotopic: more strongly, Toda's external-smash formula and the exact
low-stem groups show that \(x\wedge y:S^{12}\to S^6\) is based
null-homotopic for every pair of based maps \(x,y:S^6\to S^3\).
Neither conclusion computes either one-copy sphere-map class, and neither
implies nullity of a map into one of the earlier conically stratified
targets before the final sphere-valued postcomposition.
For the specified \(S^6\to S^3\) map we now compute that datum:
its complete regular inverse image and actual derivative framing bound
an explicit collared \(D^2\times T^2\) inside one original cusp chart
times an interval. The full holomorphic period correction is retained
by an exact invertible chart map. Thus the one-copy map itself is
based null-homotopic, as are all its Cartesian powers. The general
framing-torsor calculation remains valid, but its actual coefficient
for this map is zero. The separate one-copy map to \(S^2\) remains open.
Retaining the angle before the final sphere quotient instead gives
\(H=p_3q_2:X\to S^4\) with \([H]=\eta_4\eta_5\ne0\).
We compute its entire normally four-framed inverse torus and its
Arf invariant one, proving that the original conical map \(q_2\)
is nonnull. The ordered Hopf smash \(\eta_2\wedge H:S^9\to S^6\)
represents \(12\nu_6\ne0\) in \(\mathbb Z/24\), and its
threefold suspension represents \(12\nu_9\) in \(\pi_{12}(S^9)\).
Its exact peripheral character has values \((12,0,12)\) and kernel
\(\chi^{-1}(2\mathbb Z/12)\). These statements use the specific
retained-angle map, not the null angle-forgetting postcomposition.
Their targets contain moving embedded metric four-spheres whose
radius-\(R\) balls are exactly Euclidean for all sufficiently deep cusp
parameters.  The same radial quotients also form explicit Hausdorff
second-countable five-manifolds with boundary \(\R^4\).  Each admits a
continuous surjection to a compact cone on \(T^4\), homeomorphic away from
that boundary and constant on it.  These comparison maps are nonproper and
are not quotient maps; in particular they are not topological blow-downs or
resolutions, and the compact cone does not carry the topology obtained by
collapsing the boundary \(\R^4\).

We prove the point-set quotient statements, optimality, target homotopy
types, local homology, and metric assertions in full.  For one fixed
cusp-adapted metric on each compact target, the four-dimensional slices
appear as vertical \(\R^4\) subspaces of moving-base-point bubbles; the
full bubbles are \(\R^8\) for the \(S^{12}\) route and \(\R^{12}\) for
the \(S^{24}\) routes.  Any unrescaled fixed compatible target metric
instead collapses the slices to the cone vertex, and no metric on a source
sphere is produced.  A separate finite-cover construction also
decompactifies the original four-tori to \(\R^4\), but it is ancillary:
it neither uses nor replaces the \(S^{12}/S^{24}\) sphere-carrier
construction studied here.  We also give an explicit \(SU(3)\) clutching
calculation: nonzero second-Chern bundles on the moving \(T^4\) slices
cannot extend over the contractible cone carrier or over a null ambient
four-sphere, while pullback along either null sphere map is trivial.  These
facts are then compared with the currently separately constructed
higher-rung Gram quotient
\(M_3(\C)/U(3)\cong\operatorname{Herm}_3^+\), whose one-point
compactification is a nine-ball with rank-deficient boundary an \(S^8\).
The smash-map, product, join, bundle, almost-complex, and fixed-metric
obstructions do not remove that topological carrier; they do rule out
the exact stronger identifications to which they apply.  The chosen based
homeomorphism of the two abstract \(S^8\)'s and the resulting null composite
are the strongest bridge proved here.  Constructing a map that also respects
rank strata, metrics, or bundle data remains an active bridge problem, not
a claim of ontological separation.  Nothing below constructs a complex
\(S^{12}\) or \(S^{24}\), an ambient four-dimensional compact target, or
a physical spacetime.  For \(S^8\), the almost-complex obstruction follows
from the \(p_2\)-formula and the Hirzebruch signature theorem (or, for the
standard sphere, stable tangent triviality).  For \(S^{12}\), the Borel--Serre
mod-\(5\) reduced-power calculation, Bott integrality, and the
stable-tangent \(p_3\)-identity give three independent obstructions.
\end{abstract}

\paragraph{Status and provenance.}
The compact complex threefold $X$, its torus fibration, its period functions,
its three local fillings, and the recognition of its underlying smooth
manifold as $S^6$ belong to the unpublished candidate construction of
Levent Alp\"oge \citep{AlpogeS6_2026}.  They are the geometric input of this
note, have been reconstructed and separately verified in the accompanying
workbench, and are restated below in the exact form used rather than reproved
here.  The quotient, smash, join, framing, clutching,
homology, metric, optimality, and obstruction statements are proved here
from that displayed input.  The symbols \(S^{12}\) and \(S^{24}\) refer
to based homeomorphism types of smash powers of the underlying real
\(S^6\), not to Cartesian powers and not to new complex manifolds.  The
central positive result is an exact four-dimensional limit carried by
moving target slices and by a noncompact decompactification boundary.  It
is not an assertion that either compact source sphere, or either fixed
compact quotient target, becomes \(\R^4\).

\section{The exact input data}

\subsection{The candidate fibration and its periods}

Let $B\cong\mathbb P^1$ have three distinguished values
$p_0,p_1,p_2$, with $p_0$ the cusp and with finite local orders
$m_1=3$ and $m_2=4$.  The candidate construction supplies a connected
compact complex threefold and a proper surjective holomorphic map
\begin{equation}
 f:X\longrightarrow B
 \label{eq:fibration}
\end{equation}
with connected fibres.  Its cusp fibre $W=f^{-1}(p_0)$ is reduced,
irreducible, nonnormal, and toroidal; the fibres at $p_1,p_2$ have
scheme-theoretic multiplicities $3,4$.  The same construction identifies
the underlying based smooth manifold with the standard sphere:
\begin{equation}
 (X,x_0)\cong_{\mathrm{diff}}(S^6,*),\qquad x_0\in W.
 \label{eq:source-sphere}
\end{equation}
Only this real based diffeomorphism is used below.

Put $B^\circ=B\setminus\{p_0,p_1,p_2\}$ and
$J=f^{-1}(B^\circ)$.  On the marked universal cover of $B^\circ$ take the
ordered lattices
\begin{equation}
 V=\Z\langle\gamma,u,w,\delta\rangle,
 \qquad
 \Lambda=V^*=\Z\langle\widehat\gamma,\widehat u,
                              \widehat w,\widehat\delta\rangle.
 \label{eq:lattices}
\end{equation}
There are holomorphic functions $\tau,\mu,\beta$ for which
\begin{equation}
 \Pi=[Z\mid I_2],\qquad
 Z=\begin{pmatrix}6\mu&\tau\\ \beta&\mu\end{pmatrix}
 \label{eq:period-matrix}
\end{equation}
maps $\Lambda$ isomorphically to a lattice in $\C^2$.  Thus a marked
regular fibre is $\C^2/\Pi(z)\Lambda$.  The nondegeneracy inequalities are
\begin{equation}
 \Im\tau>0,
 \qquad
 \Im\beta-\frac{6(\Im\mu)^2}{\Im\tau}<0.
 \label{eq:period-nondegeneracy}
\end{equation}
The equivariant descent fixes the zero section.  Consequently $J$ is a
holomorphic family of complex Lie groups, rather than merely an affine
torus bundle.

Choose a cusp coordinate $t$ with $t(p_0)=0$.  The canonical cusp
calculation gives functions $\mu(t),b(t),h(t)$ holomorphic on a full disc,
where
\begin{equation}
 b=\beta+\tau,\qquad s=\tau-h,\qquad e^{2\pi i s}=t,
 \label{eq:cusp-coordinates}
\end{equation}
and the exact identity is
\begin{equation}
 Z=s\begin{pmatrix}0&1\\-1&0\end{pmatrix}
   +\begin{pmatrix}6\mu&h\\b-h&\mu\end{pmatrix}.
 \label{eq:cusp-period-identity}
\end{equation}
On the positive radial ray set $t=e^{-2\pi T}$ and choose $s=iT$.  Write
\begin{equation}
 \begin{aligned}
 \mu_T&=\alpha_T+i m_T,&
 h_T&=\eta_T+iH_T,& b_T&=\xi_T+iB_T,\\
 c_T&=\xi_T-\eta_T,&
 L_T&=T+H_T,&q_T&=L_T-B_T,\\
 D_T&=L_Tq_T+6m_T^2.&&
 \end{aligned}
 \label{eq:radial-notation}
\end{equation}
If $m=\Im\mu(0)$, $H=\Im h(0)$, and $B=\Im b(0)$, holomorphy gives
\begin{equation}
 m_T=m+O(e^{-2\pi T}),\quad
 H_T=H+O(e^{-2\pi T}),\quad
 B_T=B+O(e^{-2\pi T}).
 \label{eq:radial-errors}
\end{equation}

Identify $\C^2$ with $\R^4$ in the order
$(x_1,y_1,x_2,y_2)$.  The exact real period matrix is
\begin{equation}
 P_T=
 \begin{pmatrix}
 6\alpha_T&\eta_T&1&0\\
 6m_T&L_T&0&0\\
 c_T&\alpha_T&0&1\\
 -q_T&m_T&0&0
 \end{pmatrix},
 \qquad \det P_T=-D_T,\qquad D_T>0.
 \label{eq:real-period-matrix}
\end{equation}
Equip
\begin{equation}
 F_T=\R^4/P_T\Z^4
 \label{eq:radial-fibre}
\end{equation}
with the quotient of the coordinate Euclidean metric.  This is a metric
chosen on the marked radial fibres; it is not asserted to extend as one
canonical global metric on $X$.

\subsection{Integral monodromy and its invariant flag}

In the ordered basis of $\Lambda$, the three monodromies are
\begin{equation}
 A_1=
 \begin{pmatrix}
 1&0&0&0\\6&0&1&0\\-6&-1&-1&0\\-2&1&0&1
 \end{pmatrix},\quad
 A_2=
 \begin{pmatrix}
 1&0&0&0\\0&0&-1&0\\-6&1&0&0\\3&0&1&1
 \end{pmatrix},\quad
 M_0=
 \begin{pmatrix}
 1&0&0&0\\0&1&0&0\\0&1&1&0\\-1&0&0&1
 \end{pmatrix}.
 \label{eq:monodromy-matrices}
\end{equation}
They satisfy $A_1^3=A_2^4=I$ and $A_1A_2M_0=I$.  Direct multiplication
gives
\begin{equation}
 (M_0-I)(a,b,c,d)^{\mathsf T}=(0,0,b,-a)^{\mathsf T}.
 \label{eq:cusp-monodromy-action}
\end{equation}
Hence the saturated cusp lattice
\begin{equation}
 K_\Lambda:=\ker(M_0-I)=\im(M_0-I)
 =\Z\langle\widehat w,\widehat\delta\rangle
 \label{eq:cusp-lattice}
\end{equation}
is both the invariant and vanishing image lattice for the unipotent cusp
operator.

Define the primitive covector and its kernel by
\begin{equation}
 \varepsilon_\gamma
 (a\widehat\gamma+b\widehat u+c\widehat w+d\widehat\delta)=a,
 \qquad
 L_\Lambda:=\ker\varepsilon_\gamma
 =\Z\langle\widehat u,\widehat w,\widehat\delta\rangle.
 \label{eq:invariant-character}
\end{equation}
The first row of each matrix in \cref{eq:monodromy-matrices} shows
$\varepsilon_\gamma A_j=\varepsilon_\gamma$ and
$\varepsilon_\gamma M_0=\varepsilon_\gamma$.  Moreover
\begin{equation}
 \begin{aligned}
 A_1\widehat u&=-\widehat w+\widehat\delta,&
 A_1\widehat w&=\widehat u-\widehat w,&
 A_1\widehat\delta&=\widehat\delta,\\
 A_2\widehat u&=\widehat w,&
 A_2\widehat w&=-\widehat u+\widehat\delta,&
 A_2\widehat\delta&=\widehat\delta.
 \end{aligned}
 \label{eq:invariant-flag}
\end{equation}
Thus $L_\Lambda$ is invariant and is the smallest invariant primitive
lattice containing $K_\Lambda$.  Solving invariance for a row covector
$(a,b,c,d)$ gives $b=2c$, $d=3c$ from $A_1$ and $b=c$, $d=2c$ from
$A_2$; hence $b=c=d=0$.  Therefore
\begin{equation}
 \operatorname{Hom}(\Lambda,\Z)^{\pi_1(B^\circ)}
 =\Z\varepsilon_\gamma.
 \label{eq:unique-invariant-character}
\end{equation}

\subsection{The native cusp quotient and the invariant circle quotient}

Let $D_0=\{|t|\le\epsilon\}$ be small enough to contain no singular
value except $p_0$, and put $N_0=f^{-1}(D_0)$.  Set
\begin{equation}
 \overline\Lambda=\Lambda/K_\Lambda,\qquad
 \overline T^2=(\overline\Lambda\otimes\R)/\overline\Lambda,
 \label{eq:bar-torus}
\end{equation}
with ordered basis induced by $(\widehat\gamma,\widehat u)$.

\begin{lemma}[Native quotient on the punctured cusp]
\label{lem:native-quotient}
There is a real-analytic fibrewise group epimorphism
\begin{equation}
 \overline\pi:N_0\setminus W\longrightarrow D_0^*\times\overline T^2
 \label{eq:bar-pi}
\end{equation}
with connected two-torus kernel.  Its target local system is trivial.  If
\begin{equation}
 Y_T=\begin{pmatrix}6m_T&L_T\\-q_T&m_T\end{pmatrix},
 \qquad Q_T=\R^2/Y_T\Z^2,
 \label{eq:Y-and-Q}
\end{equation}
then
\begin{equation}
 \iota_T:\overline T^2\longrightarrow Q_T,\qquad
 [(a,b)]\longmapsto[Y_T(a,b)^{\mathsf T}]
 \label{eq:iota}
\end{equation}
is an isomorphism and
\begin{equation}
 \pi_T:=\iota_T\circ\overline\pi|_{F_T}:
 F_T\longrightarrow Q_T
 \label{eq:pi-T}
\end{equation}
is the quotient induced by
$\rho(x_1,x_2,x_3,x_4)=(x_2,x_4)$.  With the Euclidean quotient metric,
\begin{equation}
 \sys(Q_T)\ge T-C
 \label{eq:QT-systole}
\end{equation}
for a constant $C$ independent of sufficiently large $T$.
\end{lemma}

\begin{proof}
On the universal cover of $D_0^*$ use the real-analytic flat marking
$(z,x+\Lambda)\mapsto[z,\Pi(z)x]$.  A deck generator acts by
$(z,x+\Lambda)\mapsto(g_0z,M_0x+\Lambda)$.  By
\cref{eq:cusp-monodromy-action}, $M_0$ preserves $K_\Lambda$ and induces
the identity on $\overline\Lambda$.  Hence
\[
 [z,\Pi(z)x]\longmapsto([z],[x\bmod K_\Lambda])
\]
is deck invariant and descends to \cref{eq:bar-pi}.  The sublattice
$K_\Lambda$ is generated by two members of a basis and is primitive, so
the fibre kernel is connected.

The four columns of $P_T$ are
$v_{\gamma,T},v_{u,T},e_1,e_3$.  Projection by $\rho$ kills $e_1,e_3$
and maps the first two columns to the columns of $Y_T$.  Since
\begin{equation}
 \det Y_T=L_Tq_T+6m_T^2=D_T>0,
 \label{eq:Y-determinant}
\end{equation}
$Y_T$ is invertible, \cref{eq:iota} is a torus isomorphism and $\pi_T$
has kernel
$(\R e_1\oplus\R e_3)/(\Z e_1\oplus\Z e_3)$.  Finally
\begin{equation}
 Y_T=T\begin{pmatrix}0&1\\-1&0\end{pmatrix}
 +\begin{pmatrix}6m_T&H_T\\B_T-H_T&m_T\end{pmatrix}.
 \label{eq:Y-asymptotic}
\end{equation}
The second matrix is uniformly bounded by \cref{eq:radial-errors}.  If
$C$ bounds its operator norm, the reverse triangle inequality gives
\[
 \|Y_T(a,b)^{\mathsf T}\|\ge(T-C)\sqrt{a^2+b^2}
\]
for every $(a,b)\in\Z^2$.  This proves \cref{eq:QT-systole}.
\end{proof}

The full monodromy has a different, one-dimensional quotient.  In the flat
marking define
\begin{equation}
 \chi_\gamma([z,\Pi(z)x])=
 ([z],\varepsilon_\gamma(x)\bmod\Z).
 \label{eq:global-character}
\end{equation}
\Cref{eq:unique-invariant-character} makes this a global
real-analytic fibrewise group epimorphism
\begin{equation}
 \chi_\gamma:J\longrightarrow B^\circ\times(\R/\Z)
 \label{eq:global-circle-map}
\end{equation}
and put
\begin{equation}
 \theta_\gamma:=\operatorname{pr}_{\R/\Z}\circ\chi_\gamma:
 J\longrightarrow\R/\Z.
 \label{eq:global-circle-coordinate}
\end{equation}
Thus \(\chi_\gamma=(f,\theta_\gamma)\); in particular, scalar circle
coordinates below use \(\theta_\gamma\), not the product-valued map
\(\chi_\gamma\).
with connected three-torus kernel determined by $L_\Lambda$.  It is a real
quotient, not a holomorphic quotient of complex tori.

On $F_T$, its differential is the dual vector
\begin{equation}
 P_T^{-\mathsf T}e_\gamma
 =\left(0,\frac{m_T}{D_T},0,-\frac{L_T}{D_T}\right)^{\mathsf T},
 \qquad e_\gamma=(1,0,0,0)^{\mathsf T}.
 \label{eq:circle-dual}
\end{equation}
Thus the quotient circle, written in arclength coordinates as
\begin{equation}
 C_T=\R/\ell_T\Z,
 \label{eq:CT}
\end{equation}
has circumference
\begin{equation}
 \ell_T=\frac{1}{\|P_T^{-\mathsf T}e_\gamma\|}
 =\frac{D_T}{\sqrt{L_T^2+m_T^2}}
 =T+H-B+O(T^{-1}).
 \label{eq:circle-length}
\end{equation}
Indeed, \cref{eq:radial-errors} gives
$D_T=T^2+(2H-B)T+O(1)$ and
$\sqrt{L_T^2+m_T^2}=T+H+O(T^{-1})$.

\subsection{The two finite Seifert fillings}

We record the filling data because the full-base cone construction uses the
actual target overlap maps.  Let $\Delta_j$ be a linearizing disc above
$D_j$ with coordinate $s$, where
$s\mapsto e^{-2\pi i/m_j}s$.  Put
\begin{equation}
 \begin{aligned}
 (m_1,v_1,\kappa_1)
  &=(3,\widehat\gamma+2\widehat u-4\widehat w,1),\\
 (m_2,v_2,\kappa_2)
  &=(4,-\widehat\gamma-3\widehat u+3\widehat w,-1),
 \end{aligned}
 \qquad \kappa_j=\varepsilon_\gamma(v_j).
 \label{eq:finite-data}
\end{equation}
Matrix multiplication using \cref{eq:monodromy-matrices} gives
$A_jv_j=v_j$.  On the un-conjugated flat family the generator is
\begin{equation}
 \Phi_j(s,x)=
 \left(e^{-2\pi i/m_j}s,A_jx+\frac{v_j}{m_j}\right).
 \label{eq:finite-affine-action}
\end{equation}
It has order $m_j$.  Every nonidentity power is free: away from $s=0$
the base point moves, while over $s=0$ a fixed point would imply that the
primitive invariant character changes integrally, contrary to
\begin{equation}
 \varepsilon_\gamma
 \left(A_j^kx+\frac{k v_j}{m_j}\right)
 =\varepsilon_\gamma(x)+\frac{k\kappa_j}{m_j}\pmod{\Z},
 \quad 0<k<m_j.
 \label{eq:finite-character-shift}
\end{equation}
The quotient is the smooth finite filling $N_j\to D_j$.

The character descends equivariantly to
\begin{equation}
 C_j=\frac{\Delta_j\times(\R/\Z)}
 {(s,\theta)\sim
  (e^{-2\pi i/m_j}s,\theta+\kappa_j/m_j)}.
 \label{eq:Seifert-target}
\end{equation}
The action is free because $\gcd(\kappa_j,m_j)=1$.  To see the topology
without an external classification theorem, put $u=e^{2\pi i\theta}$ and
choose $r_j\kappa_j\equiv1\pmod{m_j}$; one may take
$r_1=1,r_2=-1$.  The invariant coordinates
\begin{equation}
 w=su^{r_j},\qquad U=u^{m_j}
 \label{eq:solid-torus-chart}
\end{equation}
identify $C_j$ real-analytically with a solid torus, and
\begin{equation}
 t_j=s^{m_j}=w^{m_j}U^{-r_j}.
 \label{eq:Seifert-multiplicity}
\end{equation}
Thus $C_j\to D_j$ has a multiple central circle of order $m_j$.  The map
$N_j\to C_j$ has fibres modeled on the connected torus
$(L_\Lambda\otimes\R)/L_\Lambda$ and is a smooth proper $\T^3$-bundle;
no group zero is asserted on the affine central filling.

It remains important that this local target agrees with the product circle
on the punctured overlap.  Let
$\Pi_\R(z):\Lambda\otimes\R\to\C^2$ be the underlying real isomorphism.
For a logarithm branch $\ell$ define $b_j^\ell(z)$ uniquely by
\begin{equation}
 \Pi_\R(z)b_j^\ell(z)
 =\frac{\ell(z)}{2\pi i}\Pi(z)v_j,
 \qquad
 \beta_j(z)=\varepsilon_\gamma(b_j^\ell(z))\bmod\Z.
 \label{eq:beta-overlap}
\end{equation}
Changing $\ell$ by $2\pi i n$ adds $nv_j$ and changes the last expression
by $n\kappa_j\in\Z$, so $\beta_j$ is branch independent.  Directly from
the changed logarithm branch,
\begin{equation}
 \beta_j(g_jz)=\beta_j(z)-\frac{\kappa_j}{m_j}\pmod\Z.
 \label{eq:beta-cocycle}
\end{equation}
Therefore
\begin{equation}
 h_j^{S^1}(z,\theta)=(z,\theta+\beta_j(z))
 \label{eq:target-overlap}
\end{equation}
conjugates the affine target action in \cref{eq:Seifert-target} to the
linear product-circle action.  If
$h_j:N_j|_{D_j^*}\to J|_{D_j^*}$ is the logarithmic source overlap and
$\overline\chi_j:N_j\to C_j$ the descended map, equivariance gives the
exact identity
\begin{equation}
 \chi_\gamma\circ h_j=h_j^{S^1}\circ\overline\chi_j.
 \label{eq:overlap-identity}
\end{equation}

\subsection{Why the literal circle map cannot cross the cusp}

The cone used later changes the codomain at $p_0$; this is forced by a
pointwise obstruction.  On the radial ray write
\begin{equation}
 Z_T=iTB_0+C(e^{-2\pi T}),\qquad
 B_0=\begin{pmatrix}0&1\\-1&0\end{pmatrix}.
 \label{eq:toric-period}
\end{equation}
For $a=(a_1,a_2)^{\mathsf T}\in\R^2$, let $x_T(a)$ be represented by
$Z_Ta$ with zero short-period coordinates.  In the dense cusp torus,
$x_T(a)=\exp(2\pi iZ_Ta)$, and its normalized logarithmic position is
\begin{equation}
 y_T(a)=B_0a+\frac{\Im C(e^{-2\pi T})a}{T},
 \qquad
 \chi_{\gamma,T}(x_T(a))=a_1\bmod\Z.
 \label{eq:toric-paths}
\end{equation}
In the affine toric chart of
$\operatorname{conv}(0,e_1,e_2)$ the coordinates are
\begin{equation}
 z_0=\frac{t}{x_1x_2},\qquad z_1=x_1,\qquad z_2=x_2.
 \label{eq:toric-chart}
\end{equation}
If $B_0a=(y_1,y_2)$ is in the interior of this triangle, the three
absolute values are
$|t|^{1-y_{T,1}-y_{T,2}},|t|^{y_{T,1}},|t|^{y_{T,2}}$ and all tend to
zero.  Hence all such paths approach the same point of $W$.

Take
\begin{equation}
 a=(-\tfrac14,\tfrac14),\qquad
 a'=(-\tfrac12,\tfrac14).
 \label{eq:two-obstruction-paths}
\end{equation}
Their limiting positions are $(1/4,1/4)$ and $(1/4,1/2)$, both in the
triangle interior, whereas their character values are $3/4$ and $1/2$.
Thus $\chi_\gamma$ has no continuous extension through $W$.  Repeating
one of these two alternatives in one factor and fixing the paths in the
other factors proves the same for every product retaining this literal
character.  The conical target below succeeds precisely because it collapses
the entire limiting boundary torus.

\begin{lemma}[Exact balls in Euclidean and collapsed-set quotients]
\label{lem:exact-quotient-balls}
Let \(L\subset\R^n\) be a full lattice, let
\[
 d_L([x],[y])=\min_{\lambda\in L}\|x-y+\lambda\|_2,
 \qquad
 \sys(L)=\min_{0\ne\lambda\in L}\|\lambda\|_2,
\]
and let \(q_L:\R^n\to\R^n/L\) be the quotient projection.  If
\(\sys(L)>4R\), then \(q_L\) restricts to an isometry from the Euclidean
closed ball \(\overline B_R(0)\) \emph{onto} the closed radius-\(R\) ball
about \([0]\) in \((\R^n/L,d_L)\).

If \((E,d)\) is a compact length space, \(A\subset E\) is closed, and
\(q_A:E\to E/A\) is equipped with the quotient length metric, then
\begin{equation}
 d_{E/A}(q_A(x),q_A(y))
 =\min\{d(x,y),d(x,A)+d(y,A)\}.
 \label{eq:collapsed-set-metric}
\end{equation}
If \(d(p,A)>2R\), the restriction of \(q_A\) is an isometry from
\(\overline B_R(p)\) onto \(\overline B_R(q_A(p))\).  Thus any Euclidean
metric ball on the left remains exactly the same Euclidean ball after the
collapse.
\end{lemma}

\begin{proof}
For \(x,y\in\overline B_R(0)\) and \(0\ne\lambda\in L\),
\[
 \|x-y+\lambda\|_2
 \geq\|\lambda\|_2-\|x-y\|_2
 >4R-2R\geq\|x-y\|_2.
\]
Hence the zero translate is the unique shortest translate on the ball.
Conversely, the minimum defining \(d_L([0],[y])\) is attained because
\(L\) is discrete.  If that distance is at most \(R\), a minimizing
representative lies in \(\overline B_R(0)\).  This proves both isometry and
surjectivity onto the quotient ball.

For the second assertion, every quotient path either avoids the collapsed
point and has length at least \(d(x,y)\), or meets it and has length at
least \(d(x,A)+d(y,A)\); both bounds are attained as infima by paths in
\(E\).  This proves \cref{eq:collapsed-set-metric}.  If
\(x,y\in\overline B_R(p)\), then
\[
 d(x,A)+d(y,A)
 \geq2d(p,A)-d(p,x)-d(p,y)>2R\geq d(x,y),
\]
so no path through the collapsed point shortens their distance.  If
\(d_{E/A}(q_A(p),q_A(y))\leq R\), the second entry in the minimum is
greater than \(2R\), and therefore \(d(p,y)\leq R\).  This proves the
onto-ball assertion.
\end{proof}

\section{The native two-torus route and its
\texorpdfstring{$S^{12}$}{S12} source}

Define the compact Hausdorff space
\begin{equation}
 \mathcal A=
 \frac{[0,\epsilon]\times S^1\times\overline T^2}
 {\{0\}\times S^1\times\overline T^2\sim v_0,
  \quad\{\epsilon\}\times S^1\times\overline T^2\sim v_1}.
 \label{eq:A-definition}
\end{equation}
Rescaling the interval gives
\begin{equation}
 \mathcal A\cong\Susp(S^1\times\overline T^2)\cong\Susp(T^3).
 \label{eq:A-suspension}
\end{equation}
The first circle remembers the base angle.  Forgetting it would give a
smaller target with the same map on a fixed radial fibre, but would not
record the quotient over the whole punctured cusp.

\begin{theorem}[Exact localized $S^{12}$ route]
\label{thm:S12-route}
There is a continuous surjective based quotient map
\begin{equation}
 q_2:(X,x_0)\longrightarrow(\mathcal A,v_0)
 \label{eq:q2}
\end{equation}
defined by
\begin{equation}
 \begin{aligned}
 q_2(x)&=[r,e^{i\phi},\overline\pi_f(x)]
 &&\text{if }x\in N_0\setminus W,
 \quad t(f(x))=re^{i\phi},\ 0<r<\epsilon,\\
 q_2(W)&=v_0,
 &q_2(X\setminus\operatorname{int}N_0)&=v_1,
 \end{aligned}
 \label{eq:q2-formula}
\end{equation}
where $\overline\pi_f$ is the fibre component of \cref{eq:bar-pi}; the
boundary $f^{-1}(|t|=\epsilon)$ maps to $v_1$.

Its smash square is a global based map
\begin{equation}
 q_2^{\wedge2}:X\wedge X\longrightarrow
 \mathcal A\wedge\mathcal A,
 \qquad X\wedge X\cong S^{12}.
 \label{eq:S12-map}
\end{equation}
Let \(\sigma_Z^{(m)}:Z^m\to Z^{\wedge m}\) denote the smash projection.
On a synchronized radial product, the restrictions
\begin{align*}
 e_{X,T}^{(2)}&:=\sigma_X^{(2)}|_{F_T^2}:F_T^2\hookrightarrow X^{\wedge2},\\
 e_{\mathcal A,T}^{(2)}(y_1,y_2)
 &:=\sigma_{\mathcal A}^{(2)}
 \bigl([r_T,1,\iota_T^{-1}(y_1)],
       [r_T,1,\iota_T^{-1}(y_2)]\bigr):
 Q_T^2\hookrightarrow\mathcal A^{\wedge2}
\end{align*}
are embeddings, where \(r_T=e^{-2\pi T}\).  The restriction of the smash
map factors exactly as
\begin{equation}
 q_2^{\wedge2}\circ e_{X,T}^{(2)}
 =e_{\mathcal A,T}^{(2)}\circ(\pi_T\times\pi_T).
 \label{eq:S12-fibre-map}
\end{equation}
The product quotient \(\pi_T\times\pi_T\) is not an embedding: its
connected kernel is the four-torus
\(\ker(\pi_T)^2\cong T^4\).

Equip \(\R^4=\R^2\oplus\R^2\) with its Euclidean \(\ell^2\) norm and
equip
\[
 Q_T^2=\R^4/\operatorname{diag}(Y_T,Y_T)\Z^4
\]
with the explicit quotient distance
\begin{equation}
 d_{Q_T^2}([x],[y])
 =\min_{n\in\Z^4}
   \|x-y+\operatorname{diag}(Y_T,Y_T)n\|_2.
 \label{eq:Q-product-metric}
\end{equation}
This is the intrinsic Euclidean \(\ell^2\)-product metric; it is assigned
to the slice independently of any ambient metric on
\(\mathcal A^{\wedge2}\).  With this metric,
\begin{equation}
 \sys(Q_T^2)\ge T-C,\qquad
 (Q_T^2,(0,0))\longrightarrow(\R^4,0)
 \label{eq:S12-limit}
\end{equation}
in pointed local-isometric, pointed smooth, and pointed
Gromov--Hausdorff topology.  Once $T-C>4R$, the universal quotient maps
the Euclidean radius-$R$ ball isometrically \emph{onto} the corresponding
metric ball in \(Q_T^2\).

The reduced integral homology is
\begin{equation}
 \widetilde H_k(\mathcal A;\Z)\cong
 \begin{cases}
 \Z^3,&k=2,3,\\
 \Z,&k=4,\\
 0,&\text{otherwise},
 \end{cases}
 \label{eq:A-homology}
\end{equation}
and
\begin{equation}
 \widetilde H_k(\mathcal A\wedge\mathcal A;\Z)\cong
 \begin{cases}
 \Z^9,&k=4,\\
 \Z^{18},&k=5,\\
 \Z^{15},&k=6,\\
 \Z^6,&k=7,\\
 \Z,&k=8,\\
 0,&\text{otherwise}.
 \end{cases}
 \label{eq:A-smash-homology}
\end{equation}
In particular, the target is not $S^4$, $S^{12}$, or a compact
four-manifold.
\end{theorem}

\begin{proof}
The equivalence relation in \cref{eq:A-definition} is the union of the
diagonal and the squares of the two closed end sections, hence is closed.
A compact Hausdorff space modulo a closed equivalence relation is compact
Hausdorff, and the quotient is exactly the unreduced suspension.

The interior formula is single valued by \cref{lem:native-quotient}; a
turn in $\phi$ changes the marking by $M_0$, which is the identity modulo
$K_\Lambda$.  Let $U$ be a neighborhood of $v_0$.  Its inverse image in
the cylinder contains the compact end section
$\{0\}\times S^1\times\overline T^2$, and hence contains
$[0,\delta)\times S^1\times\overline T^2$ for some $\delta>0$.
Therefore $q_2^{-1}(U)$ contains $f^{-1}(|t|<\delta)$, a neighborhood of
$W$.  The same compactness argument at $v_1$ gives uniform continuity at
the outer boundary; together with the interior continuity and closed
pasting, this proves continuity on $X$.

Every interior target point has a preimage by fibrewise surjectivity of
$\overline\pi$, while the nonempty cusp fibre and exterior map to the two
vertices.  Thus $q_2$ is surjective.  Since $X$ is compact and
$\mathcal A$ Hausdorff, it is a quotient map.

Because $q_2(x_0)=v_0$, functoriality of the smash quotient produces
\cref{eq:S12-map}.  The standard quotient cell structure gives
$S^m\wedge S^n\cong S^{m+n}$; applying the based diffeomorphism
\cref{eq:source-sphere} yields $X\wedge X\cong S^{12}$
\citep[Chapter~0, Smash Products]{Hatcher2002}.

The radial fibre does not contain $x_0\in W$, and its image lies in an
interior cross-section of $\mathcal A$, away from $v_0$.  Hence the source
and target product slices miss the respective fat wedges.  The restrictions
of the smash projections to those compact slices are continuous and
one-to-one.  Their targets are Hausdorff because each fat wedge is closed
in the corresponding compact Hausdorff product, so the restrictions are
the displayed embeddings.  Functoriality of the smash quotient and the
definitions of $q_2$ and $\iota_T$ give precisely
\cref{eq:S12-fibre-map}.

The period lattice of the product target is
$\operatorname{diag}(Y_T,Y_T)\Z^4$, whose systole equals that of
$Y_T\Z^2$.  If $T-C>4R$, \cref{lem:exact-quotient-balls} applied to
\cref{eq:Q-product-metric} says that the universal quotient projection is
an exact isometry \emph{onto} its radius-$R$ quotient ball.  Since this
holds for every fixed \(R\), so it proves all three convergences.

Reduced suspension shifts reduced homology by one.  The homology of $T^3$
has ranks $3,3,1$ in positive degrees $1,2,3$, proving
\cref{eq:A-homology}.  All groups are free, so the smash K\"unneth theorem
has no Tor terms.  Convolving the ranks $(3,3,1)$ with themselves gives
\[
 9,\quad18,\quad 2(3)(1)+9=15,\quad6,\quad1
\]
in degrees $4,5,6,7,8$, which proves \cref{eq:A-smash-homology}.
\end{proof}

\section{Two circle routes and their \texorpdfstring{$S^{24}$}{S24} sources}

\subsection{The localized one-circle construction}

Define
\begin{equation}
 \mathcal B=
 \frac{[0,\epsilon]\times S^1\times S^1}
 {\{0\}\times S^1\times S^1\sim b_0,
  \quad\{\epsilon\}\times S^1\times S^1\sim b_1}
 \cong\Susp(T^2).
 \label{eq:B-definition}
\end{equation}
The two interior circle coordinates are the base angle and
$\theta_\gamma$.

\begin{proposition}[Localized $S^{24}$ route]
\label{prop:S24-local}
There is a continuous surjective based quotient map
\begin{equation}
 q_{1,\mathrm{cusp}}:(X,x_0)\longrightarrow(\mathcal B,b_0)
 \label{eq:q1-cusp}
\end{equation}
given by
\begin{equation}
 \begin{aligned}
 q_{1,\mathrm{cusp}}(x)&=[r,e^{i\phi},\theta_\gamma(x)]
 &&(x\in N_0\setminus W,\ t(f(x))=re^{i\phi}),\\
 q_{1,\mathrm{cusp}}(W)&=b_0,
 &q_{1,\mathrm{cusp}}(X\setminus\operatorname{int}N_0)&=b_1.
 \end{aligned}
 \label{eq:q1-local-formula}
\end{equation}
Its fourth smash power has source
$X^{\wedge4}\cong S^{24}$.  Identify the normalized character circle with
the arclength circle by
\(\vartheta_T:\R/\Z\to C_T\),
\(\vartheta_T([\theta])=[\ell_T\theta]\).  In the following radial
factorizations, \(\chi_{\gamma,T}\) denotes the \(C_T\)-valued map
\(\vartheta_T\circ\theta_\gamma|_{F_T}\).  Define
\begin{align*}
 e_{X,T}^{(4)}&:=\sigma_X^{(4)}|_{F_T^4}:F_T^4\hookrightarrow X^{\wedge4},\\
 e_{\mathcal B,T}^{(4)}(c_1,\ldots,c_4)
 &:=\bigwedge_{i=1}^4
 [r_T,1,\vartheta_T^{-1}(c_i)]:
 C_T^4\hookrightarrow\mathcal B^{\wedge4}.
\end{align*}
Then the synchronized radial restriction factors exactly as
\begin{equation}
 q_{1,\mathrm{cusp}}^{\wedge4}\circ e_{X,T}^{(4)}
 =e_{\mathcal B,T}^{(4)}\circ\chi_{\gamma,T}^{\times4}.
 \label{eq:q1-cusp-fibre}
\end{equation}
The two displayed \(e\)-maps are embeddings, but the intervening product
quotient is not: its connected kernel is
\(\ker(\chi_{\gamma,T})^4\cong T^{12}\).
Its intrinsic target slices have the pointed $\R^4$ limit in
\cref{eq:S24-limit} below.  Moreover
\begin{equation}
 \widetilde H_k(\mathcal B;\Z)\cong
 \begin{cases}\Z^2,&k=2,\\\Z,&k=3,\\0,&\text{otherwise},\end{cases}
 \label{eq:B-homology}
\end{equation}
and
\begin{equation}
 \widetilde H_k(\mathcal B^{\wedge4};\Z)\cong
 \begin{cases}
 \Z^{16},&k=8,\\
 \Z^{32},&k=9,\\
 \Z^{24},&k=10,\\
 \Z^8,&k=11,\\
 \Z,&k=12,\\
 0,&\text{otherwise}.
 \end{cases}
 \label{eq:B-smash-homology}
\end{equation}
\end{proposition}

\begin{proof}
Monodromy invariance of $\varepsilon_\gamma$ makes the interior formula
single valued.  The two endpoint-continuity and surjectivity arguments are
the proof of \cref{thm:S12-route} with compact cross-section $T^2$ in
place of $T^3$.  Compact-to-Hausdorff again makes the map a quotient map.
The base point lies in $W$, so the fourth smash power exists.  The radial
source and target product slices avoid the closed fat wedges.  As in the
proof of \cref{thm:S12-route}, the smash projections restrict to
embeddings on those compact slices, and functoriality gives
\cref{eq:q1-cusp-fibre}.  Finally,
suspension shifts the free homology of $T^2$.  In the fourth smash power,
choosing $j$ of the four factors in degree $3$ and the remaining factors
in degree $2$ gives degree $8+j$ and rank
$\binom4j2^{4-j}$, proving the formulas.
\end{proof}

\subsection{Full-base Seifert gluing and cusp coning}

Choose pairwise disjoint closed coordinate discs $D_0,D_1,D_2$ and put
\begin{equation}
 P=B\setminus\bigcup_{j=0}^2\operatorname{int}D_j.
 \label{eq:pair-of-pants}
\end{equation}
Glue the two solid-torus targets by their exact overlap maps:
\begin{equation}
 \mathcal M=(P\times S^1)
 \underset{h_1^{S^1}}{\cup}C_1
 \underset{h_2^{S^1}}{\cup}C_2.
 \label{eq:M-definition}
\end{equation}
The overlap maps restrict to homeomorphisms of the closed boundary tori.
Consequently this is a compact orientable three-manifold with its one remaining boundary
\begin{equation}
 \partial\mathcal M=\partial D_0\times S^1\cong T^2.
 \label{eq:M-boundary}
\end{equation}
Attaching a cone along its closed boundary gives the compact Hausdorff finite
CW space
\begin{equation}
 \mathcal C=\mathcal M\cup_{\partial\mathcal M}
 \Cone(\partial\mathcal M),
 \label{eq:C-definition}
\end{equation}
and denote the cone vertex by $c_0$.

\begin{lemma}[The cone attachment as a collapsed-boundary quotient]
\label{lem:C-as-boundary-quotient}
There is a based homeomorphism
\begin{equation}
 (\mathcal C,c_0)
 \cong
 (\mathcal M/\partial\mathcal M,[\partial\mathcal M]).
 \label{eq:C-as-boundary-quotient}
\end{equation}
Moreover
\(\mathcal C\setminus\{c_0\}\cong\operatorname{int}\mathcal M\).
\end{lemma}

\begin{proof}
Choose a collar
\(h:\partial\mathcal M\times[0,1]\hookrightarrow\mathcal M\) with
\(h(y,0)=y\).  Write the attached cone as
\[
 \frac{\partial\mathcal M\times[-1,0]}
      {\partial\mathcal M\times\{-1\}\text{ collapsed}},
\]
with its \(0\)-face glued to \(h(\partial\mathcal M\times\{0\})\).
Together with the internal collar, this is
\[
 \frac{\partial\mathcal M\times[-1,1]}
      {\partial\mathcal M\times\{-1\}\text{ collapsed}}.
\]
The linear reparametrization \(t\mapsto(t+1)/2\) identifies it with
\[
 \frac{\partial\mathcal M\times[0,1]}
      {\partial\mathcal M\times\{0\}\text{ collapsed}},
\]
the image of the collar in
\(\mathcal M/\partial\mathcal M\), and it is the identity on the common
\(t=1\) face.  Extending by the identity on the complement of the open
collar gives \cref{eq:C-as-boundary-quotient}.  Removing the collapsed
point leaves \(\mathcal M\) with an external open collar attached; the
same collar reparametrization absorbs it into the internal collar and
gives \(\mathcal C\setminus\{c_0\}\cong\operatorname{int}\mathcal M\).
\end{proof}

\begin{lemma}[Homology of the Seifert target and its cone]
\label{lem:C-homology}
For boundary loops $\rho_0,\rho_1,\rho_2$ of $P$, oriented so that
$\rho_0+\rho_1+\rho_2=0$, and the product-circle class $\eta$, the filling
relations are
\begin{equation}
 3\rho_1-\eta=0,\qquad4\rho_2+\eta=0.
 \label{eq:meridian-relations}
\end{equation}
Consequently
\begin{equation}
 H_1(\mathcal M;\Z)\cong\Z\langle g\rangle,\qquad
 \rho_0=-g,\quad\rho_1=4g,\quad\rho_2=-3g,\quad\eta=12g,
 \label{eq:M-H1}
\end{equation}
$H_2(\mathcal M;\Z)=0$, and
\begin{equation}
 \widetilde H_k(\mathcal C;\Z)\cong
 \begin{cases}\Z,&k=2,3,\\0,&\text{otherwise}.
 \end{cases}
 \label{eq:C-homology}
\end{equation}
\end{lemma}

\begin{proof}
The solid-torus charts \cref{eq:solid-torus-chart}--\cref{eq:Seifert-multiplicity}
identify the two meridian slopes, with the displayed orientation, as
$3\rho_1-\eta$ and $4\rho_2+\eta$.  Hence
\[
 H_1(\mathcal M;\Z)=
 \frac{\Z\langle\eta,\rho_1,\rho_2\rangle}
 {\langle3\rho_1-\eta,4\rho_2+\eta\rangle}.
\]
The relation matrix has rank two and its $2\times2$ minors have greatest
common divisor one, so the quotient is torsion free of rank one.  The
assignments in \cref{eq:M-H1} satisfy the two relations and are primitive,
which proves the exact identification.  They also show that
$H_1(\partial\mathcal M)\to H_1(\mathcal M)$ is surjective.

An orientable compact three-manifold with one torus boundary has
$H_3(\mathcal M)=0$ and
$\chi(\mathcal M)=\tfrac12\chi(\partial\mathcal M)=0$.  Since
$b_0=b_1=1$, it follows that $b_2=0$.  Integrally,
Poincar\'e--Lefschetz duality and the universal coefficient theorem give
\[
 H_2(\mathcal M;\Z)\cong H^1(\mathcal M,\partial\mathcal M;\Z)
 \cong\operatorname{Hom}
 \bigl(H_1(\mathcal M,\partial\mathcal M;\Z),\Z\bigr),
\]
because $H_0(\mathcal M,\partial\mathcal M;\Z)=0$.  Thus
$H_2(\mathcal M;\Z)$ is free.  Its rank is $b_2=0$, so it vanishes.

The cone attachment is the homotopy cofiber of the boundary inclusion, so
\[
 \widetilde H_k(\mathcal C;\Z)
 \cong H_k(\mathcal M,\partial\mathcal M;\Z).
\]
The universal coefficient theorem gives
$H^0(\mathcal M)=H^1(\mathcal M)=\Z$ and
$H^2(\mathcal M)=H^3(\mathcal M)=0$.  Poincar\'e--Lefschetz duality now
gives \cref{eq:C-homology}.
\end{proof}

\begin{theorem}[Exact cone-compactified $S^{24}$ route]
\label{thm:S24-route}
There is a continuous surjective based quotient map
\begin{equation}
 q_1:(X,x_0)\longrightarrow(\mathcal C,c_0)
 \label{eq:q1}
\end{equation}
whose restrictions are:
\begin{enumerate}[label=\textup{(\roman*)}]
\item $(f,\theta_\gamma):f^{-1}(P)\to P\times S^1$, equivalently
$\chi_\gamma|_{f^{-1}(P)}$;
\item the finite-filling map $\overline\chi_j:N_j\to C_j$ over $D_j$,
$j=1,2$;
\item on $N_0\setminus W$, with $t(f(x))=re^{i\phi}$,
\begin{equation}
 x\longmapsto[r,(e^{i\phi},\theta_\gamma(x))]
 \in\Cone(\partial D_0\times S^1),
 \label{eq:q1-cusp-formula}
\end{equation}
while all of $W$ maps to $c_0$.
\end{enumerate}
Its fourth smash power is
\begin{equation}
 q_1^{\wedge4}:X^{\wedge4}\longrightarrow\mathcal C^{\wedge4},
 \qquad X^{\wedge4}\cong S^{24}.
 \label{eq:S24-map}
\end{equation}
On four synchronized radial fibres define the target-product embedding
\[
 e_{\mathcal C,T}^{(4)}(c_1,\ldots,c_4)
 :=\bigwedge_{i=1}^4
 [r_T,(1,\vartheta_T^{-1}(c_i))]:
 C_T^4\hookrightarrow\mathcal C^{\wedge4}.
\]
Then the radial restriction factors exactly as
\begin{equation}
 q_1^{\wedge4}\circ e_{X,T}^{(4)}
 =e_{\mathcal C,T}^{(4)}\circ\chi_{\gamma,T}^{\times4}.
 \label{eq:S24-fibre-map}
\end{equation}
The two \(e\)-maps are embeddings; the middle product quotient has the
nontrivial connected kernel \(T^{12}\) already identified above.

Equip \(C_T^4=\R^4/\ell_T\Z^4\) with the Euclidean \(\ell^2\)-product
quotient distance
\begin{equation}
 d_{C_T^4}([x],[y])
 =\min_{n\in\Z^4}\|x-y+\ell_Tn\|_2.
 \label{eq:C-product-metric}
\end{equation}
This intrinsic slice metric is assigned independently of any ambient
metric on \(\mathcal C^{\wedge4}\).  With this metric,
\begin{equation}
 \sys(C_T^4)=\ell_T=T+H-B+O(T^{-1}),
 \qquad (C_T^4,0)\longrightarrow(\R^4,0)
 \label{eq:S24-limit}
\end{equation}
in pointed local-isometric, pointed smooth, and pointed
Gromov--Hausdorff topology.  Once $\ell_T>4R$, the universal quotient maps
the Euclidean radius-$R$ ball isometrically onto the corresponding metric
ball in \(C_T^4\).

Finally,
\begin{equation}
 \widetilde H_k(\mathcal C^{\wedge4};\Z)\cong
 \begin{cases}
 \Z,&k=8,12,\\
 \Z^4,&k=9,11,\\
 \Z^6,&k=10,\\
 0,&\text{otherwise}.
 \end{cases}
 \label{eq:C-smash-homology}
\end{equation}
\end{theorem}

\begin{proof}
The first two formulas agree on their punctured overlaps by
\cref{eq:overlap-identity}.  On $\partial D_0$, the base of the cone is
glued to the remaining boundary of $P\times S^1$, so the first and third
formulas agree there.  If $U$ is a neighborhood of $c_0$, its inverse
image in the cone contains $[0,\delta)\times T^2$ for some $\delta>0$,
because the cone base is compact.  Thus $q_1^{-1}(U)$ contains
$f^{-1}(|t|<\delta)$, a neighborhood of $W$.  Continuity elsewhere comes
from the input maps; the closed pasting lemma proves continuity on $X$.

The circle quotients and both finite Seifert maps are surjective, and $W$
maps onto the cone vertex.  Hence $q_1$ is a continuous surjection from a
compact space to the Hausdorff space $\mathcal C$ and is a quotient map.
It is based, so its fourth smash power exists.  Iterating
$S^m\wedge S^n\cong S^{m+n}$ and using \cref{eq:source-sphere} identifies
the source with $S^{24}$.

The base point belongs to $W$, whereas each radial regular fibre is
disjoint from $W$.  The source and target product slices therefore avoid
the respective closed fat wedges.  The smash projections restrict to
embeddings on these compact slices, exactly as in
\cref{thm:S12-route}, and functoriality gives
\cref{eq:S24-fibre-map}.  The product target lattice is
$\ell_T\Z^4$.  Applying \cref{lem:exact-quotient-balls} to
\cref{eq:C-product-metric} proves the exact onto-ball assertion and hence
\cref{eq:S24-limit}.

The homology in \cref{eq:C-homology} is free.  In the fourth smash power,
choose $j$ factors in degree $3$ and $4-j$ factors in degree $2$.  The
result has degree $8+j$ and rank $\binom4j$, which is exactly
\cref{eq:C-smash-homology}.
\end{proof}

\section{Optimality and genuine sphere carriers}
\label{sec:optimal-sphere-carriers}

The constructions above land in conical smash targets.  We now prove two
stronger statements.  First, within the class of connected fibrewise group
quotients defined by fixed primitive sub-local systems and carrying the
unscaled Euclidean quotient norm, the numbers of copies, and hence the
sphere exponents \(12\) and \(24\), are forced by the ranks of the
expanding native quotients.  Second, canonical further quotients of the conical
targets produce honest sphere codomains containing moving metric
four-spheres whose pointed limits are exactly \(\R^4\).

\begin{theorem}[Optimality within native connected group quotients]
\label{thm:optimal-exponents}
For each \(r\geq1\), let
\(N\subset\Lambda^{\oplus r}\) be a fixed primitive sublattice which is
invariant under the monodromy group on the base under consideration, so
that it defines a connected fibrewise torus-group quotient of the native
\(r\)-fold product.  Equip that quotient with the Riemannian quotient norm
induced by the unscaled coordinate Euclidean metric on the marked original
fibres.
\begin{enumerate}[label=\textup{(\roman*)}]
\item For a connected torus quotient of \(r\) copies over the punctured
cusp, divergence of the radial quotient systole forces the quotient
lattice to have rank at most \(2r\).  A rank-four quotient therefore
requires \(r\geq2\).  Equality is attained by
\begin{equation}
 (\Lambda/K_\Lambda)^{\oplus2},
 \qquad
 L_T^{(12)}=\operatorname{diag}(Y_T,Y_T)\Z^4,
 \label{eq:S12-optimal-lattice}
\end{equation}
and the corresponding smash source is
\(X^{\wedge2}\cong S^{12}\).

\item For a connected torus quotient of \(r\) copies over the full regular
base, or over any finite cover of that base, divergence of the radial
quotient systole forces the quotient lattice to have rank at most \(r\).
A rank-four quotient therefore requires \(r\geq4\).  Equality is attained
by four copies of the full-monodromy-invariant character
\(\chi_\gamma\), whose radial period lattice is
\begin{equation}
 L_T^{(24)}=\ell_T\Z^4,\qquad
 \ell_T=T+H-B+O(T^{-1}),
 \label{eq:S24-optimal-lattice}
\end{equation}
 and the corresponding smash source is
 \(X^{\wedge4}\cong S^{24}\).
\end{enumerate}
The asserted minimality is confined to this class.  It does not exclude
nonlinear or nongroup maps, disconnected kernels, parameter-dependent
kernels or target rescalings, non-torus targets, or constructions which
replace the native product family by a different source.
\end{theorem}

\begin{proof}
Fix the primitive kernel \(N\) from the statement.  Every fixed element of
\[
 K_\Lambda^{\oplus r}
 =\Z\langle\widehat w,\widehat\delta\rangle^{\oplus r}
\]
has uniformly bounded physical norm on the radial cusp.  If a vector
\(v\) in that lattice had nonzero class in
\(\Lambda^{\oplus r}/N\), primitivity identifies that class with a nonzero
period of the connected quotient lattice, whose quotient norm is at most
the original norm and would therefore remain bounded.  The quotient
systole could not diverge.  Hence
\[
 K_\Lambda^{\oplus r}\subset N
\]
and
\[
 \rank(\Lambda^{\oplus r}/N)\leq4r-2r=2r.
\]
For \(r=2\), the primitive kernel \(K_\Lambda^{\oplus2}\) attains this
bound.  The cusp monodromy acts trivially modulo \(K_\Lambda\), so the
quotient descends over \(D_0^*\).  Its radial lattice is exactly
\cref{eq:S12-optimal-lattice}, and
\[
 \sys(L_T^{(12)})=\sys(Y_T\Z^2)\geq T-C\longrightarrow\infty.
\]

For a quotient over the full regular base, monodromy invariance gives the
stronger bound.  The inclusion \(K_\Lambda^{\oplus r}\subset N\), together
with \(A_1\widehat w=\widehat u-\widehat w\), implies
\[
 L_\Lambda^{\oplus r}
 =\Z\langle\widehat u,\widehat w,\widehat\delta\rangle^{\oplus r}
 \subset N.
\]
Consequently
\(\rank(\Lambda^{\oplus r}/N)\leq4r-3r=r\).

The conclusion survives every finite cover of \(B^\circ\).  Put
\[
 G_*=A_2A_1=I+N_*,
 \qquad N_*^2=0.
\]
Every finite-index subgroup of \(\pi_1(B^\circ)\) contains a positive
power of the loop with monodromy \(G_*\).  Direct multiplication gives
\[
 G_*^m\widehat w
 =\widehat w+m(\widehat u-\widehat\delta).
\]
If \(N\) is invariant under the finite-index monodromy group, the already
proved inclusion of \(\widehat w,\widehat\delta\) gives
\(m\widehat u\in N\).  Since \(N\) is primitive,
\(\widehat u\in N\), proving the same rank bound.

The quotient by \(L_\Lambda\) is precisely the rank-one character
\(\chi_\gamma\).  The constructions of
\cref{prop:S24-local,thm:S24-route} show that four copies attain rank
four, including the two finite Seifert fillings and the coned cusp.
Finally, the based diffeomorphism \cref{eq:source-sphere} gives
\[
 X^{\wedge r}\cong(S^6)^{\wedge r}\cong S^{6r},
\]
which proves the asserted sphere dimensions.
\end{proof}

\begin{theorem}[Sphere-to-sphere carriers with moving four-spheres]
\label{thm:sphere-carriers}
There are continuous surjective based quotient maps
\[
 \widehat q_3:(X,x_0)\longrightarrow(S^3,*),
 \qquad
 \widehat q_\gamma:(X,x_0)\longrightarrow(S^2,*),
\]
whose smash powers give continuous surjective quotient maps
\begin{equation}
 \widehat q_3^{\wedge2}:S^{12}\twoheadrightarrow S^6,
 \qquad
 \widehat q_\gamma^{\wedge4}:S^{24}\twoheadrightarrow S^8.
 \label{eq:sphere-carrier-maps}
\end{equation}
For every sufficiently large \(T\), the first target contains a based
embedded metric four-sphere
\[
 \mathbb S_{12,T}^4=S_T^2\wedge S_T^2\hookrightarrow S^6,
\]
and the second contains a based embedded metric four-sphere
\[
 \mathbb S_{24,T}^4=C_T^{\wedge4}\hookrightarrow S^8.
\]
Here the embedded subspaces carry the \(T\)-dependent quotient length
metrics specified below; no ambient metric on \(S^6\) or \(S^8\) is being
silently imposed.
On synchronized radial fibres, the exact factorizations are
\begin{align}
 F_T^2&\xrightarrow{\ \pi_T\times\pi_T\ }Q_T^2
 \xrightarrow{\ \kappa_{12,T}\ }
 S_T^2\wedge S_T^2=\mathbb S_{12,T}^4
 \lhook\joinrel\longrightarrow S^6,
 \label{eq:S12-moving-sphere-map}\\
 F_T^4&\xrightarrow{\ \chi_{\gamma,T}^{\times4}\ }
 C_T^4\xrightarrow{\ \kappa_{24,T}\ }
 C_T^{\wedge4}=\mathbb S_{24,T}^4
 \lhook\joinrel\longrightarrow S^8,
 \label{eq:S24-moving-sphere-map}
\end{align}
Here
\[
 \kappa_{12,T}(y_1,y_2)
 =\bigl[c(\iota_T^{-1}(y_1)),c(\iota_T^{-1}(y_2))\bigr]
\]
first collapses \(K^{(1)}\) in each marked two-torus and then collapses the
resulting product fat wedge, whereas
\(\kappa_{24,T}:C_T^4\to C_T^{\wedge4}\) is the smash projection and
collapses its fat wedge.  The final arrows are topological embeddings.
Neither radial map from the source product is an embedding: already its
first arrow has connected kernel \(T^4\) in the \(S^{12}\) row and
\(T^{12}\) in the \(S^{24}\) row, and the following collapse has further
nontrivial fibres.
With the quotient length metrics defined in the proof and the specified
midpoint basepoints,
\begin{equation}
 (\mathbb S_{12,T}^4,\zeta_{12,T})\longrightarrow(\R^4,0),
 \qquad
 (\mathbb S_{24,T}^4,\zeta_{24,T})\longrightarrow(\R^4,0)
 \label{eq:moving-sphere-limits}
\end{equation}
in pointed local-isometric and pointed Gromov--Hausdorff topology.  For
every \(R>0\), each radius-\(R\) ball is exactly Euclidean once
\(T-C>4R\), respectively once \(\ell_T>4R\).  The quotient metrics on
these compact four-spheres are singular at their smash points, so no
global pointed-smooth topology is invoked; the smooth assertion is exactly
\(C^\infty\) convergence on the displayed, eventually Euclidean regular
neighborhoods, which exhaust every fixed radius about the midpoint
basepoints.
\end{theorem}

\begin{proof}
For a based compact space \((Z,z_0)\), write
\[
 \Sigma_r(Z,z_0)=
 \frac{[0,\epsilon]\times Z}
 {(\{0,\epsilon\}\times Z)\cup([0,\epsilon]\times\{z_0\})}
\]
for its reduced suspension.  Use the marked presentation
\(\overline T^2=\R^2/\Z^2\), put
\[
 K^{(1)}
 =((\R/\Z)\times\{0\})\cup(\{0\}\times(\R/\Z)),
\]
and let
\[
 c:\overline T^2\longrightarrow\overline T^2/K^{(1)}\cong S^2
\]
be the degree-one quotient which collapses the one-skeleton.  The formula
\[
 \widehat p_2:\mathcal A\longrightarrow
 \Sigma_r(\overline T^2/K^{(1)},c(K^{(1)}))\cong S^3,
 \qquad
 [r,e^{i\phi},u]\longmapsto[r,c(u)]
\]
is well defined.  Both suspension ends of \(\mathcal A\), as well as the
entire cylinder over \(K^{(1)}\), map to the reduced-suspension basepoint.
It is continuous and surjective.  Define
\(\widehat q_3=\widehat p_2\circ q_2\).

Let \(y_*=[(1/2,1/2)]\in\overline T^2\).  Pull the Euclidean quotient
metric of \(Q_T\) back along \(\iota_T\); explicitly,
\begin{equation}
 \bar d_T([a],[b])
 =\min_{n\in\Z^2}\|Y_T(a-b+n)\|_2.
 \label{eq:marked-Q-metric}
\end{equation}
Give
\[
 S_T^2=\overline T^2/K^{(1)}
\]
the quotient length metric of \cref{eq:collapsed-set-metric}.  From
\cref{eq:Y-asymptotic}, for every
\(v\in\R^2\),
\[
 \|Y_Tv\|\geq(T-C)\|v\|.
\]
The standard flat distance from \(y_*\) to \(K^{(1)}\) is \(1/2\);
therefore
\begin{equation}
 d_T:=\operatorname{dist}(y_*,K^{(1)})
 \geq\frac{T-C}{2}
 \label{eq:distance-to-one-skeleton}
\end{equation}
in the pulled-back metric.  Translation of
\cref{lem:exact-quotient-balls} shows that the radius-\(R\) ball about
\(y_*\) in \((\overline T^2,\bar d_T)\) is exactly Euclidean when
\(T-C>4R\).  Under the same inequality, \(d_T>2R\), so the collapsed-set
part of that lemma maps this ball isometrically \emph{onto} the radius-
\(R\) ball about \(z_T=c(y_*)\) in \(S_T^2\).  Thus that quotient ball is
exactly a Euclidean two-ball.

At any fixed interior suspension coordinate, \(S_T^2\) is a based
embedded subspace of \(S^3\): the point \(c(K^{(1)})\) is the global
reduced-suspension basepoint.  Consequently
\[
 \mathbb S_{12,T}^4=S_T^2\wedge S_T^2
\]
embeds in \(S^3\wedge S^3=S^6\).  Equip \((S_T^2)^2\) with the
Euclidean \(\ell^2\)-product length metric
\[
 d_{2,T}((u_1,u_2),(v_1,v_2))
 =\bigl(d_{S_T^2}(u_1,v_1)^2+d_{S_T^2}(u_2,v_2)^2\bigr)^{1/2},
\]
and give the smash its quotient length metric after collapsing
\[
 (S_T^2\times\{c(K^{(1)})\})
 \cup(\{c(K^{(1)})\}\times S_T^2).
\]
The point \(\zeta_{12,T}=[z_T,z_T]\) has distance \(d_T\) from this fat
wedge.  Before the collapse, its product radius-\(R\) ball is the Euclidean
four-ball.  The collapsed-set part of
\cref{lem:exact-quotient-balls} maps that ball isometrically onto the
radius-\(R\) ball in the smash when \(T-C>4R\).  On a radial fibre,
\[
 \widehat p_2\circ q_2
 =[r,c(\iota_T^{-1}(\pi_T(\,\cdot\,)))],
\]
which proves \cref{eq:S12-moving-sphere-map}.

For the four-copy construction, choose \(1\in S^1\) as character
basepoint and define
\[
 \widehat p_1:\mathcal B\longrightarrow\Sigma_r(S^1,1)\cong S^2,
 \qquad
 [r,e^{i\phi},\theta]\longmapsto[r,\theta].
\]
This formula forgets the base angle and collapses both suspension ends
and the entire cylinder with \(\theta=1\) to the reduced-suspension
basepoint.  It is continuous and surjective.  Put
\(\widehat q_\gamma=\widehat p_1\circ q_{1,\mathrm{cusp}}\).

At an interior suspension coordinate the character circle is a based
embedded \(S^1\subset S^2\), and therefore
\[
 \mathbb S_{24,T}^4=C_T^{\wedge4}
 \subset(S^2)^{\wedge4}=S^8.
\]
Here the normalized circle \(\R/\Z\) is identified with the arclength
circle \(C_T=\R/\ell_T\Z\) by
\([\theta]\mapsto[\ell_T\theta]\).  Equip \(C_T^4\) with the Euclidean
\(\ell^2\)-product quotient metric \cref{eq:C-product-metric}, and give
the smash the quotient length metric obtained by collapsing the fat wedge on which at
least one coordinate is \(0\) modulo \(\ell_T\).  The midpoint
\[
 \zeta_{24,T}
 =[(\ell_T/2,\ell_T/2,\ell_T/2,\ell_T/2)]
\]
has distance \(\ell_T/2\) from that fat wedge.  If \(\ell_T>4R\),
\cref{lem:exact-quotient-balls} first identifies the product radius-\(R\)
ball with the Euclidean four-ball and then maps it isometrically
\emph{onto} the radius-\(R\) ball in the smash quotient.  The radial
formula is
\cref{eq:S24-moving-sphere-map}.

Both one-copy maps are continuous surjections from compact \(X\) to
Hausdorff spheres, hence quotient maps.  Their smash powers are
surjective; compact-to-Hausdorff again makes them quotient maps.  Using
\cref{eq:source-sphere},
\[
 X^{\wedge2}\cong S^{12},\quad(S^3)^{\wedge2}\cong S^6,
 \qquad
 X^{\wedge4}\cong S^{24},\quad(S^2)^{\wedge4}\cong S^8.
\]
This proves \cref{eq:sphere-carrier-maps}.  The exact onto-ball statements
prove the pointed local-isometric and Gromov--Hausdorff limits in
\cref{eq:moving-sphere-limits}, as well as \(C^\infty\) convergence on the
eventually Euclidean regular neighborhoods specified in the theorem.
\end{proof}

\begin{corollary}[The full external smash pairing is zero]
\label{cor:full-S12-map-null}
For every pair of based maps \(x,y:S^6\to S^3\),
\[
 [x\wedge y]=0\qquad\text{in}\qquad\pi_{12}(S^6).
\]
Consequently every repeated smash \(g^{\wedge r}:S^{6r}\to S^{3r}\) is
based null-homotopic for \(r\geq2\).  In particular, the continuous
surjective quotient
\(\widehat q_3^{\wedge2}:S^{12}\twoheadrightarrow S^6\) in
\cref{eq:sphere-carrier-maps} is based null-homotopic.
\end{corollary}

\begin{proof}
Bifunctoriality of the based smash product supplies the exact factorization
\[
 S^6\wedge S^6
 \xrightarrow{\operatorname{id}_{S^6}\wedge y}
 S^6\wedge S^3
 \xrightarrow{x\wedge\operatorname{id}_{S^3}}
 S^3\wedge S^3 .
\]
Under the displayed ordered sphere homeomorphisms, Toda's Barratt--Hilton
formula identifies this composite exactly as
\[
 [x\wedge y]=(E^3[x])\circ(E^6[y]):
 S^{12}\longrightarrow S^9\longrightarrow S^6.
\]
Indeed, the sign in Toda's first formula is
\((-1)^{(p+k)h}=(-1)^{18}=+1\) when \(p=q=k=h=3\)
\citep[Proposition~3.1]{Toda1962}.
Swapping the two smash inputs, rather than reversing an ill-typed
composition, introduces the source--target twist sign
\((-1)^{ph}=-1\).  Thus the same ordered sphere conventions also give
\[
 [x\wedge y]=-(E^3[y])\circ(E^6[x]).
\]
No swap is used in the proof below.

The three exact groups needed are
\[
 \pi_6(S^3)\cong\mathbb Z/12,
 \qquad
 \pi_{12}(S^9)\cong\mathbb Z/24,
 \qquad
 \pi_{12}(S^6)\cong\mathbb Z/2,
\]
\citep[opening table, p.~339]{Hatcher2002};
\citep[Table~1, p.~113]{IsaksenWangXu2023}.
Put \(\beta=E^6[y]\in\pi_{12}(S^9)\).  Since \(12[y]=0\),
functoriality gives \(12\beta=0\).  In the cyclic group
\(\mathbb Z/24\), every element annihilated by \(12\) is even; hence
\(\beta=2\gamma\) for some \(\gamma\in\pi_{12}(S^9)\).  Postcomposition
with \(E^3x:S^9\to S^6\) is a group homomorphism, and its target has
exponent two.  Therefore
\[
 [x\wedge y]=(E^3[x])_*(2\gamma)
 =2(E^3[x])_*(\gamma)=0.
\]
This generator-free argument deliberately avoids a common notation trap.
In Toda's Chapter~V, \(\nu'\) denotes the order-four \(2\)-primary
component, not an integral generator of
\(\pi_6(S^3)\cong\mathbb Z/12\).  If
\(\alpha_1(3)\) denotes the order-three component, an integral generator
may instead be written
\[
 \omega=\nu'+\alpha_1(3),
\]
as recorded explicitly by
\href{https://doi.org/10.2969/jmsj/06030767}{Maruyama--Oshima,
p.~769}.  After choosing the stable generator \(\nu_9\) compatibly, the
corresponding integral formula is
\([E^6(a\omega)]=2a\nu_9\), not the false assertion that Toda's
order-four \(\nu'\) alone generates the integral group.
Taking \(x=y=g\) proves \(g^{\wedge2}\simeq_*0\).  For \(r\geq2\),
associativity writes
\(g^{\wedge r}=g^{\wedge2}\wedge g^{\wedge(r-2)}\), and smashing the
based null-homotopy with the remaining factors proves the final assertion.
\end{proof}

\begin{remark}[The one-copy invariant and its computed value]
\label{rem:S12-one-copy-invariant}
The preceding vanishing theorem does not determine a one-copy class
\([g]\in\pi_6(S^3)\cong\Z/12\).  Stabilization identifies this group with
the even subgroup of the stable three-stem.  An integral
Serre-spectral-sequence calculation gives the same injectivity and identifies
twice the stable Hopf--\(\nu\) generator with a generator of the suspended
\(\Z/12\)-summand
\citep[Theorem~5.41(b), p.~573, and pp.~577--578]
{HatcherSpectralSequences}.  Thus, after choosing a stable generator
\(\nu_9\), there is a unique \(a\in\Z/12\) such that
\begin{equation}
 [E^6g]=2a\nu_9\in\pi_{12}(S^9)\cong\Z/24.
 \label{eq:S12-sixfold-suspension-class}
\end{equation}
Because \(a\mapsto2a\) embeds \(\Z/12\) as the even subgroup of
\(\Z/24\), \(E^6g\) is null-homotopic if and only if \(g\) is
null-homotopic.  The one-copy class and its sixfold suspension therefore
encode the same integer modulo \(12\), now computed as zero in
Theorem~\ref{cone-cusp-null:thm}.

For the specified quotient \(g=\widehat q_3\), that integer has an exact
geometric representative.  Choose \(0<r_0<\epsilon\) and
\(u_0\in\overline T^2\setminus K^{(1)}\), and let
\(y_0=[r_0,c(u_0)]\) lie in the open three-cell of \(\Sigma_rS^2\).  The
literal formula for \(q_2\) gives
\begin{equation}
 g^{-1}(y_0)
 =\{x\in f^{-1}(|t|=r_0):\overline\pi_f(x)=u_0\}.
 \label{eq:S12-regular-inverse-image}
\end{equation}
The fibre of \(\overline\pi\) is the two-torus defined by
\(K_\Lambda=\Z\langle\widehat w,\widehat\delta\rangle\), and the cusp
monodromy fixes both generators.  Hence
\cref{eq:S12-regular-inverse-image} is the mapping torus of an affine translation of the kernel torus,
which is diffeomorphic to \(S^1\times T^2\cong T^3\): if the translation is
\(k\mapsto k+v\), choose a real lift \(\widetilde v\); the map
\([\theta,k]\mapsto(e^{2\pi i\theta},k+\theta\widetilde v)\) respects
\((1,k)\sim(0,k+v)\) and has the inverse obtained by subtracting
\(\theta\widetilde v\). After a smooth chart is chosen in the open
target cell, the differential of \(g\) supplies its normal three-framing.
Pontryagin--Thom identifies the stabilized class of \(g\) with the framed
bordism class of this specifically framed \(T^3\).  Since stabilization on
\(\pi_6(S^3)\) is injective, computing that framing class computes \(a\)
exactly.  The unframed diffeomorphism type \(T^3\), point-set surjectivity,
and the zero ordinary homology map do not compute \(a\).
\end{remark}

\begin{theorem}[Exact framed-fibre reduction for the one-copy map]
\label{thm:one-copy-framing-reduction}
Let \(g=\widehat q_3:S^6\to S^3\), smooth it without changing it near the
regular inverse image in \cref{eq:S12-regular-inverse-image}, and orient
source and target.  For an oriented regular value \(y_0\), put
\(M=g^{-1}(y_0)\cong T^3\).  Then the following assertions hold.
\begin{enumerate}[label=\textup{(\roman*)}]
\item The derivative supplies an honest oriented normal framing
\(\varphi_g\) of \(M\subset S^6\), and its Pontryagin--Thom class is
\([g]\in\pi_6(S^3)\).  After adjoining the outward normal of
\(S^6\subset\R^7\), the resulting stable normal framing represents
\begin{equation}
 \operatorname{PT}(M,\varphi_g^{\mathrm{st}})
 =E^\infty[g]=2a\nu\in\pi_3^{\mathrm S}\cong\Z/24,
 \label{eq:framed-fibre-stable-class}
\end{equation}
where \(a\in\Z/12\) is the same integer as in
\cref{eq:S12-sixfold-suspension-class} and the stable generator
\(\nu\) is chosen compatibly with \(\nu_9\).

\item For every closed connected oriented three-manifold \(N\), the set
\(\operatorname{Fr}^{\mathrm{st}}_\infty(N)\) of direct-limit stable
framing classes maps onto the spin structures.  The latter form an affine
\(H^1(N;\Z/2)\)-space, and, after a spin structure is fixed, its inverse
image is an affine \(H^3(N;\Z)\)-space.  Equivalently, after one stable
framing is chosen, differences lie in
\([N,SO]\), whose three-dimensional obstruction filtration has precisely
the \(H^1(N;\Z/2)\) and \(H^3(N;\Z)\) stages.  This splitting is not
canonical.

\item Fix the spin component of \(\varphi_g^{\mathrm{st}}\).  Changing a
direct-limit stable framing by a local generator
\(\sigma\in\pi_3(SO)\cong\Z\) changes its Pontryagin--Thom class by
\(J(\sigma)=\nu\).  Changing the honest rank-three normal framing in a
small ball by a generator
\(\rho\in\pi_3(SO(3))\cong\Z\) changes the stable difference by
\(2\sigma\), and hence changes \cref{eq:framed-fibre-stable-class} by
\(2\nu\).  Consequently the honest-framing \(H^3(T^3;\Z)\)-coordinate,
reduced modulo \(12\), can realize every possible change of the
class \(a\); neither the diffeomorphism \(M\cong T^3\) nor its spin
component determines \(a\).
\end{enumerate}
Here \emph{direct-limit stable} means stabilization through
\(SO=\varinjlim SO(n)\).  It must not be confused with a framing of the
particular rank-four bundle \(\varepsilon^1\oplus TN\): the latter has
difference torsor \([N,SO(4)]\), and, after its spin component is fixed,
an affine \(H^3(N;\pi_3(SO(4)))\)-fibre.  Since
\(\pi_3(SO(4))\cong\Z\oplus\Z\), this fibre is
\(H^3(N;\Z\oplus\Z)\cong\Z\oplus\Z\) for connected closed oriented
\(N\), before passage to the direct limit.
\end{theorem}

\begin{proof}
The literal quotient formula is smooth and a submersion on a neighborhood
of the displayed regular fibre.  Choose a smaller such neighborhood \(U\).
Its complement is compact and its image omits \(y_0\), so a relative smooth
approximation, fixed on \(U\) and sufficiently close on the complement,
still omits \(y_0\) there.  It therefore has exactly the same regular
inverse image and represents the same homotopy class.

Choose an oriented orthonormal basis
\((e_1,e_2,e_3)\) of \(T_{y_0}S^3\).  At \(x\in M\), regularity makes
\[
 dg_x:\nu_x(M\subset S^6)\longrightarrow T_{y_0}S^3
\]
is an isomorphism.  Orient the normal bundle by pulling back the target
orientation through \(dg\), and then give \(M\) the unique orientation for
which \(TM\oplus\nu_{M/S^6}=TS^6|_M\) as oriented bundles.  The three vector fields
\((dg|_\nu)^{-1}e_i\) are therefore the asserted honest normal framing.
The tubular-neighborhood collapse followed by this framed identification
with \(M\times\R^3\) is the usual Pontryagin--Thom representative of the
regular-value map \(g\); choosing the displayed orientation convention
fixes its sign.  Stabilizing \(S^6\subset\R^7\) sends this representative
to \(E^\infty[g]\).  The injective stabilization
\(\pi_6(S^3)=\Z/12\hookrightarrow\pi_3^{\mathrm S}=\Z/24\) has image the
even subgroup, as recalled in
\cref{rem:S12-one-copy-invariant}; this proves (i).

After an ambient stable trivialization is fixed, a stable tangent framing
determines a stable normal framing, and conversely, by taking the
complementary summand; the two difference maps are inverse in \(SO\).
Thus the following obstruction calculation applies equally to the stable
normal framing in (i), with only the already fixed sign convention changed.
First note that every spin structure is realized by a stable framing.
Indeed, if \(v_2\) is the second Wu class, then for every
\(x\in H^1(N;\Z/2)\),
\[
 \langle v_2x,[N]\rangle=\langle Sq^2x,[N]\rangle=0
\]
by the instability axiom.  Poincar\'e duality gives \(v_2=0\), and
\(w_2(TN)=v_2+w_1(TN)^2=0\) because \(N\) is oriented.  For any chosen
spin lift, its stable classifying map lands in \(BSpin\), which is
three-connected; obstruction theory on the three-complex \(N\) makes
that lift null-homotopic, and a chosen null-homotopy is a compatible
stable framing.

Now fix one direct-limit stable framing.  Every other one is obtained
by a map \(N\to SO\), so the difference set is a torsor for \([N,SO]\).
On a three-dimensional CW structure the only relevant homotopy groups are
\[
 \pi_1(SO)=\Z/2,\qquad \pi_2(SO)=0,
 \qquad \pi_3(SO)=\Z.
\]
The first obstruction changes the spin lift and lies in
\(H^1(N;\Z/2)\).  Once that lift is fixed, there is no degree-two
obstruction and the final difference lies in \(H^3(N;\Z)\).  The
fundamental group of a topological group acts trivially on its higher
homotopy groups, so the last coefficient system is untwisted.  This is the
explicit obstruction tower stated by
\citet[Example~1.1, p.~245]{MattheyPitschScherer2010}; for honest
three-framings, the fixed-spin affine
\(H^3(N;\pi_3SO(3))\)-description and the noncanonical union over
\(H^1(N;\Z/2)\) are given by
\citet[pp.~92--93]{KirbyMelvin1999}.

For (iii), support the change of framing in an oriented embedded
three-ball and collapse its boundary.  The difference is then represented
by a map \(S^3\to SO\).  Pontryagin--Thom additivity identifies the changed
framed manifold with the connected sum with the correspondingly framed
\(S^3\), so the change of class is exactly the stable \(J\)-image; this is
\citet[Lemma~5.1 and Corollary~5.2, pp.~89--90]{Levine1985}, and in
direct-limit stable-parallelization language it is
\citet[Lemma~1.10, p.~247]{MattheyPitschScherer2010}.
Adams's image-of-\(J\) theorem makes
\(J:\pi_3(SO)=\Z\twoheadrightarrow\pi_3^{\mathrm S}=\Z/24\)
surjective \citep[Theorem~1.5, p.~22]{Adams1966}; choose signs so that
\(J(\sigma)=\nu\).  Kirby and Melvin's explicit quaternionic generators
\[
 \rho(q)x=qxq^{-1},\qquad \sigma(q)x=qx
\]
satisfy
\(\pi_3(SO(3))\to\pi_3(SO),\ \rho\mapsto2\sigma\)
\citep[pp.~91--92]{KirbyMelvin1999}.  Thus an honest local twist changes
the stable Pontryagin--Thom class by \(2\nu\).  Since \(2\nu\) generates
the even subgroup of order \(12\), its coefficient modulo \(12\) is exactly
the missing datum \(a\).

For calibration, give each factor of \(T^3=(S^1)^3\) either its bounding
stable framing or its Lie framing.  If a factor is bounding, the product
is framed-null-bordant, for instance by replacing that factor with \(D^2\).
If all three are Lie framed, multiplicativity in the framed-bordism ring
gives \(\eta^3\), and the integral stable-stem relation is
\begin{equation}
 [T^3,\varphi_{\mathrm{Lie}}]=\eta^3=12\nu
 \quad\text{in }\pi_3^{\mathrm S}=\Z/24.
 \label{eq:lie-torus-framed-class}
\end{equation}
This relation is recorded in
\citet[Proposition~4.56 and the discussion on p.~385]{Hatcher2002}.
It supplies useful reference framings in the eight spin components, but
does not identify the derivative framing \(\varphi_g\) with any one of
them.  The relative \(H^3\)-coordinate still has to be calculated from the
literal differential of \(g\).
\end{proof}

\begin{corollary}[Null-homotopy of every repeated circle-quotient smash]
\label{cor:full-S24-map-null}
Let \(f=\widehat q_\gamma:S^6\to S^2\) be the based sphere map in
\cref{eq:sphere-carrier-maps}.  Then, for every integer \(r\geq2\),
\[
 [f^{\wedge r}]=0\in\pi_{6r}(S^{2r}).
\]
In particular,
\[
 [\widehat q_\gamma^{\wedge2}]
 =0\in\pi_{12}(S^4),
 \qquad
 [\widehat q_\gamma^{\wedge4}]
 =0\in\pi_{24}(S^8).
\]
Thus the displayed continuous surjection
\(\widehat q_\gamma^{\wedge4}:S^{24}\twoheadrightarrow S^8\) is based
null-homotopic.
\end{corollary}

\begin{proof}
Bifunctoriality and associativity of the based smash product give the exact
factorization
\[
\begin{tikzcd}[column sep=large]
 S^6\wedge S^6
 \arrow[r,"f\wedge\operatorname{id}_{S^6}"]&
 S^2\wedge S^6
 \arrow[r,"\operatorname{id}_{S^2}\wedge f"]&
 S^2\wedge S^2 .
\end{tikzcd}
\]
Under the standard based homeomorphisms
\(S^a\wedge S^b\cong S^{a+b}\), the first arrow represents
\(\Sigma^6f:S^{12}\to S^8\); with the opposite suspension convention it
is conjugate to that map by the relevant twist homeomorphisms.  Freudenthal's
suspension theorem is in its isomorphism range because
\(12\leq2\cdot8-2\), and it remains in that range after every further
suspension.  Consequently
\[
 \pi_{12}(S^8)\cong\pi_4^{\mathrm S}.
\]
Here the stable-range identification follows from Freudenthal
\citep[Corollary~4.24, p.~360]{Hatcher2002}, while the fourth stable stem
is zero \citep[Table~1, p.~113]{IsaksenWangXu2023}.  Hence
\(f\wedge\operatorname{id}_{S^6}\) is based null-homotopic, so its
composite \(f\wedge f\) is based null-homotopic.

For \(r\geq2\), associativity identifies
\[
 f^{\wedge r}=(f\wedge f)\wedge f^{\wedge(r-2)},
\]
where \(f^{\wedge0}=\operatorname{id}_{S^0}\).  Smashing a based
null-homotopy with any based map gives a based null-homotopy, proving the
general assertion.  Null-homotopy does not contradict point-set
surjectivity or the exact radial metric restrictions: those are properties
of the specified representative, whereas its homotopy class forgets them.
\end{proof}

\begin{corollary}[Smash, join, and Cartesian product are distinct tests]
\label{cor:smash-join-product-separation}
Let \(g:S^6\to S^3\) and \(f:S^6\to S^2\) be based maps.  For every
\(r\geq2\), the based reduced joins of the specified quotient maps satisfy
\begin{align*}
 \underbrace{g*\cdots*g}_{r} &:S^{7r-1}\longrightarrow S^{4r-1},
 &g*\cdots*g&\simeq_*0,\\
 \underbrace{f*\cdots*f}_{r} &:S^{7r-1}\longrightarrow S^{3r-1},
 &f*\cdots*f&\simeq_*0,
\end{align*}
provided \(g=\widehat q_3\) in the first line and
\(f=\widehat q_\gamma\) in the second.  In particular, the two-fold join
of \(\widehat q_3\) is a null map \(S^{13}\to S^7\), and the four-fold
join of \(\widehat q_\gamma\) is a null map \(S^{27}\to S^{11}\).

By contrast, for every \(r\geq1\) and every based map \(h:A\to B\),
\begin{equation}
 h^{\times r}:A^r\longrightarrow B^r\text{ is null-homotopic}
 \quad\Longleftrightarrow\quad h\text{ is null-homotopic}.
 \label{eq:cartesian-power-null-iff-one-copy}
\end{equation}
Consequently the proven smash and join nullities do not determine the
Cartesian product maps
\(\widehat q_3^{\times2}\) or \(\widehat q_\gamma^{\times4}\): their
nullity is exactly the one-copy question. Theorem~\ref{cone-cusp-null:thm}
now proves all Cartesian powers of \(\widehat q_3\) null; the
one-copy question for \(\widehat q_\gamma\) remains open.
\end{corollary}

\begin{proof}
The natural based equivalence between reduced join and suspended smash is
\[
 A*B\simeq_*\Sigma(A\wedge B).
\]
Iterating it identifies the \(r\)-fold join map with
\(\Sigma^{r-1}(h^{\wedge r})\).  The sphere dimensions are
\(r(6+1)-1=7r-1\), \(r(3+1)-1=4r-1\), and
\(r(2+1)-1=3r-1\).  The asserted nullities now follow from
\cref{cor:full-S12-map-null,cor:full-S24-map-null}, since suspension
preserves a based null-homotopy.

If \(h\simeq_*0\), taking the Cartesian product of \(r\) based
null-homotopies contracts \(h^{\times r}\).  Conversely, let
\(i_j:A\to A^r\) insert the basepoint in every factor except the \(j\)-th,
and let \(p_j:B^r\to B\) be projection.  Then
\[
 p_j\circ h^{\times r}\circ i_j=h.
\]
A null-homotopy of the product therefore contracts \(h\), proving
\cref{eq:cartesian-power-null-iff-one-copy} without any cohomological or
dimensional hypothesis.
\end{proof}

\begin{corollary}[The literal open cells in the moving four-spheres]
\label{cor:moving-open-R4-cells}
Put \(K_T^{(1)}=\iota_T(K^{(1)})\subset Q_T\), and denote by \(0\in C_T\)
the character basepoint.  Removing the smash point from each moving
four-sphere gives
\begin{align}
 \mathcal U_{12,T}
 &:=(Q_T\setminus K_T^{(1)})^2
 \cong\mathbb S_{12,T}^4\setminus\{*\}
 \cong\R^4
 \hookrightarrow S^6,
 \label{eq:U12-open-cell}\\
 \mathcal U_{24,T}
 &:=(C_T\setminus\{0\})^4
 \cong\mathbb S_{24,T}^4\setminus\{*\}
 \cong\R^4
 \hookrightarrow S^8.
 \label{eq:U24-open-cell}
\end{align}
These are topologically embedded, nonclosed four-cells.  They are open in
the corresponding moving \(S^4\), but not open in the ambient \(S^6\) or
\(S^8\).  Under any fixed compatible metric on the compact sphere target,
the ambient restriction to either cell has finite diameter and therefore
cannot be globally isometric to Euclidean \(\R^4\).  The exact metric
statement is the exhaustion by the midpoint Euclidean balls of
\cref{eq:moving-sphere-limits}.
\end{corollary}

\begin{proof}
For any based Hausdorff spaces \((Z_i,z_i)\), the smash projection restricts
to a homeomorphism
\[
 \prod_i(Z_i\setminus\{z_i\})
 \xrightarrow{\ \cong\ }
 \left(\bigwedge_i Z_i\right)\setminus\{*\},
\]
because the complement of the fat wedge is saturated and the smash
equivalence relation is equality there.  Also
\[
 S_T^2\setminus\{c(K^{(1)})\}
 \cong Q_T\setminus K_T^{(1)}
 \cong(0,1)^2\cong\R^2,
\]
and
\[
 C_T\setminus\{0\}\cong(0,\ell_T)\cong\R.
\]
Taking two and four products gives the two copies of \(\R^4\).  Restricting
the established embeddings
\(\mathbb S_{12,T}^4\hookrightarrow S^6\) and
\(\mathbb S_{24,T}^4\hookrightarrow S^8\) gives the displayed embeddings.
Each cell accumulates at the omitted smash point, so it is not closed.
Its ambient nonopenness follows from invariance of dimension in local
manifold charts, since \(4<6\) and \(4<8\).  Every fixed compatible metric
on the compact ambient sphere is bounded, and so is its restriction to the
cell; Euclidean \(\R^4\) is unbounded.
\end{proof}

\begin{remark}[Null-homotopy of the moving four-spheres]
\label{rem:moving-sphere-null}
For a fixed interior suspension coordinate, the inclusion
\(S_T^2\hookrightarrow S^3\) contracts through based maps by sliding that
coordinate to a suspension endpoint.  The collapsed one-skeleton remains
at the global basepoint throughout.  Hence
\(S_T^2\wedge S_T^2\hookrightarrow S^6\) is null-homotopic.  The same
sliding argument contracts \(C_T\hookrightarrow S^2\), and therefore
contracts \(C_T^{\wedge4}\hookrightarrow S^8\).  Thus the moving metric
four-spheres do not contradict the fixed-target orientation obstruction
in \cref{thm:obstructions}: their fundamental classes die in the ambient
targets.
\end{remark}

\begin{remark}[Homotopical status of the two sphere maps]
\label{rem:sphere-map-homotopy}
The maps in \cref{eq:sphere-carrier-maps} induce zero on all
positive-degree integral homology and cohomology groups and have zero
ordinary Hurewicz image.  They have no classical degree because source
and target dimensions differ.  The class of
\(\widehat q_\gamma^{\wedge4}:S^{24}\to S^8\) is zero by
\cref{cor:full-S24-map-null}; the class of the distinct map
\(\widehat q_3^{\wedge2}:S^{12}\to S^6\) is zero by
\cref{cor:full-S12-map-null}.  These are theorems about the full specified
sphere-valued quotient maps, not merely their Hurewicz images or their
moving-\(S^4\) inclusions. The one-copy class
\([\widehat q_3]\in\pi_6(S^3)\) is zero by
\cref{cone-cusp-null:thm}, which computes the actual derivative framing.
The one-copy class \([\widehat q_\gamma]\in\pi_6(S^2)\) remains open.
The unstable based homotopy classes of the
conical-target maps
\[
 q_2^{\wedge2}:S^{12}\to\mathcal A^{\wedge2},\qquad
 q_{1,\mathrm{cusp}}^{\wedge4}:S^{24}\to\mathcal B^{\wedge4},\qquad
 q_1^{\wedge4}:S^{24}\to\mathcal C^{\wedge4}
\]
remain open.  Nullity after the final one-skeleton or base-angle quotient
does not imply nullity before that postcomposition.
\end{remark}

\subsection{Clutching classes and the exact nonextension statements}

\begin{theorem}[The \(SU(3)\) clutching calculation and its scope]
\label{thm:SU3-clutching-scope}
For each \(k\in\Z\) there is a principal \(SU(3)\)-bundle
\(P_k\to S^4\) satisfying
\begin{equation}
 \langle c_2(P_k),[S^4]\rangle=k.
 \label{eq:SU3-clutching-c2}
\end{equation}
Here and below \(c_2(P)\) means the second Chern class of the associated
standard rank-three complex vector bundle.
It may be placed on either moving four-sphere
\(\mathbb S_{12,T}^4\subset S^6\) or
\(\mathbb S_{24,T}^4\subset S^8\).  With orientations chosen so that the
corresponding quotient collapse
\[
 c_T:L_T\longrightarrow\mathbb S_{\nu,T}^4,
 \qquad
 L_T=Q_T^2\ \text{or}\ C_T^4,
 \qquad \nu=12\ \text{or}\ 24,
\]
has degree \(+1\), the pullback to the radial four-torus obeys
\begin{equation}
 \langle c_2(c_T^*P_k),[L_T]\rangle=k.
 \label{eq:radial-T4-clutching-c2}
\end{equation}
For \(k\ne0\), these bundles have the nonextension properties
\textup{(i)} and \textup{(ii)} below.  Independently, the full sphere maps
satisfy \textup{(iii)}.
\begin{enumerate}[label=\textup{(\roman*)}]
\item \(P_k\) on a moving \(S^4\) is not the restriction of a principal
\(SU(3)\)-bundle on the ambient \(S^6\) or \(S^8\).
\item \(c_T^*P_k\) is not the restriction of a principal \(SU(3)\)-bundle
on the fixed compact conical carrier
\(\mathcal D_\nu\cong\Cone(T^4)\).
\item If \(F:S^N\to S^m\) is either full sphere map in
\cref{eq:sphere-carrier-maps} and \(Q\to S^m\) is any principal
\(SU(3)\)-bundle, then \(F^*Q\to S^N\) is trivial.  Thus neither \(P_k\)
on its moving \(S^4\) nor \(c_T^*P_k\) on \(L_T\) is bundle data supplied
by that pullback: those bundles live on different base spaces and are
additional clutching choices.
\end{enumerate}
Moreover every principal \(SU(3)\)-bundle on an abstract \(S^8\) is
trivial, whereas the principal \(SU(3)\)-bundles on \(S^{12}\) form
\(\pi_{11}(SU(3))\cong\Z/4\); those four source-side classes become trivial
after stabilization to \(SU\) and are independent of the clutching integer
in \cref{eq:SU3-clutching-c2}.
\end{theorem}

\begin{proof}
Write \(S^4=D^4_+\cup_{S^3}D^4_-\), identify
\(S^3\cong SU(2)\), and let
\(i:SU(2)\hookrightarrow SU(3)\) be
\(q\mapsto\operatorname{diag}(q,1)\).
Choose a smooth based degree-\(k\) map \(d_k:S^3\to S^3\) and use
\(g_k=i\circ d_k\) as the equatorial transition map.  Since
\(SU(3)/SU(2)\cong S^5\), the long exact homotopy sequence and
\(\pi_4(S^5)=\pi_3(S^5)=0\) show that
\(i_*:\pi_3(SU(2))\to\pi_3(SU(3))\) is an isomorphism.  Hence \(g_k\) represents
\(k\in\pi_3(SU(3))\cong\Z\).  The clutching construction gives \(P_k\).

There is also a direct differential-form check of its integer.  On the
punctured hemisphere \(D^4_-\setminus\{0\}\cong(0,1]\times S^3\), let
\(\theta_k\) denote the radial pullback of \(g_k^{-1}dg_k\).  On
\(D^4_+\) take \(A_+=0\).  On the other hemisphere, choose a smooth
radial cutoff \(\chi\), identically zero on a neighborhood of its center
and one near its equator, and take \(A_-=\chi\theta_k\), extended by zero
across that neighborhood.  On the equatorial collar
the gauge-change formula is
\(A_-=g_k^{-1}A_+g_k+g_k^{-1}dg_k\), so these are the local forms of a
global connection.  The Maurer--Cartan identity gives
\begin{equation}
 F_-=d\chi\wedge\theta_k+(\chi^2-\chi)\theta_k^2.
 \label{eq:clutching-curvature}
\end{equation}
The square of the first term vanishes, the square of the second has four
tangential legs on \(S^3\) and vanishes, and the two cross terms agree.
Since
\(2\int_0^1(\chi^2-\chi)d\chi=-1/3\), Chern--Weil and transgression give
\begin{equation}
 \frac1{8\pi^2}\int_{S^4}\operatorname{tr}(F\wedge F)
 =-\frac1{24\pi^2}\int_{S^3}
       \operatorname{tr}\bigl((g_k^{-1}dg_k)^3\bigr)=k,
 \label{eq:clutching-Chern-Weil}
\end{equation}
Here \(\operatorname{tr}\) is the ordinary matrix trace in the defining
three-dimensional representation; it is not rescaled.  We choose the
orientation and the sign of the generator so that the last integral is
the degree \(k\).  Reversing both choices changes both displayed signs and
not the classification.  This proves \cref{eq:SU3-clutching-c2}.

For completeness, \(BSU(3)\) is three-connected and
\[
 \bigl(\pi_4(BSU(3)),\pi_5(BSU(3))\bigr)=(\Z,0).
\]
Hence obstruction theory on any four-dimensional CW complex \(K\)
identifies
\begin{equation}
 [K,BSU(3)]\xrightarrow{\ c_2\ }H^4(K;\Z)
 \label{eq:SU3-bundles-on-four-complex}
\end{equation}
bijectionally.  In the two-copy construction, \(L_T=Q_T^2\cong T^4\)
and \(c_T\) collapses the product fat subcomplex; in the four-copy
construction, \(L_T=C_T^4\cong T^4\) and it collapses the fat wedge.
In either case the quotient is the displayed moving \(S^4\), and its top
cell is the product of the oriented top cells of the factors.  Thus
\((c_T)_*[L_T]=[S^4]\), and naturality of \(c_2\) proves
\cref{eq:radial-T4-clutching-c2}.

The moving-sphere inclusions are null-homotopic by
\cref{rem:moving-sphere-null}.  If \(P_k\) extended over the ambient
sphere, its classifying map restricted along that null inclusion would be
null, making \(P_k\) trivial, contrary to \(k\ne0\).  Similarly, the radial
slice inclusion \(T^4\hookrightarrow\Cone(T^4)\) contracts by sliding its
radial coordinate to the cone vertex.  An extending bundle would therefore
have trivial restriction, contradicting
\cref{eq:radial-T4-clutching-c2}.  Finally, if
\(F:S^N\to S^m\) is either full sphere map and \(b:S^m\to BSU(3)\)
classifies a bundle, then the pullback is classified by \(b\circ F\).
The nullities in
\cref{cor:full-S12-map-null,cor:full-S24-map-null} make this composite
null.  These three arguments prove (i)--(iii) and show explicitly that the
clutching bundles are additional choices, not data descended through the
quotient maps.

Clutching on \(S^n\) gives
\([S^n,BSU(3)]=\pi_{n-1}(SU(3))\).  The primary computations of Mimura and
Toda give
\begin{equation}
 \pi_7(SU(3))=0,
 \qquad
 \pi_{11}(SU(3))=\Z/4
 \label{eq:SU3-low-groups}
\end{equation}
\citep[opening table, p.~217]{MimuraToda1963}.  This proves the \(S^8\)
and \(S^{12}\) assertions.  Bott periodicity gives
\(\pi_{11}(SU)\cong\Z\); the stabilization homomorphism sends the finite
group \(\Z/4\) to this torsion-free group and is therefore zero.  No
principal bundle, connection, curvature form, or transport rule has been
specified by the sphere quotient itself.
\end{proof}

\begin{proposition}[Compact cone-on-\(T^4\) carriers]
\label{prop:compact-T4-carriers}
Fix \(T_0\) sufficiently large and put
\(r_0=e^{-2\pi T_0}<\epsilon\).  Define
\[
 \mathcal D_{12}
 =\frac{[0,r_0]\times(\overline T^2)^2}
 {\{0\}\times(\overline T^2)^2\sim d_{12}},
 \qquad
 \mathcal D_{24}
 =\frac{[0,r_0]\times(S^1)^4}
 {\{0\}\times(S^1)^4\sim d_{24}}.
\]
Both spaces are homeomorphic to \(\Cone(T^4)\).  There are topological
embeddings
\begin{align}
 \delta_{12}:\mathcal D_{12}&\hookrightarrow\mathcal A\wedge\mathcal A,
 &
 [r,u_1,u_2]&\longmapsto
 [\,[r,1,u_1],[r,1,u_2]\,],
 \label{eq:D12-embedding}\\
 \delta_{24}:\mathcal D_{24}&\hookrightarrow\mathcal C^{\wedge4},
 &
 [r,\theta_1,\ldots,\theta_4]&\longmapsto
 \bigwedge_{j=1}^4[r,(1,\theta_j)].
 \label{eq:D24-embedding}
\end{align}
Here \(1\in S^1\) is the positive radial base angle.  For
\(T\geq T_0\), the restrictions of the sphere quotient maps to the
synchronized radial fibres are exactly
\begin{align}
 F_T^2&\xrightarrow{\ \pi_T\times\pi_T\ }Q_T^2
 \cong\{e^{-2\pi T}\}\times T^4\subset\mathcal D_{12},
 \label{eq:D12-slice}\\
 F_T^4&\xrightarrow{\ \chi_{\gamma,T}^{\times4}\ }C_T^4
 \cong\{e^{-2\pi T}\}\times T^4\subset\mathcal D_{24}.
 \label{eq:D24-slice}
\end{align}
Every displayed four-torus slice is null-homotopic in its carrier.  At
the vertex \(d_\nu\), for \(\nu\in\{12,24\}\), the local integral homology
is
\begin{equation}
 H_k(\mathcal D_\nu,\mathcal D_\nu\setminus\{d_\nu\};\Z)
 \cong
 \begin{cases}
  \Z^4,&k=2,\\
  \Z^6,&k=3,\\
  \Z^4,&k=4,\\
  \Z,&k=5,\\
  0,&\text{otherwise}.
 \end{cases}
 \label{eq:carrier-local-homology}
\end{equation}
In particular, neither carrier is locally a four-manifold at its vertex.
\end{proposition}

\begin{proof}
The two displayed domains are cones on
\((\overline T^2)^2\cong T^4\) and \((S^1)^4\cong T^4\).  At \(r=0\),
every factor on the right of
\cref{eq:D12-embedding,eq:D24-embedding} is the corresponding smash
basepoint, so the formulas descend from the cylinders to the cones.  For
\(0<r\leq r_0\), none of the factors is a smash basepoint.  Equality in
the smash quotient is then equality in every factor, and the coordinate
descriptions recover \(r\) and every torus coordinate.  Each
\(\delta_\nu\) is therefore injective.  Its domain is compact and its
target Hausdorff, so it is an embedding.

The restriction formulas are
\(q_2|_{F_T}=\iota_T^{-1}\circ\pi_T\) and
\(q_1|_{F_T}=\chi_{\gamma,T}\), applied in two and four factors.
Sliding the cone coordinate from \(e^{-2\pi T}\) to zero contracts each
slice to the cone vertex.  Finally, excision and the long exact sequence
of the pair give, for \(L=T^4\),
\[
 H_k(\Cone(L),\Cone(L)\setminus\{v\};\Z)
 \cong\widetilde H_{k-1}(L;\Z).
\]
The ranks \(4,6,4,1\) of \(H_1,H_2,H_3,H_4\) of \(T^4\) prove
\cref{eq:carrier-local-homology}.
\end{proof}

\begin{theorem}[Target homotopy types and point-set quotient maps]
\label{thm:target-homotopy-types}
The smash maps are continuous surjective quotient maps
\[
 q_2^{\wedge2}:S^{12}\twoheadrightarrow\mathcal A^{\wedge2},
 \qquad
 q_1^{\wedge4}:S^{24}\twoheadrightarrow\mathcal C^{\wedge4}.
\]
Their targets have homotopy types
\begin{align}
 \mathcal A^{\wedge2}
 &\simeq
 \bigvee^9S^4\vee\bigvee^{18}S^5\vee
 \bigvee^{15}S^6\vee\bigvee^6S^7\vee S^8,
 \label{eq:A-smash-wedge}\\
 \mathcal C^{\wedge4}
 &\simeq
 S^8\vee\bigvee^4S^9\vee\bigvee^6S^{10}
 \vee\bigvee^4S^{11}\vee S^{12}.
 \label{eq:C-smash-wedge}
\end{align}
Both maps induce zero on every positive-degree integral homology and
cohomology group and have zero ordinary Hurewicz image.  Neither has a
continuous homotopy section.  These conclusions do not decide whether
either unstable homotopy class is zero; both conical-target classes remain
open.  The same is true of the localized class
\(q_{1,\mathrm{cusp}}^{\wedge4}:S^{24}\to\mathcal B^{\wedge4}\).
\end{theorem}

\begin{proof}
If \(g:(K,k_0)\to(L,\ell_0)\) is a surjective based map between compact
Hausdorff spaces, every non-basepoint class of \(L^{\wedge r}\) has a
representative \((y_1,\ldots,y_r)\) with all \(y_i\ne\ell_0\).  Choosing
\(x_i\in g^{-1}(y_i)\) proves that \(g^{\wedge r}\) is surjective.  Its
source is compact, and its target is Hausdorff because the fat wedge is
closed.  Hence \(g^{\wedge r}\) is a quotient map.  Apply this to
\(q_2\) and \(q_1\).

The standard suspension splitting
\[
 \Sigma(X_1\times\cdots\times X_n)
 \simeq
 \bigvee_{\varnothing\ne I\subseteq\{1,\ldots,n\}}
 \Sigma\!\left(\bigwedge_{i\in I}X_i\right)
\]
gives
\[
 \mathcal A=\Sigma T^3
 \simeq\bigvee^3S^2\vee\bigvee^3S^3\vee S^4.
\]
Finite smash products distribute over finite wedges.  Taking ordered
degree sums gives multiplicities \(9,18,15,6,1\) in dimensions
\(4,5,6,7,8\), proving \cref{eq:A-smash-wedge}.

For \(\mathcal C\), write \(x=\rho_1\), \(y=\rho_2\), and
\(h=\eta\).  The pair-of-pants product and the two Seifert fillings give
\[
 \pi_1(\mathcal M)=
 \langle x,y,h\mid[x,h]=[y,h]=1,\,
 x^3h^{-1}=1,\,y^4h=1\rangle.
\]
The remaining boundary torus is generated by \(h\) and
\(\rho_0=(xy)^{-1}\).  Coning that torus kills their normal closure, so
\[
 \pi_1(\mathcal C)
 \cong\langle x,y\mid x^3=1,\ y^4=1,\ xy=1\rangle=1,
\]
because \(y=x^{-1}\) and \(\gcd(3,4)=1\).  By
\cref{lem:C-homology},
\[
 \widetilde H_2(\mathcal C;\Z)=\Z,\qquad
 \widetilde H_3(\mathcal C;\Z)=\Z,
\]
with no other reduced homology.  Since \(\mathcal C\) is a simply
connected finite three-complex, Hurewicz supplies a sphere representing
the \(H_2\)-generator, and the Whitehead exact sequence makes
\(\pi_3(\mathcal C)\to H_3(\mathcal C)\) surjective.  A wedge of
representing maps \(S^2\vee S^3\to\mathcal C\) is a homology equivalence
between simply connected CW complexes, hence a homotopy equivalence.
Thus \(\mathcal C\simeq S^2\vee S^3\).  Its fourth smash power gives
\cref{eq:C-smash-wedge} with binomial multiplicities \(1,4,6,4,1\).

The first target has no reduced homology above degree eight and the
second none above degree twelve.  The source spheres have no positive
homology below degrees twelve and twenty-four.  All positive-degree
homology and cohomology maps therefore vanish, as do the ordinary
Hurewicz images.  A homotopy section would split a nonzero lower homology
group of the target into the zero lower homology of the source sphere,
which is impossible.  The displayed wedges can have nonzero higher
homotopy groups, so this argument does not settle unstable null-homotopy.
\end{proof}

\begin{theorem}[Exact decompactification spaces with boundary \(\R^4\)]
\label{thm:exact-decompactification}
For \(\nu\in\{12,24\}\), set
\[
 A_T^{(12)}=\operatorname{diag}(Y_T,Y_T),\qquad
 A_T^{(24)}=\ell_T I_4,\qquad s_0=T_0^{-1}.
\]
On \([0,s_0]\times\R^4\), define a fibrewise equivalence relation
\(\sim_\nu\) by
\[
 (0,x)\sim_\nu(0,y)\iff x=y,
\]
and, for \(0<s\leq s_0\), by
\[
 (s,x)\sim_\nu(s,y)\iff
 x-y\in A_{1/s}^{(\nu)}\Z^4.
\]
Let
\[
 \mathfrak E_\nu=([0,s_0]\times\R^4)/\!\sim_\nu,
 \qquad
 p_\nu:\mathfrak E_\nu\longrightarrow[0,s_0].
\]
Then \(\mathfrak E_\nu\) is a Hausdorff, second-countable topological
five-manifold with boundary, and
\begin{align}
 p_{12}^{-1}(s)&\cong Q_{1/s}^2\quad(s>0),
 &p_{12}^{-1}(0)&\cong\R^4,
 \label{eq:E12-fibres}\\
 p_{24}^{-1}(s)&\cong C_{1/s}^4\quad(s>0),
 &p_{24}^{-1}(0)&\cong\R^4.
 \label{eq:E24-fibres}
\end{align}
The other boundary component is the compact four-torus over \(s=s_0\).

Give each positive fibre the explicit Euclidean \(\ell^2\)-quotient metric
\begin{equation}
 d_{\nu,s}([x],[y])
 =\min_{n\in\Z^4}
 \|x-y+A_{1/s}^{(\nu)}n\|_2
 \qquad(s>0),
 \label{eq:E-fibre-metric}
\end{equation}
and give the zero fibre its Euclidean metric.  For every \(R>0\), there is
\(s_R>0\) such that, for \(0\leq s<s_R\), the universal quotient
projection restricts to an exact isometry
\[
 \overline B_R(0)
 \xrightarrow{\ \cong\ }
 \overline B_R([0])\subset p_\nu^{-1}(s)
\]
onto the indicated metric ball.  Thus the positive fibres converge to the zero fibre in pointed
local-isometric, pointed smooth, and pointed Gromov--Hausdorff topology.

There are continuous surjective comparison maps
\[
 \R^4=p_{12}^{-1}(0)\hookrightarrow\mathfrak E_{12}
 \xrightarrow{\ \beta_{12}\ }\mathcal D_{12}
 \xhookrightarrow{\ \delta_{12}\ }\mathcal A^{\wedge2}
\]
and
\[
 \R^4=p_{24}^{-1}(0)\hookrightarrow\mathfrak E_{24}
 \xrightarrow{\ \beta_{24}\ }\mathcal D_{24}
 \xhookrightarrow{\ \delta_{24}\ }\mathcal C^{\wedge4}.
\]
For \(s>0\), \(\beta_\nu\) identifies the fibre with the radial torus
slice at \(r=e^{-2\pi/s}\), and
\begin{equation}
 \beta_\nu:
 \mathfrak E_\nu\setminus p_\nu^{-1}(0)
 \xrightarrow{\ \cong\ }
 \mathcal D_\nu\setminus\{d_\nu\}
 \label{eq:beta-positive-homeomorphism}
\end{equation}
is a homeomorphism.  On \(p_\nu^{-1}(0)\), it is constant with value
\(d_\nu\), and
\(\beta_\nu^{-1}(d_\nu)=p_\nu^{-1}(0)\cong\R^4\).
Nevertheless, \(\beta_\nu\) is neither proper nor a quotient map.  The
induced continuous bijection
\begin{equation}
 \overline\beta_\nu:
 \mathfrak E_\nu/p_\nu^{-1}(0)\longrightarrow\mathcal D_\nu
 \label{eq:beta-induced-bijection}
\end{equation}
is not a homeomorphism.  Thus \(\R^4\) is an actual boundary fibre of
\(\mathfrak E_\nu\), but \(\mathcal D_\nu\) is not the topological quotient
obtained from \(\mathfrak E_\nu\) through \(\beta_\nu\) by collapsing that
fibre; no topological blow-down or resolution is asserted.
\end{theorem}

\begin{proof}
Write
\[
 L_s=A_{1/s}^{(\nu)}\Z^4\quad(s>0),\qquad L_0=\{0\}.
\]
The matrices vary continuously for \(s>0\), and
\[
 \sys(L_s)\longrightarrow\infty\qquad(s\downarrow0)
\]
by \cref{eq:QT-systole,eq:circle-length}.  The equivalence relation is
closed.  Indeed, suppose
\((s_j,x_j)\sim_\nu(s_j,y_j)\) and
\((s_j,x_j,y_j)\to(s,x,y)\).  If \(s>0\), write
\(y_j-x_j=A_{1/s_j}^{(\nu)}n_j\).  On a compact neighborhood of \(s\),
the inverse matrices are uniformly bounded, so boundedness of
\(y_j-x_j\) forces \(n_j\) into a finite set.  After passage to a
subsequence it is constant, and \(x\sim_\nu y\).  If \(s=0\), every
nonzero \(n_j\) would give
\[
 \|y_j-x_j\|\geq\sys(L_{s_j})\longrightarrow\infty,
\]
contrary to convergence.  Hence \(x=y\).

The quotient projection
\[
 q_\nu:[0,s_0]\times\R^4\longrightarrow\mathfrak E_\nu
\]
is open.  If \(U\) is open, its saturation is \(U\) together with, for
each \(n\in\Z^4\), the image of
\(U\cap((0,s_0]\times\R^4)\) under the homeomorphism
\[
 (s,x)\longmapsto(s,x+A_{1/s}^{(\nu)}n).
\]
That saturation is open.  A closed equivalence relation with an open
quotient map has Hausdorff quotient: inequivalent points admit a product
neighborhood disjoint from the relation, and the saturations of the two
factor neighborhoods are disjoint.  Since the source is second countable
and the quotient is open, images of a countable base form a countable
base of \(\mathfrak E_\nu\).

Fix \(x_0\in\R^4\).  Choose \(a>0\) and \(s_1>0\) such that
\(\sys(L_s)>2a\) for \(0<s<s_1\).  The quotient projection is one-to-one
on \([0,s_1)\times B_a(x_0)\), and its restriction is an open continuous
bijection onto an open subset; hence it is a homeomorphism.  These are
half-space charts at \(s=0\).  Evenly covered small balls give product
charts for \(s>0\), and the corresponding one-sided charts at \(s=s_0\)
give the second boundary component.  This proves the manifold and fibre
claims.

For the metric assertion, choose \(s_R\) so that
\(\sys(L_s)>4R\) for \(s<s_R\).  Formula
\cref{eq:E-fibre-metric} and \cref{lem:exact-quotient-balls} show that the
universal quotient maps \(\overline B_R(0)\) isometrically onto the entire
radius-\(R\) quotient ball.  At \(s=0\) the assertion is the identity.
This proves all three pointed convergences.

For \(s>0\), put \(T=1/s\) and \(r=e^{-2\pi T}\).  In the \(S^{12}\)
case, write \(x=(x_1,x_2)\in\R^2\oplus\R^2\) and set
\[
 \beta_{12}([s,x])
 =\bigl[r,\iota_T^{-1}([x_1]_{Y_T\Z^2}),
           \iota_T^{-1}([x_2]_{Y_T\Z^2})\bigr].
\]
In the \(S^{24}\) case, write \(x=(x_1,\ldots,x_4)\) and set
\[
 \beta_{24}([s,x])
 =\bigl[r,[x_1/\ell_T],\ldots,[x_4/\ell_T]\bigr].
\]
The formulas are well defined and continuous on positive fibres.  Define
\(\beta_\nu([0,x])=d_\nu\) for every \(x\in\R^4\).  Every neighborhood
of \(d_\nu\) contains all cone cross-sections with sufficiently small
\(r\), and \(e^{-2\pi/s}\to0\) as \(s\downarrow0\).  Continuity at the
zero fibre is therefore uniform in \(x\).  Each positive fibre maps onto
its complete cross-section and the zero fibre onto the vertex, proving
surjectivity.

The homeomorphism on the positive part is also explicit.  The map
\[
 (0,s_0]\times\R^4\longrightarrow(0,s_0]\times\R^4,
 \qquad
 (s,x)\longmapsto\bigl(s,(A_{1/s}^{(\nu)})^{-1}x\bigr)
\]
is a homeomorphism and conjugates the varying lattice action to the
constant translation action of \(\Z^4\).  It therefore induces
\[
 \Phi_\nu:
 \mathfrak E_\nu\setminus p_\nu^{-1}(0)
 \xrightarrow{\ \cong\ }(0,s_0]\times(\R^4/\Z^4).
\]
The punctured cone has the corresponding product coordinates
\(\mathcal D_\nu\setminus\{d_\nu\}\cong
(0,r_0]\times(\R^4/\Z^4)\).  In these coordinates \(\beta_\nu\) is
\[
 (s,[a])\longmapsto(e^{-2\pi/s},[a]).
\]
Because \(s\mapsto e^{-2\pi/s}\) is a homeomorphism from
\((0,s_0]\) onto \((0,r_0]\), this proves
\cref{eq:beta-positive-homeomorphism}.  It also proves that the inverse
image of \(d_\nu\) is exactly the zero fibre.  Since that fibre is
\(\R^4\), the compact singleton \(\{d_\nu\}\) has noncompact preimage;
hence \(\beta_\nu\) is not proper.

It remains to prove that \(\beta_\nu\) is not a quotient map.  Fix
\[
 a=\tfrac12e_1=(\tfrac12,0,0,0)\in\R^4\setminus\Z^4,
\]
choose a strictly decreasing sequence \(s_j\downarrow0\), and set
\begin{equation}
 z_j=[s_j,A_{1/s_j}^{(\nu)}a]\in\mathfrak E_\nu.
 \label{eq:beta-counterexample}
\end{equation}
For every \(n\in\Z^4\), the vector \(e_1+2n\) is a nonzero lattice
vector, so
\begin{equation}
 \|A_{1/s_j}^{(\nu)}(a+n)\|_2
 =\tfrac12\|A_{1/s_j}^{(\nu)}(e_1+2n)\|_2
 \geq\tfrac12\sys(L_{s_j})\longrightarrow\infty.
 \label{eq:beta-counterexample-escape}
\end{equation}
Consequently \(Z=\{z_j:j\geq1\}\) is closed in \(\mathfrak E_\nu\).
Indeed, its radial coordinates have no positive accumulation, and
\cref{eq:beta-counterexample-escape} says that it eventually misses every
bounded half-space chart about every point of the zero fibre.  On the
other hand, in the product coordinates above,
\[
 \beta_\nu(z_j)=[e^{-2\pi/s_j},[a]]
 \longrightarrow d_\nu,
\]
while \(d_\nu\notin\beta_\nu(Z)\).  Thus \(\beta_\nu(Z)\) is not closed.
Because \(\beta_\nu\) is injective away from the zero fibre and no \(z_j\)
maps to \(d_\nu\),
\[
 \beta_\nu^{-1}(\beta_\nu(Z))=Z.
\]
If \(\beta_\nu\) were a quotient map, closedness of this saturated inverse
image would force \(\beta_\nu(Z)\) to be closed, a contradiction.

The constant value on \(p_\nu^{-1}(0)\) makes \(\beta_\nu\) factor through
the quotient projection to give the continuous bijection
\(\overline\beta_\nu\) in \cref{eq:beta-induced-bijection}.  Were that
bijection a homeomorphism, its composite with the quotient projection
would make \(\beta_\nu\) a quotient map.  Hence
\(\overline\beta_\nu\) is not a homeomorphism, completing every assertion.
\end{proof}

\begin{corollary}[Exact global meaning of the Euclidean limit]
\label{cor:global-R4-scope}
The unmodified four-torus fibres themselves have native pointed limit
\(\T^2\times\R^2\).  Without finite-cover unspooling, the sphere-carrier
constructions derived from that cusp produce \(\R^4\) in three other precise
senses:
\begin{enumerate}[label=\textup{(\roman*)}]
\item as the pointed intrinsic limit of the four-torus slices in the
compact quotient carriers of \(S^{12}\) and \(S^{24}\); and
\item as each nonclosed open cell
\(\mathcal U_{12,T}\subset S^6\) and
\(\mathcal U_{24,T}\subset S^8\) inside a moving four-sphere; and
\item as the decompactification boundary fibre of
\(\mathfrak E_{12}\) and \(\mathfrak E_{24}\).
\end{enumerate}
It does not produce \(\R^4\) as the whole compact sphere or compact
conical target, as a continuous quotient target of either compact source
sphere, as the quotient of \(\mathfrak E_\nu\) through \(\beta_\nu\), or
as the unrescaled ambient pointed limit of either compact conical target.
\end{corollary}

\begin{proof}
The positive assertions are
\cref{prop:compact-T4-carriers,cor:moving-open-R4-cells,thm:exact-decompactification}.
A
continuous image of a compact sphere is compact, so it cannot be
\(\R^4\).  No finite slice \(T^4\) embeds in \(\R^4\): invariance of
domain would make its image open, while compactness would make it compact,
and no nonempty open subset of \(\R^4\) is compact.  The local homology
\cref{eq:carrier-local-homology} shows that the carrier vertex is not a
four-manifold point.  The nonquotient conclusion for \(\beta_\nu\) is
\cref{eq:beta-counterexample}--\cref{eq:beta-counterexample-escape}.
Finally, the unrescaled ambient collapse is proved in
\cref{thm:uniform-collapse} below.  The comparison maps send the boundary
\(\R^4\) to the vertex, but do not give the vertex the topology of a
collapsed boundary quotient.
\end{proof}

\section{Intrinsic expansion and ambient collapse}

The metric limit just proved is a statement about the varying intrinsic
flat quotient metrics on the regular target slices.  Here ``flat'' refers
to the Riemannian metrics; no Sormani--Wenger intrinsic-flat convergence
is being asserted.  The limit cannot be interpreted as an unrescaled
limit inside one fixed metric on a compact cone.

Two metrics on an embedded slice must therefore be distinguished.  The
intrinsic slice metric is the quotient or induced Riemannian path metric
already specified on \(Q_T^2\), \(C_T^4\), or a moving \(S^4\).  If \(D\)
is a fixed ambient metric and \(i_T\) is the topological slice embedding,
the ambient restriction is instead
\[
 d^{\mathrm{amb}}_T(x,y)=D(i_T(x),i_T(y)).
\]
It allows paths which leave the slice, including shortcuts toward a cone
vertex or a collapsed fat wedge, and need not equal the intrinsic path
metric.  The next result concerns this ambient restriction; the preceding
Euclidean-ball results concern the intrinsic metric.

\begin{proposition}[Exact metric scope]
\label{prop:metric-scope}
Let fixed metrics compatible with the topologies be placed on
$\mathcal A,\mathcal B,\mathcal C$ and on their relevant smash powers.
If $j_T(Q_T^2)$ is the two-copy slice in
$\mathcal A^{\wedge2}$ and
$k_T^{\mathcal B}(C_T^4)$,
$k_T^{\mathcal C}(C_T^4)$ are the four-copy slices in the two circle
targets, then
\begin{equation}
 \begin{aligned}
 \operatorname{diam}_{\mathrm{amb}}j_T(Q_T^2)&\longrightarrow0,\\
 \operatorname{diam}_{\mathrm{amb}}k_T^{\mathcal B}(C_T^4)&\longrightarrow0,\\
 \operatorname{diam}_{\mathrm{amb}}k_T^{\mathcal C}(C_T^4)&\longrightarrow0.
 \end{aligned}
 \label{eq:ambient-collapse}
\end{equation}
In the intrinsic quotient metrics, by contrast,
\begin{equation}
 \operatorname{diam}(Q_T^2)\ge\tfrac12(T-C)\longrightarrow\infty,
 \qquad
 \operatorname{diam}(C_T^4)\ge\tfrac12\ell_T\longrightarrow\infty.
 \label{eq:intrinsic-expansion}
\end{equation}
\end{proposition}

\begin{proof}
Every neighborhood of a cone vertex contains all sufficiently small radial
cross-sections, uniformly in the compact angular and torus coordinates.
In a compatible metric, given $\delta>0$, an open $\delta/2$-ball about
the vertex therefore contains each sufficiently small cross-section, whose
diameter is then less than $\delta$.  Finite products and the continuous
smash quotient preserve this uniform convergence, proving
\cref{eq:ambient-collapse}.

Let a flat torus have systole $s$, and choose a shortest nonzero period
$\lambda$.  It is primitive.  For every period $\mu$, the nonzero period
$\lambda-2\mu$ has norm at least $s$, so
\[
 \operatorname{dist}([\lambda/2],[0])
 =\frac12\inf_\mu\|\lambda-2\mu\|\ge\frac{s}{2}.
\]
Applying this to \cref{eq:QT-systole,eq:circle-length} and observing that
a product torus contains each factor isometrically proves
\cref{eq:intrinsic-expansion}.
\end{proof}

\begin{definition}[The radial target-slice embeddings]
\label{def:target-slice-embeddings}
Put \(r_T=e^{-2\pi T}\), and define the arclength parametrization
\[
 \vartheta_T:\R/\Z\longrightarrow C_T=\R/\ell_T\Z,
 \qquad [\theta]\longmapsto[\ell_T\theta].
\]
The three synchronized radial target slices are embedded by
\begin{align*}
 j_T:Q_T^2&\longrightarrow\mathcal A^{\wedge2},\\
 (y_1,y_2)&\longmapsto
 [\,[r_T,1,\iota_T^{-1}(y_1)],
    [r_T,1,\iota_T^{-1}(y_2)]\,],
\end{align*}
\begin{align*}
 k_T^{\mathcal B}:C_T^4&\longrightarrow\mathcal B^{\wedge4},\\
 (c_1,\ldots,c_4)&\longmapsto
 \bigwedge_{i=1}^4
 [r_T,1,\vartheta_T^{-1}(c_i)],
\end{align*}
and
\begin{align*}
 k_T^{\mathcal C}:C_T^4&\longrightarrow\mathcal C^{\wedge4},\\
 (c_1,\ldots,c_4)&\longmapsto
 \bigwedge_{i=1}^4
 [r_T,(1,\vartheta_T^{-1}(c_i))].
\end{align*}
The notation in the final line uses the cusp cone
\(\Cone(\partial D_0\times S^1)\subset\mathcal C\).
\end{definition}

\begin{theorem}[Topology-forced collapse for every fixed compatible metric]
\label{thm:uniform-collapse}
Let \(D_{12}\), \(D_{24}^{\mathcal B}\), and
\(D_{24}^{\mathcal C}\) be arbitrary fixed metrics compatible with the
topologies of \(\mathcal A^{\wedge2}\),
\(\mathcal B^{\wedge4}\), and \(\mathcal C^{\wedge4}\), respectively,
and let \(*\) be the corresponding smash basepoint.  Then
\begin{align}
 d_H^{D_{12}}\bigl(j_T(Q_T^2),\{*\}\bigr)&\longrightarrow0,\notag\\
 d_H^{D_{24}^{\mathcal B}}
  \bigl(k_T^{\mathcal B}(C_T^4),\{*\}\bigr)&\longrightarrow0,\notag\\
 d_H^{D_{24}^{\mathcal C}}
  \bigl(k_T^{\mathcal C}(C_T^4),\{*\}\bigr)&\longrightarrow0.
 \label{eq:fixed-metric-collapse}
\end{align}
Thus, with metrics restricted from any one fixed ambient target, all
three target slices converge in Gromov--Hausdorff topology to a point.
\end{theorem}

\begin{proof}
We prove the first assertion.  Let
\[
 G:[0,r_0]^2\times(T^3)^2\longrightarrow\mathcal A^{\wedge2}
\]
be the product of the two inner-cylinder quotient maps followed by the
smash quotient.  It is continuous and
\(G(0,0,z_1,z_2)=*\) for every \((z_1,z_2)\).  Compactness of the domain
and uniform continuity with respect to \(D_{12}\) give
\[
 \sup_{(z_1,z_2)\in(T^3)^2}
 D_{12}\bigl(G(r,r,z_1,z_2),*\bigr)\longrightarrow0
 \qquad(r\downarrow0).
\]
The set \(j_T(Q_T^2)\) is contained in the \(r=r_T\)
cross-section.  This proves the first Hausdorff convergence.  Replacing
\(T^3\) by \(T^2\) and taking four factors proves the other two.
Hausdorff convergence to a singleton implies Gromov--Hausdorff
convergence to that singleton.
\end{proof}

\begin{theorem}[One fixed cusp-adapted metric and its complete bubbles]
\label{thm:fixed-cusp-metric}
There are fixed compatible length metrics
\[
 D_{12}\text{ on }\mathcal A^{\wedge2},\qquad
 D_{24}^{\mathcal B}\text{ on }\mathcal B^{\wedge4},\qquad
 D_{24}^{\mathcal C}\text{ on }\mathcal C^{\wedge4}
\]
with the following properties.  On the fixed marked torus
\(\overline T^2=\R^2/\Z^2\), put
\[
 h_T:=\iota_T^*g_{Q_T}
 =dx^{\mathsf T}Y_T^{\mathsf T}Y_Tdx,\qquad
 T(r):=-\frac{1}{2\pi}\log r.
\]
On the inner cusp portion of \(\mathcal A\), in normalized base-angle
coordinate \(u\in\R/\Z\), the metric is induced by
\begin{equation}
 g_{\mathcal A}
 =dr^2+r^2du^2+\frac{r^2}{T(r)^2}h_{T(r)}.
 \label{eq:fixed-A-metric}
\end{equation}
On the cusp cone of \(\mathcal B\) or \(\mathcal C\), in normalized
character coordinate \(\theta\in\R/\Z\), it is induced by
\begin{equation}
 g_{\mathcal D}=dr^2+r^2(du^2+d\theta^2),
 \qquad\mathcal D\in\{\mathcal B,\mathcal C\}.
 \label{eq:fixed-BC-metric}
\end{equation}
On every finite product use the \(\ell^2\)-product length metric
\[
 d_2((x_i),(y_i))=\left(\sum_i d_i(x_i,y_i)^2\right)^{1/2},
\]
and on smash powers take its metric quotient after collapsing the closed
fat wedge.

Let \(g_Z\) denote the resulting Riemannian tensor on each open regular
stratum, and put
\[
 p_T^{12}=j_T(0,0),\qquad
 p_T^{24,\mathcal D}=k_T^{\mathcal D}(0,0,0,0).
\]
Let \(d_{12,T}^{\mathrm{sl}}\) and
\(d_{24,T}^{\mathcal D,\mathrm{sl}}\) denote the intrinsic Riemannian path
metrics of \(j_T^*g_Z\) and \((k_T^{\mathcal D})^*g_Z\), transported to
the embedded slice images.  These are not notation for the ambient
restrictions of \(D_{12}\) and \(D_{24}^{\mathcal D}\).  The first
fundamental forms satisfy the exact identities
\begin{equation}
 \left(\frac{T}{r_T}\right)^2j_T^*g_Z
 =g_{Q_T}\oplus g_{Q_T},
 \label{eq:fixed-A-slice}
\end{equation}
and
\begin{equation}
 \left(\frac{\ell_T}{r_T}\right)^2
 (k_T^{\mathcal D})^*g_Z
 =g_{C_T}^{\oplus4}.
 \label{eq:fixed-BC-slice}
\end{equation}
Consequently
\begin{align}
 \left(j_T(Q_T^2),\frac{T}{r_T}d_{12,T}^{\mathrm{sl}},p_T^{12}\right)
 &\longrightarrow(\R^4,0),\notag\\
 \left(k_T^{\mathcal D}(C_T^4),
       \frac{\ell_T}{r_T}d_{24,T}^{\mathcal D,\mathrm{sl}},
       p_T^{24,\mathcal D}\right)
 &\longrightarrow(\R^4,0)
 \label{eq:fixed-slice-bubbles}
\end{align}
in pointed local-isometric, pointed smooth, and pointed
Gromov--Hausdorff topology.  No equality with the ambient subspace
distances is asserted.

The full ambient pointed blow-ups instead converge in pointed
Gromov--Hausdorff topology as
\begin{align}
 \left(\mathcal A^{\wedge2},
       \frac{T}{r_T}D_{12},p_T^{12}\right)
 &\longrightarrow(\R^8,0),\notag\\
 \left(\mathcal D^{\wedge4},
       \frac{\ell_T}{r_T}D_{24}^{\mathcal D},
       p_T^{24,\mathcal D}\right)
 &\longrightarrow(\R^{12},0),
 \qquad \mathcal D\in\{\mathcal B,\mathcal C\}.
 \label{eq:fixed-ambient-bubbles}
\end{align}
For every prescribed radius, the corresponding balls are eventually
contained in unwrapped regular-stratum coordinate neighborhoods, and the
pulled-back Riemannian tensors converge there in \(C^\infty\).  This is the
only smooth-convergence assertion for the globally singular smash targets.
The slice \(\R^4\) is the coordinate
subspace obtained by fixing the two radial and two base-angle coordinates
in the first line, or the four radial and four base-angle coordinates in
the second.
\end{theorem}

\begin{proof}
By \cref{eq:Y-asymptotic},
\[
 T^{-2}h_T
 =dx^{\mathsf T}(T^{-1}Y_T)^{\mathsf T}(T^{-1}Y_T)dx
 \longrightarrow dx^{\mathsf T}dx
\]
in every \(C^k\) norm.  Indeed, all nonleading cusp coefficients are
holomorphic in \(t=e^{-2\pi T}\), so their \(T\)-derivatives are
exponentially small.  Therefore the compact cross-section metrics in
\cref{eq:fixed-A-metric} remain uniformly bounded and converge to a
positive definite flat metric.  Their diameters are \(O(r)\).  The radial
coordinate is \(1\)-Lipschitz and a radial segment from level \(r\) to the
vertex has length \(r\); the length completion adds exactly one point at
\(r=0\), with the prescribed cone topology.  The same statements are
immediate for \cref{eq:fixed-BC-metric}.  Interpolation on a closed
annulus away from the vertex extends these tensors to fixed compatible
length metrics on \(\mathcal A,\mathcal B,\mathcal C\).

For a compact metric space \((E,d)\) and closed subset \(W\subset E\),
the metric quotient collapsing \(W\) to one point is
\[
 d_{E/W}([x],[y])
 =\min\{d(x,y),d(x,W)+d(y,W)\}.
\]
Applying this formula to the fat wedges in the product length spaces
gives fixed compatible metrics on the smash targets.

On the \(r_T\)-slice of \(\mathcal A\), with \(u=0\), the first
fundamental form is \(r_T^2T^{-2}h_T\).  Since
\(h_T=\iota_T^*g_{Q_T}\), this proves \cref{eq:fixed-A-slice}.  On the
\(r_T\)-slice of \(\mathcal B\) or \(\mathcal C\), with \(u=0\), the
form is \(r_T^2d\theta^2\).  Under \(\vartheta_T\), the arclength metric
of \(C_T\) pulls back to \(\ell_T^2d\theta^2\), proving
\cref{eq:fixed-BC-slice}.  The systolic estimates
\[
 \sys(Q_T^2)\geq T-C,\qquad \sys(C_T^4)=\ell_T
\]
then prove \cref{eq:fixed-slice-bubbles}: every fixed-radius ball is
eventually exactly Euclidean.  These are the intrinsic path metrics of
curves constrained to the slices; this argument does not identify them
with the restrictions of the ambient distances.

For the ambient limits, use one \(\mathcal A\)-factor and put
\[
 \lambda_T=\frac{T}{r_T},\quad
 s=\lambda_T(r-r_T),\quad \eta=Tu,\quad z=Y_Tx
\]
on the universal cover of the torus directions.  On every fixed bounded
set in \((s,\eta,z)\),
\[
 \frac r{r_T}=1+\frac sT,\qquad
 T(r)=T+O(T^{-1}),\qquad
 Y_{T(r)}Y_T^{-1}=I+O(T^{-2}),
\]
with the same estimates after any fixed number of derivatives.  Hence
\(\lambda_T^2g_{\mathcal A}\) converges to
\[
 ds^2+d\eta^2+dz_1^2+dz_2^2.
\]
The inward radial boundary is at distance \(T\), the \(u\)-period is
\(T\), and the \(z\)-period lattice is \(Y_T\Z^2\), with systole at least
\(T-C\).  The outward interface \(r=r_0\), beyond which the metric was
interpolated, is at scaled distance
\[
 \frac{T}{r_T}(r_0-r_T)\longrightarrow\infty.
\]
Thus both inward and outward exits from the displayed regular coordinate
neighborhood, as well as every periodic identification, recede to infinity.
One factor converges to \(\R^4\), and two factors to \(\R^8\).

For one \(\mathcal B\)- or \(\mathcal C\)-factor, put
\[
 \lambda_T=\frac{\ell_T}{r_T},\quad
 s=\lambda_T(r-r_T),\quad
 \eta=\ell_Tu,\quad\zeta=\ell_T\theta.
\]
Then
\[
 \lambda_T^2g_{\mathcal D}
 =ds^2+\left(\frac r{r_T}\right)^2(d\eta^2+d\zeta^2)
\]
converges smoothly on bounded coordinate sets to the Euclidean metric on
\(\R^3\).  Both circle periods and the inward radial boundary are at
scaled distance \(\ell_T\to\infty\).  The outward interface is at distance
\[
 \frac{\ell_T}{r_T}(r_0-r_T)\longrightarrow\infty,
\]
so four factors converge to \(\R^{12}\).

The product-metric distance from a synchronized point to the fat wedge is
exactly \(r_T\).  After the displayed rescalings it is \(T\), respectively
\(\ell_T\).  Thus every fixed-radius ball eventually avoids the fat
wedge.  Nor can a distance between two points of such a ball be shortened
by an outward excursion: if \(\mathcal O_T\) is the outward interface and
the basepoint-to-\(\mathcal O_T\) distance is \(a_T\), any such excursion
between points in the radius-\(R\) ball has length at least
\(2(a_T-R)>2R\) for all sufficiently large \(T\).  The identical estimate
with the fat wedge, whose scaled distance tends to infinity, excludes the
inward collapsed-set shortcut.  Hence all relevant metric balls remain
inside the unwrapped regular neighborhoods, the smash quotient does not
change their local calculations, and the \(C^\infty\) statement has exactly
the restricted scope asserted in the theorem.
This proves \cref{eq:fixed-ambient-bubbles}.
\end{proof}

\begin{theorem}[Complete scalar-rescaling trichotomy]
\label{thm:scaling-trichotomy}
For the fixed metrics of \cref{thm:fixed-cusp-metric}, let
\(\lambda_T>0\) be arbitrary and put
\[
 \beta_T=\lambda_Tr_T.
\]
On the \(Q_T^2\) slice or either \(C_T^4\) slice, pass to a subsequence
so that \(\beta_T\) has a limit in \([0,\infty]\).  The pointed limit is
\begin{equation}
 \begin{cases}
  \{*\},&\beta_T\longrightarrow0,\\
  (\R/L\Z)^4,&\beta_T\longrightarrow L\in(0,\infty),\\
  \R^4,&\beta_T\longrightarrow\infty.
 \end{cases}
 \label{eq:scaling-trichotomy}
\end{equation}
In the middle line the lattice is orthogonally equivalent to
\(L\Z^4\).  At the same time,
\begin{equation}
 \lambda_TD(p_T,*)=\beta_T
 \label{eq:vertex-distance}
\end{equation}
for every synchronized basepoint used above.  Therefore every scalar
rescaling which unrolls the slice to \(\R^4\) sends the cone/fat-wedge
basepoint to infinite scaled distance.  Such a limit is not a tangent
cone based at that point.
\end{theorem}

\begin{proof}
For the \(Q_T^2\) route, the rescaled period matrix in each
two-dimensional factor is
\[
 \frac{\lambda_Tr_T}{T}Y_T
 =\beta_T\left(
 \begin{pmatrix}0&1\\-1&0\end{pmatrix}
 +\frac1T
 \begin{pmatrix}6m_T&H_T\\B_T-H_T&m_T\end{pmatrix}
 \right).
\]
If \(\beta_T\to L\in(0,\infty)\), this lattice converges to
\(LJ\Z^2\), where \(J=\left(\begin{smallmatrix}0&1\\-1&0\end{smallmatrix}\right)\);
two factors give the square four-torus in
\cref{eq:scaling-trichotomy}.  If \(\beta_T\to0\), the lengths of the
displayed lattice basis tend to zero, and so does the diameter.  If
\(\beta_T\to\infty\), the reverse triangle inequality gives systole
\[
 \frac{\beta_T}{T}(T-C)
 =\beta_T\left(1-\frac CT\right)\longrightarrow\infty.
\]
Every fixed-radius ball is then eventually Euclidean.

For a \(C_T^4\) slice, its four rescaled periods all have length
\(\lambda_Tr_T=\beta_T\), so the same three cases follow directly.
Before the smash quotient, a synchronized point in the product lies at
distance exactly \(r_T\) from the fat wedge: moving any one radial
coordinate to zero realizes this distance, and the \(dr^2\)-term proves
the lower bound.  The metric-quotient formula in the proof of
\cref{thm:fixed-cusp-metric} proves \cref{eq:vertex-distance}.
\end{proof}

\begin{corollary}[Unit-volume normalization does not give \(\R^4\)]
\label{cor:unit-volume}
The four-volume of \(Q_T^2\) is \((\det Y_T)^2\sim T^4\).  Its
unit-volume length factor is \((\det Y_T)^{-1/2}\sim T^{-1}\), and its
unit-volume limit is a square flat four-torus, not \(\R^4\).  Likewise,
the four-volume of \(C_T^4\) is \(\ell_T^4\), and its unit-volume
normalization is exactly the square unit four-torus.
\end{corollary}

\begin{proposition}[Vertex tangents and compact asymptotic cones]
\label{prop:tangent-scope}
The compact target topology alone determines no metric tangent cone.  For
an exact conic length metric built from fixed piecewise-Riemannian metrics
on the links, however, the vertex links have topological and Hausdorff
dimensions \(7\) for \(\mathcal A^{\wedge2}\) and \(11\) for both
fourfold circle targets.  The resulting vertex cones have dimensions
\(8\) and \(12\), and are not \(\R^4\).

More precisely, one has the based topological models
\[
 \mathcal A^{\wedge2}\cong\Susp(T^3*T^3),\qquad
 \mathcal B^{\wedge4}\cong\Susp((T^2)^{*4}).
\]
At their suspension basepoints the conic links are \(T^3*T^3\) and
\((T^2)^{*4}\).  Also
\[
 \mathcal C^{\wedge4}\cong\mathcal M^4/\partial(\mathcal M^4),
\]
so an exact conic metric at the collapsed boundary has link
\(\partial(\mathcal M^4)\).  If \(*\) is its collapsed smash point, its
complete local integral homology is
\begin{equation}
 H_i(\mathcal C^{\wedge4},
     \mathcal C^{\wedge4}\setminus\{*\};\Z)
 \cong
 \begin{cases}
  \Z^4,&i=2,\\
  \Z^6,&i=3,\\
  \Z^4,&i=4,\\
  \Z,&i=5,\\
  \Z,&i=8,\\
  \Z^4,&i=9,\\
  \Z^6,&i=10,\\
  \Z^4,&i=11,\\
  \Z,&i=12,\\
  0,&\text{otherwise}.
 \end{cases}
 \label{eq:C-smash-local-homology}
\end{equation}
In particular the collapsed point is not a manifold point of any
dimension.  For the other two exact conic models, the top local homology
occurs in degree \(8\), respectively \(12\), rather than degree \(4\).

Finally, scaling the distances of a fixed compact metric target by
\(R_n^{-1}\), where \(R_n\to\infty\), gives a point:
\[
 \operatorname{diam}(Z,R_n^{-1}D)
 =R_n^{-1}\operatorname{diam}(Z,D)\longrightarrow0.
\]
Thus \cref{eq:fixed-slice-bubbles,eq:fixed-ambient-bubbles} describes
moving-basepoint blow-ups, or bubbles, not vertex-based tangent cones and
not asymptotic cones of the compact targets.
\end{proposition}

\begin{proof}
For compact spaces \(V_i\) with closed subspaces \(A_i\), the quotient
definitions give a canonical homeomorphism
\[
 (V_1/A_1)\wedge\cdots\wedge(V_m/A_m)
 \cong
 \frac{V_1\times\cdots\times V_m}
 {\bigcup_iV_1\times\cdots\times A_i\times\cdots\times V_m}.
\]
Take \(V_i=\Cone(L_i)\) and \(A_i=L_i\), its outer boundary.  A product
of \(m\) cones is homeomorphic to the cone on the iterated join
\(L_1*\cdots*L_m\), and the displayed union of faces is the outer link.
Collapsing that link gives its suspension.  Hence
\[
 \dim(T^3*T^3)=3+3+1=7,\qquad
 \dim((T^2)^{*4})=4\cdot2+3=11.
\]

By \cref{lem:C-as-boundary-quotient} and the same quotient identity, if
\(N=\mathcal M^4\) and
\[
 \partial_\times N
 =\bigcup_{j=1}^4
   \mathcal M^{j-1}\times\partial\mathcal M
   \times\mathcal M^{4-j},
\]
then
\[
 \mathcal C^{\wedge4}\cong N/\partial_\times N.
\]
Rounding the product corners replaces \(\partial_\times N\) by a
homeomorphic eleven-dimensional closed manifold \(\partial N\).  A collar
of \(\partial N\) descends to a cone neighborhood of the collapsed point,
so its link is exactly \(\partial N\).

The collar also makes
\(\operatorname{int}\mathcal M\hookrightarrow\mathcal M\) a homotopy
equivalence.  Hence, for
\[
 U=\mathcal C^{\wedge4}\setminus\{*\}
 \cong(\operatorname{int}\mathcal M)^4,
\]
the K\"unneth theorem and \cref{lem:C-homology} give
\[
 H_j(U;\Z)\cong H_j(N;\Z)\cong
 \begin{cases}
  \Z^{\binom4j},&0\leq j\leq4,\\
  0,&\text{otherwise}.
 \end{cases}
\]
On the other hand, \cref{eq:C-smash-homology} gives the absolute homology
of \(\mathcal C^{\wedge4}\): its reduced groups have ranks
\(1,4,6,4,1\) in degrees \(8,9,10,11,12\) and vanish otherwise.  The long
exact sequence of the pair
\((\mathcal C^{\wedge4},U)\) therefore yields
\[
 H_i(\mathcal C^{\wedge4},U;\Z)
 \cong H_{i-1}(U;\Z)\quad(2\leq i\leq5)
\]
and
\[
 H_i(\mathcal C^{\wedge4},U;\Z)
 \cong H_i(\mathcal C^{\wedge4};\Z)\quad(8\leq i\leq12).
\]
The map on \(H_0\) is an isomorphism and all groups in the remaining
degrees vanish.  This proves \cref{eq:C-smash-local-homology}.

For comparison, the top relative homology groups of the two suspension
models and of \(N/\partial N\) are nonzero:
\[
 H_8(\Cone(T^3*T^3),T^3*T^3;\Z)\cong\Z
\]
and, by orientation and Poincar\'e--Lefschetz duality,
\[
 H_{12}(\mathcal M^4,\partial(\mathcal M^4);\Z)\cong\Z.
\]
The \(\mathcal B\) calculation is identical in degree \(12\).
Euclidean \(\R^4\) has local homology \(\Z\) only in degree \(4\), so
none of these conic vertex models is \(\R^4\).  The diameter computation
proves the final scaling assertion.
\end{proof}

\paragraph{Metric consequence.}
For each fixed cusp-adapted compact-codomain metric, the corresponding
moving-base-point rescalings have a Euclidean bubble containing the native
four-dimensional decompactification as its vertical \(\R^4\) factor.  The
complete bubbles are \(\R^8\) and \(\R^{12}\).  The factor
\(\beta_T=\lambda_Tr_T\) proves that a different scalar rescaling of the
same fixed cusp-adapted metric cannot remove this dimensional discrepancy:
the slice unrolls only after the smash basepoint has gone to infinite
scaled distance.  No metric on a source sphere is supplied.  Pulling a
codomain metric back through \(q_2^{\wedge2}\) or \(q_1^{\wedge4}\) gives,
where differentiable, a degenerate quadratic form along the
positive-dimensional quotient kernels.

\section{Sharp obstructions}

\begin{proposition}[Three exact almost-complex obstructions on \(S^{12}\)]
\label{prop:S12-almost-complex-obstruction}
Let \(M\) be a smooth manifold homeomorphic to \(S^{12}\).  There is no
bundle endomorphism \(J:TM\to TM\) satisfying
\(J^2=-\operatorname{Id}_{TM}\).  Two independent contradictions---the
Borel--Serre reduced-power calculation at the single prime \(5\), and Bott
integrality---use only the homotopy type of \(M\).  For the standard smooth
sphere there is also an independent stable-tangent Pontryagin-class
contradiction.

These statements concern an unstabilized complex structure on the actual
rank-twelve tangent bundle.  A stable complex structure on
\(TS^{12}\oplus\varepsilon^r_{\R}\), a complex structure on an open subset
or on the fibres of a map, or an endomorphism which fails along a singular
set does not satisfy this hypothesis.
\end{proposition}

\begin{proof}
Assume that \(J\) exists and put \(E=(TM,J)\), a complex vector bundle of
rank six.  Orient \(M\) by the interleaved complex frame
\[
 (v_1,Jv_1,\ldots,v_6,Jv_6),
\]
write \([M]_J\) for the resulting fundamental class, and choose
\(u_J\in H^{12}(M;\Z)\) by
 \(\langle u_J,[M]_J\rangle=1\).  Since \(M\) has the integral cohomology
of \(S^{12}\),
\begin{equation}
 c_1(E)=c_2(E)=c_3(E)=c_4(E)=c_5(E)=0.
 \label{eq:S12-lower-Chern-zero}
\end{equation}
The top-Chern/Euler identity uses exactly this orientation, and
Poincar\'e--Hopf therefore gives
\begin{equation}
 c_6(E)=e(E_{\R})=e(TM)=2u_J.
 \label{eq:S12-top-Chern-two}
\end{equation}
Reversing a previously chosen orientation changes the displayed generator,
not the coefficient two.

For the first contradiction, specialize the original Borel--Serre
calculation
\citep[\S12, Proposition~12.4, Lemma~12.5, and Proposition~12.6,
p.~430; \S15, Proposition~15.1, p.~434]{BorelSerre1953}.
They write \(C_{2j}\) for the universal class now denoted \(c_j\).  At
\(j=6\), the admissible prime is exactly \(p=5<6\), with \(5\nmid6\).
In \(H^*(BU(6);\mathbb F_5)\), pass to a splitting space, write
\(x_1,\ldots,x_6\) for the degree-two Chern roots, and put
\(s_r=\sum_i x_i^r\).  Since
\(\mathcal P_5^1(x_i)=x_i^5\), the Cartan formula and Newton identities
give the complete universal identity
\begin{align}
 \mathcal P_5^1(\overline c_2)
 ={}&\sum_{i<j}(x_i^5x_j+x_ix_j^5)=s_5s_1-s_6 \notag\\
 ={}&6\overline c_6+\overline c_1^4\overline c_2
 -\overline c_1^3\overline c_3
 -4\overline c_1^2\overline c_2^2
 +\overline c_1^2\overline c_4
 +7\overline c_1\overline c_2\overline c_3
 -\overline c_1\overline c_5 \notag\\
 &\quad+2\overline c_2^3-6\overline c_2\overline c_4
 -3\overline c_3^2.
 \label{eq:S12-P1-exact}
\end{align}
Modulo the lower-Chern ideal
\(I=(\overline c_1,\ldots,\overline c_5)\), this says
\[
 \mathcal P_5^1(\overline c_2)
 \equiv6\overline c_6=\overline c_6\pmod I.
\]
Pulling back along the classifying map of \(E\), using
\cref{eq:S12-lower-Chern-zero} and naturality, gives
\[
 \overline{c_6(E)}=\mathcal P_5^1(\overline{c_2(E)})=0.
\]
But \cref{eq:S12-top-Chern-two} gives
\(\overline{c_6(E)}=2\overline u_J\ne0\) in
\(H^{12}(M;\mathbb F_5)\), a contradiction.

For the second contradiction, choose an orientation-preserving homotopy
equivalence \(f:S^{12}\to M\).  Bott integrality says that the sixth Chern
class of every complex bundle on \(S^{12}\) is divisible by
\((6-1)!=120\); see
\href{https://doi.org/10.1007/978-1-4757-2261-1}{Husemoller,
\emph{Fibre Bundles}, Chapter~20, \S9, Theorem~9.6 and
Corollary~9.8, p.~308}.  There is no normalization ambiguity here:
Newton's sixth identity and \cref{eq:S12-lower-Chern-zero} give
\[
 \operatorname{ch}_6(f^*E)
 =\frac{s_6(f^*E)}{6!}
 =-\frac{c_6(f^*E)}{5!}.
\]
Hence \(c_6(f^*E)=120k f^*u_J\) for some \(k\in\Z\), while
\cref{eq:S12-top-Chern-two} gives
\(c_6(f^*E)=2f^*u_J\).  The equation \(120k=2\) has no integral solution.

Finally take the standard smooth sphere
\(S^{12}\subset\R^{13}\).  Its normal field gives the explicit stable
trivialization
\begin{equation}
 TS^{12}\oplus\varepsilon^1_{\R}
 \xrightarrow{\ \cong\ }\varepsilon^{13}_{\R},
 \qquad (x;v,a)\longmapsto(x;v+ax),
 \label{eq:S12-stable-trivialization}
\end{equation}
so \(p_3(TS^{12})=0\).  On the other hand,
\[
 E_{\R}\otimes_{\R}\C\cong E\oplus\overline E,
 \qquad p_3(E_{\R})=-c_6(E\oplus\overline E).
\]
\Cref{eq:S12-lower-Chern-zero} and
\(c_6(\overline E)=c_6(E)\) yield
\begin{equation}
 p_3(TS^{12})=-2c_6(E)=-4u_J,
 \label{eq:S12-Pontryagin-contradiction}
\end{equation}
contradicting stable triviality in the torsion-free group
\(H^{12}(S^{12};\Z)\).  This last calculation is independently recorded
in
\href{https://doi.org/10.1016/j.difgeo.2017.10.010}{Konstantis--Parton,
Lemma~4.1 and Theorem~4.2}.

Each contradiction requires a globally defined \(J\) on the actual bundle
\(TM\).  Stabilizing, deleting points, restricting to regular fibres, or
allowing failure at a singular locus removes the hypothesis rather than
producing an exception.
\end{proof}

\begin{theorem}[Obstructions to stronger globalizations]
\label{thm:obstructions}
The preceding constructions have the following exact limitations.
\begin{enumerate}[label=\textup{(\roman*)}]
\item Let $r\ge1$ and let $N\subset\Lambda^{\oplus r}$ be a fixed
primitive, monodromy-invariant kernel defining a connected fibrewise
torus-group quotient of the native product over the full regular base, or
over a finite cover of it.  Equip it with the quotient norm of the
unscaled marked Euclidean source metric.  If its radial quotient systole
tends to infinity, then
\begin{equation}
 \rank(\Lambda^{\oplus r}/N)\le r.
 \label{eq:rank-bound}
\end{equation}
Four copies are therefore minimal for a full-base, four-dimensional,
all-expanding quotient within this class.  This bound says nothing about
nonlinear maps, varying kernels, rescaled metrics, or non-torus targets.
The $S^{12}$ construction avoids the bound by localizing to $D_0^*$ and
collapsing the exterior.

\item No nontrivial finite base cover makes $K_\Lambda$ a sub-local system
over the whole regular base.  This remains true for a compact cover allowed
to branch over the original singular values.  Moreover, a local cyclic
change $t=s^M$ cannot extend as a connected nontrivial cover of
$\mathbb P^1$ branched only at $p_0$.

\item There is no map from $X^2$, $X\wedge X$, $X^4$, or
$X^{\wedge4}$ to a fixed $T^4$ whose restriction retains the four target
character classes of the corresponding displayed product quotient.  There
is no map to a fixed $S^4$ whose restriction, for even one regular fibre,
is a degree-one collapse of $Q_T^2$ or of $C_T^4$.

\item Cartesian products are not the sphere sources:
\begin{equation}
 X^2\simeq S^6\times S^6\not\simeq S^{12},
 \qquad X^4\simeq(S^6)^4\not\simeq S^{24}.
 \label{eq:products-not-spheres}
\end{equation}
The synchronized two- and fourfold fibre products have real dimensions
$10$ and $18$.  Joins give $X*X\cong S^{13}$ and
$X^{*4}\cong S^{27}$, not the required dimensions.

This product obstruction does not obstruct the preceding smash-source maps.
The global quotient $q_2:(X,x_0)\to(\mathcal A,v_0)$ is based and hence
$q_2\times q_2$ is constant on the global fat wedge
$X\vee X\subset X\times X$.  It therefore induces
$q_2^{\wedge2}:X\wedge X\to\mathcal A\wedge\mathcal A$, whose source is
$S^{12}$.  Since $x_0$ lies on the central fibre, the regular radial product
$F_T^2$ is disjoint from that fat wedge; the product-to-smash quotient is
injective there and the restriction is exactly $\pi_T\times\pi_T$ under the
markings.  This differs from trying to factor the torus map
$\pi_T\times\pi_T$ through $F_T\wedge F_T$ with the fibre identity as
basepoint, which is impossible because the map is nonconstant on that
fibrewise fat wedge.

\item Every connected unbranched finite cover of $X$, of a Cartesian
power $X^r$, or of the smash spheres $X\wedge X$ and $X^{\wedge4}$ is
trivial.  The local unspooling covers are not restrictions of nontrivial
connected unbranched covers of these compact sources.

\item Neither $S^{12}$ nor $S^{24}$, with any differentiable structure,
admits an almost-complex structure.  Hence neither smash source can be
promoted to a compact complex sphere.
\end{enumerate}
\end{theorem}

\begin{proof}
Every fixed element of $K_\Lambda^{\oplus r}$ has uniformly bounded radial
length because $P_T\widehat w=e_1$ and
$P_T\widehat\delta=e_3$.  If such an element were not in \(N\),
primitivity would make its class a nonzero period of the connected
quotient lattice, with quotient norm no greater than its uniformly bounded
source norm.  Thus
\begin{equation}
 K_\Lambda^{\oplus r}\subset N.
 \label{eq:K-in-N}
\end{equation}
For full monodromy, $A_1\widehat w=\widehat u-\widehat w$ and invariance
then put $\widehat u$ in $N$ in every copy.  Consequently
$L_\Lambda^{\oplus r}\subset N$ and
$\rank(\Lambda^{\oplus r}/N)\le4r-3r=r$.  Taking
$N=L_\Lambda^{\oplus4}$ gives the four-circle quotient, so the bound is
sharp.

For the finite-cover statements put
\begin{equation}
 H_*=A_2A_1=
 \begin{pmatrix}
 1&0&0&0\\6&1&1&0\\0&0&1&0\\-5&0&-1&1
 \end{pmatrix}=I+N_*,\qquad N_*^2=0.
 \label{eq:H-star}
\end{equation}
Hence
\begin{equation}
 H_*^n\widehat w=\widehat w+n(\widehat u-\widehat\delta),
 \qquad H_*^n\widehat\delta=\widehat\delta.
 \label{eq:H-star-orbit}
\end{equation}
Every finite-index subgroup of the base fundamental group contains a
positive power of the loop with monodromy $H_*$.  Equation
\cref{eq:H-star-orbit} shows both that it never preserves $K_\Lambda$ and,
combined with \cref{eq:K-in-N}, that $n\widehat u\in N$.  Primitivity of
$N$ implies $\widehat u\in N$, so the same rank bound holds after every
finite base cover.  A branched compact cover restricts over $B^\circ$ to
such a finite cover.  Finally,
$\mathbb P^1\setminus\{p_0\}\cong\C$ is simply connected, so a connected
cover branched only over $p_0$ has degree one.

For a fixed $T^4$ target, each target character class would pull back from
$H^1(X^2;\Z)$, $H^1(X^4;\Z)$, or the first cohomology of the relevant
smash sphere.  Every one of these groups vanishes.  On a radial fibre,
however, a surjective torus homomorphism injects the target character
lattice into the source character lattice; the four displayed target
characters therefore pull back to four independent nonzero classes.  This
contradiction proves the fixed-torus assertion.

For the fixed $S^4$ assertion, let $\omega_T\in H^2(Q_T;\Z)$ be the
orientation class.  The map $\pi_T^*:H^1(Q_T)\to H^1(F_T)$ is injective,
and torus cohomology is the exterior algebra on $H^1$.  Therefore
\begin{equation}
 (\pi_T\times\pi_T)^*(\omega_T\times\omega_T)\ne0
 \quad\text{in }H^4(F_T^2;\Z).
 \label{eq:fibre-H4}
\end{equation}
A degree-one collapse $Q_T^2\to S^4$ pulls the orientation generator back
to $\pm(\omega_T\times\omega_T)$, while
$H^4(X^2;\Z)=H^4(S^{12};\Z)=0$.  The four-circle case is identical: its
orientation class pulls back to the exterior product of four independent
circle characters, but $H^4(X^4;\Z)=H^4(S^{24};\Z)=0$.  The conical
targets evade these contradictions because the fibre orientation and
character classes are not restrictions of fixed global target classes.

The K\"unneth theorem gives $H_6(X^2)=\Z^2$, whereas $H_6(S^{12})=0$,
and gives ranks $4,6,4$ in degrees $6,12,18$ for $X^4$, whereas
$S^{24}$ has none.  An $r$-fold synchronized fibre product has dimension
\[
 6r-2(r-1)=4r+2,
\]
which gives $10$ and $18$.  The standard join identity
$S^m*S^n\cong S^{m+n+1}$ gives the stated join dimensions.

The spaces in the cover assertion are simply connected, so connected
unbranched covers are trivial.  Branched or normalized compactifications
are different objects and have not been constructed here.  Finally,
\cref{prop:S12-almost-complex-obstruction} proves the \(S^{12}\) assertion
by three exact calculations.  Borel and Serre's general theorem applies to
\(S^{24}\), with every differentiable structure, by taking \(n=12\)
\citep[\S15, Proposition~15.1 and the following remark,
p.~434]{BorelSerre1953}.
\end{proof}

\section{The nine-dimensional Gram carrier and the surviving bridge problem}
\label{sec:nine-dimensional-Gram-carrier}

The higher-rung construction that produces a topological \(S^8\) has a
different source and a different quotient mechanism from every smash map
above.  We record its topology here so that the logical implications can be
stated without identifying two constructions merely because their final
boundaries are abstractly homeomorphic spheres.  This is a statement about
what has presently been constructed and proved, not a theorem that no deeper
geometric or physical morphism exists.

\begin{theorem}[The exact \(B^9/S^8\) Gram carrier]
\label{thm:B9-S8-Gram-carrier}
Let
\[
 \mathcal V_3=\operatorname{Mat}_{3\times3}(\C),
 \qquad
 \mathcal P_3=\{H\in\operatorname{Herm}_3(\C):H\succeq0\},
\]
and let \(U(3)\) act on \(\mathcal V_3\) by left multiplication.  The Gram
map
\[
 \Gamma:\mathcal V_3\longrightarrow\mathcal P_3,
 \qquad \Gamma(A)=A^\dagger A,
\]
induces homeomorphisms
\begin{equation}
 \mathcal V_3/U(3)\cong\mathcal P_3,
 \qquad
 \mathcal V_3^+/U(3)\cong\mathcal P_3^+\cong B^9,
 \label{eq:Gram-B9-quotient}
\end{equation}
where \(+\) denotes one-point compactification.  Under the last
homeomorphism,
\begin{equation}
 \partial B^9
 =\{[A]:\rank A\leq2\}\cup\{[\infty]\}
 \cong S^8.
 \label{eq:Gram-rank-loss-S8}
\end{equation}
Thus the carrier is a nine-dimensional ball and its \(S^8\) is the
rank-loss-and-infinity boundary.  The displayed orbit construction is not
a principal \(SU(3)\)-bundle construction.
\end{theorem}

\begin{proof}
The Gram map is invariant because
\((UA)^\dagger(UA)=A^\dagger A\), and it is surjective because
\(\Gamma(H^{1/2})=H\) for \(H\succeq0\).  Suppose
\(A^\dagger A=B^\dagger B\).  The rule
\[
 T:\operatorname{im}A\longrightarrow\operatorname{im}B,
 \qquad T(Ax)=Bx,
\]
is well-defined: \(Ax=0\) implies
\(\|Bx\|^2=x^\dagger B^\dagger Bx=x^\dagger A^\dagger Ax=0\).
The same identity with two vectors shows that \(T\) is a complex-linear
isometry.  Extend \(T\) across the orthogonal complements to a unitary
\(U\in U(3)\); then \(UA=B\).  Thus Gram matrices separate orbits.
Moreover
\begin{equation}
 \|A\|_F^2=\operatorname{tr}(A^\dagger A),
 \label{eq:Gram-properness}
\end{equation}
so \(\Gamma\) is proper.  The induced map from the orbit space is a proper
continuous bijection between locally compact Hausdorff spaces and hence a
homeomorphism.

It remains to identify the compactified cone with a ball and its rank-loss
locus with the boundary.  Put
\[
 \mathcal K_3=\{H\in\mathcal P_3:\operatorname{tr}H=1\},
 \quad
 E=\{X\in\operatorname{Herm}_3(\C):\operatorname{tr}X=0\},
 \quad c=I_3/3.
\]
The real vector space \(E\) has dimension \(8\).  For
\(u\in S(E)\), trace zero and \(u\ne0\) imply
\(\lambda_{\min}(u)<0\); set
\begin{equation}
 R(u)=-\frac1{3\lambda_{\min}(u)}>0.
 \label{eq:Gram-radial-function}
\end{equation}
For \(s\geq0\), the eigenvalues of \(c+su\) are
\(1/3+s\lambda_j(u)\), and therefore
\[
 c+su\succeq0
 \quad\Longleftrightarrow\quad0\leq s\leq R(u).
\]
The continuous maps
\begin{align}
 \chi(c)&=0,
 &\chi(c+su)&=\frac{s}{R(u)}u,
 \label{eq:Gram-K3-ball-map}\\
 \chi^{-1}(0)&=c,
 &\chi^{-1}(tu)&=c+tR(u)u\quad(0<t\leq1)
 \label{eq:Gram-K3-ball-inverse}
\end{align}
are inverse homeomorphisms
\(\mathcal K_3\cong B(E)\cong B^8\).  Continuity at the origin follows
because \(R\) is positive, continuous, and bounded above and below on the
compact sphere \(S(E)\).  The equality \(s=R(u)\) holds exactly on the
singular locus, so \(\chi\) takes that locus to \(S(E)\).

Every nonzero \(H\in\mathcal P_3\) has the unique form
\(H=r\rho\), where \(r=\operatorname{tr}H>0\) and
\(\rho\in\mathcal K_3\).  Put
\[
 t=\frac{r-1}{r+1},
 \qquad
 \mathcal D^9=\{(y,t)\in E\oplus\R:\|y\|_F+|t|\leq1\}.
\]
Define on the one-point compactification
\begin{equation}
 \Xi(r\rho)=((1-|t|)\chi(\rho),t),
 \qquad
 \Xi(0)=(0,-1),
 \qquad
 \Xi(\infty)=(0,1).
 \label{eq:Gram-compactification-map}
\end{equation}
For \(|t|<1\), its inverse is
\begin{equation}
 \rho=\chi^{-1}\!\left(\frac{y}{1-|t|}\right),
 \qquad r=\frac{1+t}{1-t},
 \qquad H=r\rho.
 \label{eq:Gram-compactification-inverse}
\end{equation}
The factor \(1-|t|\) tends uniformly to zero both as \(r\downarrow0\)
and as \(r\to\infty\).  By \cref{eq:Gram-properness}, a set
in \(\mathcal P_3\) is unbounded exactly when its trace is unbounded.
Thus \cref{eq:Gram-compactification-map,eq:Gram-compactification-inverse}
give a homeomorphism \(\mathcal P_3^+\cong\mathcal D^9\).

The latter is the unit ball of the norm
\((y,t)\mapsto\|y\|_F+|t|\) on the nine-dimensional vector space
\(E\oplus\R\), hence is homeomorphic to the Euclidean \(B^9\).  Explicitly,
for \(p=(y,t)\ne0\), the radial formula
\begin{equation}
 p\longmapsto
 \frac{\|y\|_F+|t|}{\sqrt{\|y\|_F^2+t^2}}p
 \label{eq:Gram-double-cone-radial-map}
\end{equation}
and the origin is sent to the origin.  For \(q=(q_E,q_t)\ne0\), the
inverse is
\[
 q\longmapsto
 \frac{\|q\|_2}{\|q_E\|_F+|q_t|}\,q.
\]
It too sends the origin to the origin.
For finite \(H=r\rho\), its image lies on the boundary precisely when
\(\|\chi(\rho)\|_F=1\), equivalently when \(H\) and \(A\) are singular.
The zero matrix and the point at infinity are the two boundary tips, and
zero already belongs to the rank-loss locus.  This proves the second
homeomorphism and \cref{eq:Gram-rank-loss-S8}.

Because \(\Gamma\) is proper, it extends continuously by
\(\Gamma(\infty)=\infty\).  The induced map
\(\mathcal V_3^+/U(3)\to\mathcal P_3^+\) is a continuous bijection from a
compact space to a Hausdorff space, hence a homeomorphism.  This completes
\cref{eq:Gram-B9-quotient}.

Finally, this is a \(U(3)\)-orbit calculation.  The Gram matrix alone does
not separate strict \(SU(3)\)-orbits on the full-rank locus: if
\(B=UA\) with \(A\) invertible, then
\(\det B=(\det U)\det A\).  A strict \(SU(3)\)-invariant must retain the
additional determinant \(z=\det A\), subject to
\(|z|^2=\det(A^\dagger A)\).  On the rank-loss locus the action also has
nontrivial stabilizers.  Hence neither the full quotient nor its boundary
is being asserted to be a principal \(SU(3)\)-bundle.
\end{proof}

\begin{theorem}[Exact obstruction scope and surviving bridges]
\label{thm:Gram-carrier-separation}
Let \(S^8_{\mathrm{Gram}}=\partial\mathcal P_3^+\) denote the boundary in
\cref{thm:B9-S8-Gram-carrier}, and let
\(S^8_{\wedge}=(S^2)^{\wedge4}\) denote the target of
\(\widehat q_\gamma^{\wedge4}\).  Then:
\begin{enumerate}[label=\textup{(\roman*)}]
\item A based homeomorphism
\(h:S^8_{\wedge}\to S^8_{\mathrm{Gram}}\) makes
\(h\circ\widehat q_\gamma^{\wedge4}:S^{24}\to S^8_{\mathrm{Gram}}\)
null-homotopic.  This transfers only the homotopy class of that chosen map;
it does not transfer the Gram orbit relation, rank stratification, a
principal bundle, a connection, or a metric.

\item The nullity of the boundary inclusion
\(S^8_{\mathrm{Gram}}\hookrightarrow B^9\) follows immediately from the
contractibility of \(B^9\).  It does not make the boundary a point and does
not contradict \cref{eq:Gram-rank-loss-S8}.  Likewise, the smash and join
nullities and the Cartesian-product obstructions do not obstruct the
existence of \(B^9\) or its boundary, because the Gram carrier is not
constructed as a Cartesian power, smash power, or join of \(X\).

\item Every principal \(SU(3)\)-bundle on
\(S^8_{\mathrm{Gram}}\cong S^8\) is trivial by
\cref{eq:SU3-low-groups}.  This does not obstruct the orbit space: the
quotient map \(\mathcal V_3^+\to\mathcal V_3^+/U(3)\) has stabilizers on
the rank strata and is neither a principal \(U(3)\)-bundle globally nor a
principal \(SU(3)\)-bundle.

\item If \(S^8_{\mathrm{Gram}}\) is given any smooth structure, the
resulting homotopy eight-sphere admits no almost-complex structure.  This
rules out a complex tangent bundle on the boundary, but not its topology
as the boundary of \(B^9\).

\item Every fixed compatible metric on the compact
\(S^8_{\mathrm{Gram}}\) is bounded, so the unrescaled whole boundary is
not Euclidean \(\R^4\).  For a fixed smooth Riemannian metric, every
blow-up sequence with scale tending to infinity has, after passing to a
basepoint subsequence, the full tangent-space limit \(\R^8\) in pointed
\(C^\infty\) Cheeger--Gromov topology, hence also in pointed
Gromov--Hausdorff topology.  The
\(\R^4\) vertical subspaces and the \(\R^8/\R^{12}\) full bubbles in
\cref{thm:fixed-cusp-metric} belong to the specifically defined conical
smash-target metrics.  Moreover
\cref{thm:scaling-trichotomy} shows that every rescaling which unrolls
one of those slices to \(\R^4\) sends its cone/fat-wedge vertex to
infinite scaled distance, so it is not a vertex tangent cone.  No metric
identifying those bubbles or vertices with the Gram rank boundary has
been specified.
\end{enumerate}
The strongest presently proved bridge is the based homeomorphism \(h\) in
\textup{(i)} and its null composite.  Each negative statement above rules
out only the exact proposed identification or structure to which its proof
applies.  Constructing a stronger bridge remains an active problem: one
would need a specified map from \(X\), a smash/cone target, or its moving
tori to \(B^9\) or \(S^8_{\mathrm{Gram}}\), together with proofs that the
map respects the relevant Gram orbit strata, descends the torus period
data, intertwines any asserted principal-bundle or connection structures,
and compares the fixed metrics in question.
\end{theorem}

\begin{proof}
Composition with a based homeomorphism preserves a based null-homotopy,
so (i) follows from \cref{cor:full-S24-map-null}.  A homeomorphism of
underlying spheres contains none of the extra maps named there.  The first
claim of (ii) is the contraction of an inclusion into a contractible
space.  The remaining claims follow from the definitions: the source of
the Gram construction is the real vector space
\(\operatorname{Mat}_{3\times3}(\C)\), and its equivalence relation is a
left \(U(3)\)-orbit relation, not any of the product, fat-wedge, smash, or
join relations in
\cref{thm:obstructions,cor:smash-join-product-separation}.
Part (iii) is exactly the first group in
\cref{eq:SU3-low-groups}, together with the stabilizer observation at the
end of the proof of \cref{thm:B9-S8-Gram-carrier}.

For (iv), let \(M\) be a smooth manifold homeomorphic to
\(S^8_{\mathrm{Gram}}\), and suppose \(J\) made \(TM\) into a complex
rank-four bundle \(E\).  With \(u\in H^8(M;\Z)\) oriented by \(J\), sphere
cohomology and the top-Chern/Euler identity give
\[
 c_1(E)=c_2(E)=c_3(E)=0,
 \qquad c_4(E)=e(TM)=2u.
\]
The exact relation between Pontryagin and Chern classes gives
\begin{equation}
 p_2(E_\R)=c_2(E)^2-2c_1(E)c_3(E)+2c_4(E)=4u\ne0,
 \label{eq:S8-almost-complex-contradiction}
\end{equation}
while \(p_1(TM)=0\) because \(H^4(M;\Z)=0\).  Also
\(\operatorname{sign}(M)=0\), since \(H^4(M;\R)=0\).  Hirzebruch's
eight-dimensional signature formula would then give
\[
 0=\operatorname{sign}(M)
 =\left\langle\frac{7p_2(TM)-p_1(TM)^2}{45},[M]\right\rangle
 =\frac{28}{45},
\]
a contradiction.  For the standard smooth sphere, the outward normal
gives the independent identity
\(TS^8\oplus\varepsilon^1_\R\cong\varepsilon^9_\R\), and hence
\(p_2(TS^8)=0\), which also contradicts
\cref{eq:S8-almost-complex-contradiction}.  These arguments use only a
hypothetical complex structure on the actual tangent bundle and say
nothing against the topological homeomorphism in
\cref{eq:Gram-rank-loss-S8}.

For (v), compactness makes every compatible metric bounded.  For a smooth
Riemannian metric on a compact \(S^8\), basepoints have a convergent
subsequence, injectivity radius has a positive uniform lower bound, and all
curvature derivatives are uniformly bounded.  Normal coordinates then
show that rescaling distances by factors tending to infinity converges on
each fixed-radius ball to the Euclidean metric on the eight-dimensional
tangent space.  By contrast, the metrics of
\cref{thm:fixed-cusp-metric} were defined from explicit cone and torus
product coordinates and need not be smooth at their cone vertices.
The equality \(\lambda_TD(p_T,*)=\lambda_Tr_T\) in
\cref{eq:vertex-distance} proves the asserted separation between an
\(\R^4\) slice-unrolling limit and a vertex tangent.
No map in the Gram construction equates those metric tensors.  This proves
all five assertions and the final list of open bridges.
\end{proof}

\section{An ancillary finite-cover route}
\label{sec:ancillary-finite-cover}

The following construction is mathematically valid but answers a
different question.  It unrolls the two short periods of the original
four-torus after simultaneous cyclic base change and finite fibre cover.
It does not use the smash sources \(S^{12}\) or \(S^{24}\), does not
produce either sphere carrier, and is not a substitute for
\cref{thm:sphere-carriers,thm:exact-decompactification}.

\begin{proposition}[Monodromy-compatible finite-cover unspooling]
\label{prop:finite-cover-unspooling}
For \(N\geq1\), put
\[
 S_N=\operatorname{diag}(1,1,N,N),\qquad
 \Lambda_N=S_N\Z^4
 =\Z\langle\widehat\gamma,\widehat u,
 N\widehat w,N\widehat\delta\rangle.
\]
The marked radial fibre
\[
 \widetilde F_{T,N}=\R^4/P_TS_N\Z^4
\]
is an \(N^2\)-sheeted holomorphic cover of \(F_T\), with deck group
\((\Z/N\Z)^2\).  After the cyclic cusp-base change \(t=s^N\), these
fibre covers assemble over the punctured cusp because
\begin{equation}
 M_0^N\Lambda_N=\Lambda_N,\qquad
 S_N^{-1}M_0^NS_N=M_0.
 \label{eq:finite-cover-conjugacy}
\end{equation}
Their systoles satisfy
\begin{equation}
 \sys(\widetilde F_{T,N})\geq\min\{T-C,N\}.
 \label{eq:finite-cover-systole}
\end{equation}
Consequently, for every sequence \(T_j,N_j\to\infty\),
\[
 (\widetilde F_{T_j,N_j},[0])
 \longrightarrow(\R^4,0)
\]
in pointed local-isometric, pointed smooth, and pointed
Gromov--Hausdorff topology.
\end{proposition}

\begin{proof}
The inclusion \(P_TS_N\Z^4\subset P_T\Z^4\) makes the identity on
\(\R^4\) descend to a finite holomorphic cover
\(\widetilde F_{T,N}\to F_T\).  Its deck group is
\[
 \frac{P_T\Z^4}{P_TS_N\Z^4}
 \cong\frac{\Z^4}{S_N\Z^4}
 \cong(\Z/N\Z)^2.
\]

Write \(M_0=I+R\).  The identity in
\cref{eq:cusp-monodromy-action} gives
\[
 R(a,b,c,d)^{\mathsf T}=(0,0,b,-a)^{\mathsf T},
 \qquad R^2=0.
\]
One turn around \(s=0\) after \(t=s^N\) winds \(N\) times around the
original cusp, so its monodromy is
\[
 M_0^N=(I+R)^N=I+NR.
\]
In particular,
\[
 M_0^N\widehat\gamma
 =\widehat\gamma-N\widehat\delta,\qquad
 M_0^N\widehat u
 =\widehat u+N\widehat w,
\]
while \(N\widehat w\) and \(N\widehat\delta\) are fixed.  This proves the
first identity in \cref{eq:finite-cover-conjugacy}; expressing the action
in the displayed basis of \(\Lambda_N\) gives the second.

Reorder real coordinates as \((x_1,x_3;x_2,x_4)\) and put
\[
 A_T=
 \begin{pmatrix}
 6\alpha_T&\eta_T\\
 c_T&\alpha_T
 \end{pmatrix}.
\]
Every vector of the covered period lattice has the exact form
\[
 \lambda_{T,N}(n,m)
 =\begin{pmatrix}A_Tn+Nm\\Y_Tn\end{pmatrix},
 \qquad n,m\in\Z^2.
\]
If \(n\ne0\), then
\[
 \|\lambda_{T,N}(n,m)\|
 \geq\|Y_Tn\|\geq(T-C)\|n\|\geq T-C.
\]
If \(n=0\), then \(m\ne0\) and
\(\|\lambda_{T,N}(0,m)\|=N\|m\|\geq N\).  This proves
\cref{eq:finite-cover-systole}.  Once the systole exceeds \(4R\),
\cref{lem:exact-quotient-balls} says that the universal quotient sends
\(\overline B_R(0)\) isometrically onto the entire closed radius-\(R\)
quotient ball.  Letting \(R\) be arbitrary proves all
three pointed convergences.
\end{proof}

\begin{remark}[Why this route is ancillary]
For \(N=2\), the fibre cover has four sheets after a twofold cusp-base
cover.  Sending \(N\to\infty\) fills in the two momentum directions that
remain periodic in the native \(T^2\times\R^2\) limit.  This construction
is useful for analysis on covers, but its source is an auxiliary covered
cusp family.  By contrast, the main theorems of this article derive their
four-dimensional objects from smash powers of the original based
candidate \(X\cong S^6\), yielding the \(S^{12}\) and \(S^{24}\) source
spheres without replacing the native family by a tower of covers.
\end{remark}

\section{Canonicity and exact conclusion}

The lattice
$K_\Lambda=\ker(M_0-I)=\im(M_0-I)$ and the quotient on the punctured
cusp are intrinsic to the marked cusp monodromy.  The compact map $q_2$
also depends on the chosen cusp coordinate, the closed cusp disc, and the
decision to collapse its outer boundary and exterior.  The map
$q_{1,\mathrm{cusp}}$ has the analogous dependence.  The map $q_1$ uses
the invariant character over the full regular base, the two precise
Seifert overlap maps, and the deliberate replacement of the entire cusp
boundary torus by a cone.  The based identifications of the smash powers
with $S^{12}$ and $S^{24}$ depend on the based choice in
\cref{eq:source-sphere}.

Thus \(q_2^{\wedge2}\) is the exact two-copy construction from the native
\(\T^2\times\R^2\) cusp limit.  Its expanding quotient is cusp-local, but
its source is the global real smash sphere \(S^{12}\).  The map
\(q_{1,\mathrm{cusp}}^{\wedge4}\) is the corresponding localized
one-circle construction.  The stronger \(q_1^{\wedge4}\) uses the unique
primitive invariant character over every regular fibre, the two exact
finite Seifert fillings, and a cone at the nonnormal cusp.  Two copies are
minimal for the cusp-local rank-two route, while four copies are minimal
for a full-base rank-four expanding quotient, within the fixed-primitive-
kernel, connected native torus-group quotients carrying the unscaled
Euclidean quotient norm.  No broader minimality is claimed.

Further quotienting gives the honest sphere maps
\[
 S^{12}\twoheadrightarrow S^6,\qquad S^{24}\twoheadrightarrow S^8.
\]
Both full maps are based null-homotopic.  For
\(\widehat q_\gamma^{\wedge4}\), the smash square of
\(\widehat q_\gamma\) vanishes because
\(\pi_{12}(S^8)\cong\pi_4^{\mathrm S}=0\).  For
\(\widehat q_3^{\wedge2}\), the stronger external-pairing calculation uses
Toda's formula
\([x\wedge y]=(E^3[x])\circ(E^6[y])\).  Since
\(\pi_6(S^3)=\Z/12\), the class \(E^6[y]\in\pi_{12}(S^9)=\Z/24\) is
even.  Postcomposition by \(E^3[x]\) is a homomorphism into
\(\pi_{12}(S^6)=\Z/2\), so it sends this even class to zero.  Thus
\(x\wedge y\simeq_*0\) for every based pair \(x,y:S^6\to S^3\), and all
higher smash powers vanish as well.  Their reduced joins vanish because
join is suspended smash.  Their Cartesian product maps vanish if and only
if the corresponding one-copy map vanishes. The explicit framed filling
now proves product nullity for \(g=\widehat q_3\); it remains open for
\(f=\widehat q_\gamma\).

The one-copy class of \(\widehat q_3\) is now zero by
Theorem~\ref{cone-cusp-null:thm}; that of \(\widehat q_\gamma\) remains open.
The general framing comparison retains the following exact structure.  Its regular inverse image is \(T^3\)
with the honest normal framing supplied by the differential.  Direct-limit
stable framing components are indexed first by the spin torsor
\(H^1(T^3;\Z/2)\), and within a fixed spin component by the affine
\(H^3(T^3;\Z)\)-coordinate.  An honest local rank-three twist maps to twice
the generator of \(\pi_3(SO)\) and changes the stable Pontryagin--Thom
class by \(2\nu\).  Thus computing that relative integer modulo \(12\),
not merely recognizing the unframed torus or its spin structure, computes
\([\widehat q_3]\in\pi_6(S^3)\). The explicit cap in
Theorem~\ref{cone-cusp-null:thm} carries out that computation and gives zero.

The based homotopy classes of
\(q_2^{\wedge2}\), \(q_{1,\mathrm{cusp}}^{\wedge4}\), and
\(q_1^{\wedge4}\) into their conically stratified targets remain open:
postcomposition with a sphere quotient cannot be reversed.  Surjectivity
and the exact moving-slice restrictions survive the sphere-valued
homotopical vanishing.
Their targets contain moving embedded four-spheres whose intrinsic
quotient metrics have exact Euclidean balls of every fixed radius
sufficiently far into the cusp.  Removing the smash point from each such
sphere gives a nonclosed, topologically embedded \(\R^4\) cell.
Separately, the native radial four-tori admit the noncompact families
\(\mathfrak E_{12}\) and \(\mathfrak E_{24}\), each a five-manifold with
boundary \(\R^4\).  The continuous maps \(\beta_\nu\) are homeomorphisms
off that boundary and send it to the cone vertex, but they are nonproper
and nonquotient; they are not topological blow-downs or resolutions.  In a
fixed cusp-adapted target metric, the same \(\R^4\)
appears as a vertical subspace of a moving-point bubble, but the full
ambient bubbles have dimensions \(8\) and \(12\).

The \(SU(3)\) clutching calculation is equally sharp.  A class-\(k\)
bundle on a moving \(S^4\) has second Chern number \(k\), and its pullback
to the radial \(T^4\) has the same number.  When \(k\ne0\), it cannot extend
over the ambient sphere or the contractible cone carrier.  Conversely,
every bundle pulled back through either null full sphere map is trivial.
These statements neither construct nor obstruct the currently separately
constructed Gram orbit space
\[
 \operatorname{Mat}_{3\times3}(\C)/U(3)
 \cong\operatorname{Herm}_3^+(\C),
 \qquad
 \operatorname{Herm}_3^+(\C)^+\cong B^9,
\]
whose rank-deficient-and-infinity boundary is \(S^8\).  That \(S^8\) is
abstractly homeomorphic to the fourfold-smash target, and a chosen based
homeomorphism gives a proved null composite from \(S^{24}\).  No
structure-preserving refinement of that homeomorphism is presently
constructed.  The Gram carrier is not
constructed as a product, smash, or join of \(X\); its orbit map is not a
principal \(SU(3)\)-bundle; and no period-lattice, connection, or fixed
metric comparison map has been given.  Its boundary inclusion into \(B^9\)
is null because the ball is contractible, but the boundary remains an
\(S^8\).  No smoothing of this homotopy \(S^8\) is almost complex, which
excludes exactly that additional tangent-bundle structure and not the
topological \(B^9/S^8\) carrier or a different future bridge.

These statements do not produce a compact complex \(S^{12}\) or
\(S^{24}\), an unrescaled ambient \(\R^4\) limit of a compact target, a
Riemannian metric on a source sphere, or a physical spacetime.  The
finite-cover tower of \cref{sec:ancillary-finite-cover} also has pointed
limit \(\R^4\), but it is logically separate from the sphere route.  A
global compactification of that auxiliary tower would require a specified
global base cover, normalization and resolution of all pulled-back
singular fibres, a lifted metric, and a new topology calculation; none is
asserted here.  Likewise, any proposed bridge to the Gram carrier still
requires an explicit map respecting its rank strata and a separate proof
of every bundle and metric compatibility claimed for that map.  Its absence
is the present proof boundary; it is not an ontological conclusion about
the carrier or its possible physical interpretation.

\section{The original cusp coefficient: an explicit framed nullbordism}


Let \(X\) be the smooth compact source of the original construction,
\(f:X\to B\) its map, and \(t\) the specified cusp coordinate.
Put \(W=f^{-1}(0)\), choose the original \(x_0\in W\), and retain the
ordered quotient lattice basis induced by
\((\widehat\gamma,\widehat u)\), with kernel generated by
\((\widehat w,\widehat\delta)\). Thus the quotient torus is
\(\overline T^2=\mathbb R^2/\mathbb Z^2\). The original periods are
\begin{equation}
 Z(s)=sB_0+C(t),\qquad e^{2\pi i s}=t,\qquad
 B_0=\begin{pmatrix}0&1\\-1&0\end{pmatrix},\qquad
 C(t)=\begin{pmatrix}6\mu(t)&h(t)\\b(t)-h(t)&\mu(t)\end{pmatrix}.
 \label{cone-cusp-null:eq:periods}
\end{equation}
Every entry of \(C\) is holomorphic at \(t=0\).
The dense torus of the original toric cusp has coordinates
\((x_1,x_2,t)\), with
\(x=\exp(2\pi i\zeta)\). The residual lattice action is exactly
\begin{equation}
 \Psi_\lambda(x,t)=
 \bigl(t^{B_0\lambda}\exp(2\pi iC(t)\lambda)x,t\bigr),
 \qquad\lambda\in\mathbb Z^2.
 \label{cone-cusp-null:eq:action}
\end{equation}
It extends freely and properly discontinuously across the original
unimodular toric filling. In particular the quotient projection is a local
biholomorphism at every toric fixed point. These are the reconstructed
data and proved local properties of the original source
\cite[\S4, source pp.~24--29]{AlpogeS6_2026}.

Let
\[
 K^{(1)}=(\mathbb R/\mathbb Z\times\{0\})\cup(\{0\}\times\mathbb R/\mathbb Z),\qquad
 c:\overline T^2\longrightarrow \overline T^2/K^{(1)}\cong S^2
\]
be the specified degree-one cell quotient. Its restriction to the open
two-cell is taken with the smooth coordinates used in the original
regular-fibre description. Both radial endpoints and the complete
basepoint cylinder are collapsed in
\[
 \Sigma_r(S^2,*)=
 \frac{[0,\epsilon]\times S^2}
 {(\{0,\epsilon\}\times S^2)\cup([0,\epsilon]\times\{*\})}
 \cong S^3.
\]
The original based map is
\begin{equation}
 g(x)=[|t|,c(\overline\pi_f(x))]\quad(0<|t|<\epsilon),\qquad
 g(W)=g(X\setminus\operatorname{int}N_0)=*.
 \label{cone-cusp-null:eq:g}
\end{equation}

\begin{theorem}[The original one-copy cusp map is null-homotopic]
\label{cone-cusp-null:thm}
The based map \(g:(X,x_0)\to(S^3,*)\) in \eqref{cone-cusp-null:eq:g} is based
null-homotopic. More precisely, a regular inverse image with its actual
ordered normal framing induced by \((r,c_1,c_2)\) is framed nullbordant
inside one original cusp chart times \([0,1]\), disjoint from
\(\{x_0\}\times[0,1]\).
\end{theorem}

\begin{proof}
The proof constructs the complete inverse image, compares its actual
framing with an explicitly extendible framing, and gives the extension.

\paragraph{1. The exact complete inverse image.}
Retain the lift
\[
 u^0=\begin{pmatrix}-1/4\\1/4\end{pmatrix},\qquad
 \overline u^0=(3/4,1/4)\pmod{\mathbb Z^2},\qquad
 B_0u^0=\begin{pmatrix}1/4\\1/4\end{pmatrix}.
\]
Thus \(\overline u^0\notin K^{(1)}\), and \(B_0u^0\) lies in the
interior of the original triangle \(\operatorname{conv}(0,e_1,e_2)\).
Use its original unimodular chart
\begin{equation}
 z_0=t/(x_1x_2),\qquad z_1=x_1,\qquad z_2=x_2,
 \qquad t=z_0z_1z_2.
 \label{cone-cusp-null:eq:chart}
\end{equation}
Set
\[
 \ell=\log|t|<0,\qquad R(t)=-2\pi\operatorname{Im}C(t),\qquad
 A(t)=\ell B_0+R(t).
\]
For a representative \(\zeta=Z(s)u+k\), where \(u,k\in\mathbb R^2\),
\(\operatorname{Im}s=-\ell/(2\pi)\) gives the exact identity
\begin{equation}
 \begin{pmatrix}\log|x_1|\\\log|x_2|\end{pmatrix}
 =A(t)u.
 \label{cone-cusp-null:eq:moduli}
\end{equation}
The quotient value is \(u\bmod\mathbb Z^2\).
Since \(B_0\) is orthogonal, for all \(v\in\mathbb R^2\),
\begin{equation}
 \|A(t)v\|\geq (|\ell|-\|R(t)\|)\|v\|.
 \label{cone-cusp-null:eq:invertibility}
\end{equation}
The boundedness of \(R\) therefore makes \(A(t)\) invertible on a
sufficiently small punctured disc.

For any point of the fibre quotient over \(\overline u^0\), subtracting
the appropriate original period, equivalently applying
\(\Psi_{-\lambda}\), gives a representative with the exact lift
\(u=u^0\). This representative is unique: another residual translate
has logarithmic moduli differing by \(A(t)\lambda\), which vanishes
only for \(\lambda=0\). Conversely every dense-torus point satisfying
\(\log|x|=A(t)u^0\) has this quotient value, since the remaining real
parts of \(\zeta\) are exactly the kernel coordinates \(k\).
For each fixed \(t\), varying \(k\bmod\mathbb Z^2\) gives all independent
phases of \(x_1,x_2\).

The following change of coordinates uses the full original period units:
\begin{align}
 h_j^0(t)&=\exp\bigl(-2\pi i(C(t)u^0)_j\bigr),\qquad j=1,2,\\
 \mathcal F(z_0,z_1,z_2)
 &=(w_0,w_1,w_2)
 =\left(\frac{z_0}{h_1^0(t)h_2^0(t)},\,
 h_1^0(t)z_1,\,h_2^0(t)z_2\right).
 \label{cone-cusp-null:eq:straightening}
\end{align}
The functions \(h_j^0\) are nowhere-zero holomorphic units, including
at \(t=0\). The product \(w_0w_1w_2\) is exactly \(t\).
The inverse of \(\mathcal F\) is
\[
 (w_0,w_1,w_2)\longmapsto
 \left(w_0h_1^0(t)h_2^0(t),\,
 \frac{w_1}{h_1^0(t)},\,\frac{w_2}{h_2^0(t)}\right),
 \qquad t=w_0w_1w_2.
\]
Thus this is a biholomorphism of the original chart restricted to the
cusp disc, not a replacement of \(C(t)\).

For \(r_0\in(0,\epsilon)\), the entire inverse image of
\(y_0=[r_0,c(\overline u^0)]\) becomes exactly
\begin{equation}
 T_{r_0}=
 \{|w_0|=\rho_0,\ |w_1|=\rho_1,\ |w_2|=\rho_2\},\qquad
 (\rho_0,\rho_1,\rho_2)=
 (r_0^{1/2},r_0^{1/4},r_0^{1/4}).
 \label{cone-cusp-null:eq:torus}
\end{equation}
Indeed, \(\log|h^0(t)|=2\pi\operatorname{Im}C(t)u^0=-R(t)u^0\),
so \eqref{cone-cusp-null:eq:moduli} gives \(\log|w_j|=\ell(B_0u^0)_j\) for
\(j=1,2\); the product identity gives the exponent \(1/2\) for
\(w_0\). Conversely these moduli recover \(u=u^0\) by invertibility
of \(A(t)\). The angular map is explicitly
\[
 (\arg t,\arg w_1,\arg w_2)\longleftrightarrow
 (\arg w_0,\arg w_1,\arg w_2),\qquad
 \arg w_0=\arg t-\arg w_1-\arg w_2\pmod{2\pi},
\]
with inverse \(\arg t=\arg w_0+\arg w_1+\arg w_2\).
Its integral matrix has determinant one. This includes the complete
angular mapping torus and all kernel phases, with their original signs.

Let \(o\) be the origin of \eqref{cone-cusp-null:eq:chart}. The original free properly
discontinuous action supplies a neighborhood \(V\) of \(o\) on which
the quotient projection is injective. To see the local injectivity
directly, first choose a relatively compact neighborhood; proper
discontinuity leaves only finitely many translates that can intersect it.
None of those nonidentity translates fixes \(o\). Shrinking finitely
many times separates the neighborhood from all of them. The quotient
map is a local biholomorphism, so its restriction to \(V\) is a chart
of the original source.
Choose a polydisc about \(0\) whose closure lies in \(\mathcal F(V)\).
For sufficiently small \(r_0\), both \(T_{r_0}\) and
\begin{equation}
 \{|w_0|\leq\rho_0,\ |w_1|=\rho_1,\ |w_2|=\rho_2\}
 \label{cone-cusp-null:eq:spatial-bound}
\end{equation}
lie in this polydisc. This proves that the whole inverse image embeds
in one chart and supplies space for the bounding manifold below.

\paragraph{2. The actual normal framing, including its angle dependence.}
Put \(L_j=\log|w_j|\), and use the unique real lift of the quotient
coordinate near \(u^0\). From the exact change of coordinates,
\begin{equation}
 \begin{pmatrix}L_1\\L_2\end{pmatrix}
 =\ell B_0u^0+A(t)(u-u^0),\qquad
 L_0=\ell-L_1-L_2.
 \label{cone-cusp-null:eq:nearby-moduli}
\end{equation}
Differentiating on \(u=u^0\) annihilates the term
\((dA(t))(u-u^0)\). Hence no derivative of \(C(t)\) has been omitted:
the exact differential there is
\begin{equation}
 \begin{pmatrix}dL_0\\dL_1\\dL_2\end{pmatrix}
 =J(t)\begin{pmatrix}d\ell\\du_1\\du_2\end{pmatrix},\quad
 J(t)=
 \begin{pmatrix}
 1/2&-A_{11}-A_{21}&-A_{12}-A_{22}\\
 1/4&A_{11}&A_{12}\\
 1/4&A_{21}&A_{22}
 \end{pmatrix}.
 \label{cone-cusp-null:eq:J}
\end{equation}
Adding the second and third rows to the first gives \((1,0,0)\),
and expansion in that row proves \(\det J(t)=\det A(t)\).

Let \(D=Dc_{\overline u^0}\) be the actual two-by-two derivative in
the ordered open-cell target coordinates \((c_1,c_2)\). It is a fixed
invertible matrix; all its entries and signs are retained. Let
\[
 H=\begin{pmatrix}r_0&0&0\\0&D_{11}&D_{12}\\0&D_{21}&D_{22}\end{pmatrix}.
\]
Let \(n_j\) be the real radial vector \(w_j\) in the \(j\)-th complex
coordinate, with other components zero. Then \(dL_i(n_j)=\delta_{ij}\).
Writing \(N=(n_0,n_1,n_2)\) as the matrix with these vector columns, the actual
\(dg\)-dual normal framing is, in the normal quotient bundle,
\begin{equation}
 \mathfrak n(t)=N J(t)H^{-1}.
 \label{cone-cusp-null:eq:actual-frame}
\end{equation}
Indeed, \(dg(N)=HJ(t)^{-1}\), so multiplication by \(J(t)H^{-1}\)
gives the identity. In the Euclidean metric of this chart the radial
vectors are perpendicular to the product torus; a different source
metric changes their representatives by the canonical projection to its
normal bundle and preserves this normal-quotient framing.

Put \(\ell_0=\log r_0\). Decrease \(r_0\) further, if necessary, so
that \(|\ell_0|>2\sup_{|t|\leq r_0}\|R(t)\|\). For
\(0\leq\sigma\leq1\), define only a homotopy of frames by
\[
 A_\sigma(t)=\ell_0B_0+\sigma R(t),
\]
and let \(J_\sigma(t)\) be \eqref{cone-cusp-null:eq:J} with \(A_\sigma\) in place
of \(A\). For every \(v\in\mathbb R^2\),
\[
 \|A_\sigma(t)v\|\geq
 (|\ell_0|-\|R(t)\|)\|v\|>0\quad(v\ne0).
\]
Thus all matrices are invertible. At \(\sigma=0\),
\begin{equation}
 J_0=
 \begin{pmatrix}
 1/2&\ell_0&-\ell_0\\
 1/4&0&\ell_0\\
 1/4&-\ell_0&0
 \end{pmatrix},\qquad \det J_0=\ell_0^2>0.
 \label{cone-cusp-null:eq:J0}
\end{equation}
Continuity also proves \(\det J_\sigma(t)>0\) for every \(\sigma,t\).
The homotopy \(N J_\sigma(t)H^{-1}\), run from \(\sigma=1\) to
\(\sigma=0\), compares the actual frame with the exact constant-matrix
frame \(N P\), where
\begin{equation}
 P=J_0H^{-1}.
 \label{cone-cusp-null:eq:P}
\end{equation}
The original embedded torus and original functions \(C(t)\) have not
changed. A framing homotopy is a framed cylinder: choose a smooth
parameter function constant near its endpoints and use the displayed
frames on \(T_{r_0}\times[0,1]\).

\paragraph{3. A collared framed nullbordism with the exact constant matrix.}
Use a new time coordinate \(v\in[0,1]\). Define
\begin{equation}
 q(v)=
 \begin{cases}
 0,&0\leq v\leq1/4,\\
 \rho_0^2\exp\bigl(4-1/(v-1/4)\bigr),&1/4<v\leq1.
 \end{cases}
 \label{cone-cusp-null:eq:q}
\end{equation}
This function is smooth: all right derivatives of
\(\exp(-1/(v-1/4))\) at \(1/4\) vanish. It is strictly increasing
for \(v>1/4\), satisfies \(q(1/2)=\rho_0^2\), and
\(q'(1/2)=16\rho_0^2>0\).
In the chart times \([0,1]\), put
\begin{equation}
 \mathcal W=
 \{|w_0|^2+q(v)=\rho_0^2,\quad
       |w_1|^2=\rho_1^2,\quad|w_2|^2=\rho_2^2\}.
 \label{cone-cusp-null:eq:bordism}
\end{equation}
No point has \(v>1/2\). The three defining differentials are independent:
the latter two have nonzero gradients in their separate complex factors;
the first has nonzero \(w_0\)-gradient whenever \(w_0\ne0\), and at
\(w_0=0\) it has the nonzero \(v\)-derivative \(q'(1/2)\).
Thus \(\mathcal W\) is a smooth compact four-manifold. Its only
boundary is \(T_{r_0}\times\{0\}\), and on \(0\leq v\leq1/4\)
it is the exact product collar \(T_{r_0}\times[0,1/4]\).
The north locus \(w_0=0,v=1/2\) is an interior cap: the implicit
function theorem expresses \(v\) smoothly in terms of
\((\operatorname{Re}w_0,\operatorname{Im}w_0)\) there. The first
factor is a disk obtained by rotating its radial profile with this smooth
cap and one boundary circle. Therefore
\(\mathcal W\cong D^2\times T^2\).
Its spatial coordinates lie in \eqref{cone-cusp-null:eq:spatial-bound}, so the complete
manifold lies in the injective original chart.

In the Euclidean chart product metric define
\[
 V_0=(w_0,0,0,q'(v)/2),\qquad
 V_1=(0,w_1,0,0),\qquad V_2=(0,0,w_2,0).
\]
These are one half of the gradients of the three defining equations
in \eqref{cone-cusp-null:eq:bordism}. They are mutually orthogonal and never zero.
For example, if \(F_0=|w_0|^2+q(v)-\rho_0^2\), then
\[
 dF_0(V_0)=2|w_0|^2+\frac12q'(v)^2>0
\]
at every point of \(\mathcal W\). On the entire collar \(q'=0\), so
\((V_0,V_1,V_2)=(N,0)\) exactly. Consequently
\begin{equation}
 (V_0,V_1,V_2)P
 \label{cone-cusp-null:eq:extended-frame}
\end{equation}
is a framing of the normal three-plane bundle of \(\mathcal W\),
and its restriction to the collar is exactly the framing \(NP\) from
\eqref{cone-cusp-null:eq:P}, with every original target-coordinate sign retained.
Concatenate the framing cylinder of the preceding paragraph with this
nullbordism, rescaling the two interval lengths. Their framings agree
exactly on the common product collar, so the concatenation is a smooth
framed nullbordism of the actual inverse image.

\paragraph{4. The based map represented by this framed inverse image.}
Choose the injective chart to exclude \(x_0\) if its center is not
\(x_0\). If its center is \(x_0\), the constructed nullbordism still
excludes \(x_0\), since \(|w_1|=\rho_1>0\) and
\(|w_2|=\rho_2>0\) everywhere on it. Compactness therefore gives
tubular neighborhoods for the framed cylinder and cap disjoint from
\(\{x_0\}\times[0,1]\).

The original continuous quotient map is smooth and a submersion near
its entire inverse image of \(y_0\). First make it constant in a
neighborhood of \(x_0\), by a based homotopy that changes nothing near
that inverse image. Explicitly, take a target chart
\(\kappa:O\to\mathbb R^3\) with \(\kappa(*)=0\), whose image is a ball,
and \(y_0\notin O\). By continuity choose a neighborhood \(U\) of
\(x_0\) with \(g(U)\subset O\), and a smooth bump function
\(\beta\) supported in \(U\) and equal to one near \(x_0\).
For \(x\in U\), replace \(g(x)\) during this homotopy by
\(\kappa^{-1}((1-a\beta(x))\kappa(g(x)))\), \(0\leq a\leq1\),
and leave \(g\) fixed off \(U\). Convexity of the target chart image
makes the formula well-defined, the vanishing of \(\beta\) near the
boundary of its support proves continuity of the pasted formula, and
\(\kappa(g(x_0))=0\) makes the homotopy based. It never introduces
an inverse image of \(y_0\). Its endpoint is smooth near \(x_0\) and
agrees with \(g\) near the entire old inverse image.

Denote the endpoint of the preceding homotopy by \(g_b\). Choose one
closed tubular neighborhood \(U_1\) of the inverse image, contained in
the smooth submersion region and with the entire inverse image contained
in \(\operatorname{int}U_1\). The compact set
\(X\setminus\operatorname{int}U_1\) has image under \(g_b\) missing
\(y_0\). In a fixed target metric put
\[
 d_0=\min_{x\in X\setminus\operatorname{int}U_1}
       d(g_b(x),y_0)>0.
\]
Relative smooth approximation, fixed on the whole closed set \(U_1\)
and on a neighborhood of \(x_0\), gives a smooth based representative
\(g_{\mathrm{sm}}\) whose error on
\(X\setminus\operatorname{int}U_1\) is less than \(d_0/2\).
It therefore has the same inverse image and differential on \(U_1\),
and \(d(g_{\mathrm{sm}}(x),y_0)>d_0/2\) on its complement. Thus there
is no uncontrolled intervening region and no additional inverse image.
Make the approximation also closer than a fixed target injectivity
radius; its short geodesics give the based homotopy.

The relative Pontryagin construction now applies to the explicit
framed nullbordism. It produces a based homotopy to a map missing
\(y_0\); stereographic contraction of \(S^3\setminus\{y_0\}\cong\mathbb R^3\)
to the image of \(*\) is based and contracts that last map.
For completeness, the dependence on the framed inverse image in this
argument is precisely the local product-neighborhood construction and
the interpolation lemma of \cite[\S7, Theorem B and Lemma 4,
pp.~43, 48--50]{Milnor1965Cusp}. Their relative versions fix the source
basepoint because all tubular modifications are away from it and
stereographic interpolation between equal values keeps that value
fixed. A framed cylinder realizes a framing homotopy; the empty
Pontryagin manifold realizes a map missing the chosen regular value.
This proves the claimed based null-homotopy of the exact map \(g\).
\end{proof}

\begin{corollary}
Under the original fixed based identification \(X\cong S^6\), the
one-copy coefficient is
\[
 [g]=0\in\pi_6(S^3)\cong\mathbb Z/12,\qquad a=0\pmod{12},\qquad
 [E^6g]=0\in\pi_{12}(S^9)\cong\mathbb Z/24.
\]
The previously defined homomorphism
\(\Theta_g:\mathbb Z/12\to\pi_6(S^3)\), \(\Theta_g([n])=n[g]\), and
its composite with the marked peripheral character are zero.
\end{corollary}
\begin{proof}
The theorem states that the actual based map is null-homotopic.
Precomposition with the fixed based identification preserves a based
null-homotopy. Suspension preserves it as well. Thus the unique
coefficient of \([g]\) in any generator of \(\mathbb Z/12\) is zero.
For every \([n]\), \(\Theta_g([n])=n\cdot0=0\), and every
composite with this homomorphism is zero.
\end{proof}

\paragraph{Exact scope.}
The proof computes this sphere-valued postcomposition. It does not
compute the homotopy class of the earlier conical-target map
\(q_2:X\to\operatorname{Susp}(S^1\times\overline T^2)\), and it
does not assert a null-homotopy for that different target. Point-set
surjectivity of \(g\) and all of its stated metric restrictions remain
properties of the original map: a homotopy-class computation does not
alter the defining formula \eqref{cone-cusp-null:eq:g}.


% BEGIN RETAINED_ANGLE_AND_HOPF_BRIDGE_20260905

\section{The retained-angle cusp quotient and its nonzero top-cell class}
\label{ret:section}

\paragraph{The exact object.}
The notation refers to the specified threefold \(X\), its fixed based
smooth identification with \(S^6\), and its cusp filling from
\cite[\S\S3.7,4.1--4.5]{retbib:Canonical}.  The native quotient and its compact
extension are the maps in \cite{retbib:ConeRoutes}; their full formulas are
recalled below.  This calculation uses the original period matrix and
native period lattices.  It computes a new postcomposition of that
particular quotient; it does not use the homotopy class of the
angle-forgetting map \(g:X\to S^3\).

\subsection{Original coordinates and the specified sphere-valued map}

The ordered quotient lattice is
\[
 \overline\Lambda=\mathbb Z\langle\overline{\widehat\gamma},
              \overline{\widehat u}\rangle,
 \qquad
 K_\Lambda=\mathbb Z\langle\widehat w,\widehat\delta\rangle,
 \qquad
 \overline T^2=\mathbb R^2/\mathbb Z^2.
\]
In these ordered bases the original period data are
\begin{equation}
 B_0=\begin{pmatrix}0&1\\-1&0\end{pmatrix},\qquad
 Z(s)=sB_0+C(t),\qquad t=e^{2\pi i s},\qquad
 C(t)=\begin{pmatrix}6\mu(t)&h_{\rm per}(t)\\
 b(t)-h_{\rm per}(t)&\mu(t)\end{pmatrix}.
 \label{ret:eq:original-periods}
\end{equation}
Here \(b=\beta+\tau\), \(s=\tau-h_{\rm per}(t)\), and every entry of
\(C\) is holomorphic at \(t=0\), exactly as in the original source.
The subscript in \(h_{\rm per}\) distinguishes that period function
from the sphere-valued map below.  Write
\begin{equation}
 R(t)=-2\pi\operatorname{Im}C(t),\qquad
 \ell=\log|t|,\qquad A(t)=\ell B_0+R(t).
 \label{ret:eq:A}
\end{equation}
Choose a closed cusp disc \(|t|\leq\epsilon\) inside the original
coordinate disc and small enough for the quotient construction.  On a
slightly larger closed disc put \(M=\sup\|R(t)\|<\infty\).  We will
choose a regular value with \(0<r_0<\epsilon\) and
\(|\log r_0|>M\).  This chooses a point of the already specified map;
it does not change that map or the cutoff radius \(\epsilon\).

The original covering model on its dense torus is
\begin{equation}
 x_j=e^{2\pi i\zeta_j},\qquad
 \Psi_n(x,t)=
 \bigl(e^{2\pi i C(t)n}t^{B_0n}x,t\bigr),\qquad n\in\mathbb Z^2.
 \label{ret:eq:deck}
\end{equation}
For a branch of \(s\), writing
\(\zeta=Z(s)u+k\) with \(u,k\in\mathbb R^2\) gives
\begin{equation}
 \log|x|=A(t)u.
 \label{ret:eq:logx}
\end{equation}
Indeed, \(\operatorname{Im}s=-\ell/(2\pi)\), so
\(-2\pi\operatorname{Im}(Z(s)u+k)=\ell B_0u+R(t)u\).
Since \(B_0\) is orthogonal,
\begin{equation}
 \bigl|\bigl(\ell B_0+vR(t)\bigr)\xi\bigr|
 \geq (|\ell|-vM)|\xi|>0
 \quad(0\leq v\leq1,\ \xi\ne0)
 \label{ret:eq:invertibility}
\end{equation}
whenever \(|\ell|>M\).  Thus \(A(t)\) is invertible there.  Equation
\eqref{ret:eq:deck} changes \(A(t)^{-1}\log|x|\) by \(n\), so the native
quotient is exactly
\begin{equation}
 \overline\pi_f([x,t])=[A(t)^{-1}\log|x|]\in\mathbb R^2/\mathbb Z^2.
 \label{ret:eq:native}
\end{equation}
This also follows from the flat marking in \cite{retbib:ConeRoutes}.

Let
\[
 T^3=S^1_\phi\times\overline T^2,
 \qquad S^1_\phi=\mathbb R/(2\pi\mathbb Z),
\]
with coordinate basepoints \(\phi=0\), \(u_1=0\), \(u_2=0\).
Its coordinate two-skeleton is the following specified closed subset:
\begin{equation}
 K^{(2)}=
 \bigl(\{0\}\times\overline T^2\bigr)
 \cup\bigl(S^1_\phi\times\{u_1=0\}\bigr)
 \cup\bigl(S^1_\phi\times\{u_2=0\}\bigr).
 \label{ret:eq:two-skeleton}
\end{equation}
The product CW structure has one open three-cell, with coordinates
\((\phi,u_1,u_2)\in(0,2\pi)\times(0,1)^2\).  Consequently
\[
 c_3:T^3\longrightarrow T^3/K^{(2)}\cong S^3
\]
is the degree-one quotient, where the orientation of \(S^3\) is fixed
by the displayed coordinate order.  One explicit identification of its
open cell with \(\mathbb R^3\) is
\[
 (\phi,u_1,u_2)\longmapsto
 \bigl(-\cot(\phi/2),-\cot(\pi u_1),-\cot(\pi u_2)\bigr).
\]
Each derivative in its own variable is strictly positive, and approach
to any boundary face leaves every compact subset of \(\mathbb R^3\).
Extending by its point at infinity gives the stated based sphere
identification and verifies its continuity on the collapsed skeleton.

The exact conical target and map are
\begin{align}
 \mathcal A&=
 \frac{[0,\epsilon]\times T^3}
 {\{0\}\times T^3\sim v_0,\quad
  \{\epsilon\}\times T^3\sim v_1},
 \label{ret:eq:A-target}\\
 q_2(x)&=[r,e^{i\phi},\overline\pi_f(x)]
       &&(t(f(x))=re^{i\phi},\ 0<r<\epsilon),\notag\\
 q_2(W)&=v_0,\qquad
 q_2\bigl(X\setminus\operatorname{int}N_0\bigr)=v_1.
 \label{ret:eq:q2}
\end{align}
Fix the same source basepoint \(x_{\rm base}\in W\).  Define
\begin{equation}
 p_3:\mathcal A\longrightarrow\Sigma_r(S^3,*),\qquad
 p_3([r,z])=[r,c_3(z)],\qquad
 H=p_3q_2:(X,x_{\rm base})\longrightarrow(S^4,*).
 \label{ret:eq:H}
\end{equation}
The reduced suspension here collapses both endpoint copies of \(S^3\)
and the entire interval over \(*\) to one point.  Thus \(p_3\) sends
both \(v_0\) and \(v_1\) to \(*\) and is well-defined and based.
It is continuous by the defining quotient topology.  An explicit
identification of its open four-cell with \(\mathbb R^4\) is
\begin{equation}
 (r,\phi,u_1,u_2)\longmapsto
 \bigl(-\cot(\pi r/\epsilon),-\cot(\phi/2),
       -\cot(\pi u_1),-\cot(\pi u_2)\bigr).
 \label{ret:eq:four-cell}
\end{equation}
This fixes the target orientation.  Replacing \(r\) locally by
\(\ell=\log r\) preserves this orientation, since \(dr/d\ell=r>0\).
The ordered coordinates in \eqref{ret:eq:four-cell} specify the sphere
identification used throughout this calculation.

\begin{theorem}
\label{ret:main}
For the exact map \eqref{ret:eq:H},
\begin{equation}
 [H]=\eta_4\circ\eta_5\ne0
       \quad\hbox{in}\quad\pi_6(S^4)\cong\mathbb Z/2,
 \label{ret:eq:main-class}
\end{equation}
where \(\eta_n:S^{n+1}\to S^n\) is the \((n-2)\)-fold reduced
suspension of the complex Hopf map.  In particular \(q_2\) is not
based null-homotopic.  Its stabilized top-cell postcomposition is the
element \(\eta^2\in\pi_2^{\rm S}\).
\end{theorem}

\subsection{The entire regular inverse image lies in one original chart}

Choose the actual lift
\begin{equation}
 u_0=(-1/4,1/4)^{\mathsf T},\qquad
 p_0=B_0u_0=(1/4,1/4)^{\mathsf T},\qquad \phi_0=\pi.
 \label{ret:eq:chosen-value}
\end{equation}
Modulo \(\mathbb Z^2\), \(u_0\) is \((3/4,1/4)\) in the open cell.  Put
\(t_0=r_0e^{i\phi_0}\) and
\(y_0=[r_0,c_3(e^{i\phi_0},[u_0])]\).  The numbers \(r_0\) and
\(t_0\) will remain nonzero throughout the framing calculation.

The original triangle \(\{0,e_1,e_2\}\) has toric chart
\begin{equation}
 \mathcal U_\sigma\cong\mathbb C^3,\qquad
 (z_0,z_1,z_2)=(t/(x_1x_2),x_1,x_2),\qquad
 t=z_0z_1z_2.
 \label{ret:eq:z-chart}
\end{equation}
These characters are dual to \((0,0,1),(1,0,1),(0,1,1)\), respectively,
as direct pairing verifies.  This is the chart of
\cite[Lemma 4.2(iii)]{retbib:Canonical}, with its original coordinates.

For \(j=1,2\), introduce the nowhere vanishing holomorphic functions
\begin{equation}
 h_j(t)=\exp\bigl(-2\pi i(C(t)u_0)_j\bigr).
 \label{ret:eq:units}
\end{equation}
The following is an exact coordinate map, with the exact inverse
displayed alongside it:
\begin{align}
 F(z)&=(w_0,w_1,w_2)
 =\bigl(z_0/(h_1(t)h_2(t)),\ h_1(t)z_1,\ h_2(t)z_2\bigr),
 \label{ret:eq:F}\\
 F^{-1}(w)&=
 \bigl(w_0h_1(t)h_2(t),\ w_1/h_1(t),\ w_2/h_2(t)\bigr),
 \quad t=w_0w_1w_2.
 \label{ret:eq:Finv}
\end{align}
Both maps are holomorphic wherever \(|t|\) is in the period disc,
both preserve \(t\), and substitution proves that both composites are
the identity.  In particular the maps extend across the origin and
across every coordinate divisor.  No period matrix is replaced by a
different matrix in this coordinate change.

\begin{lemma}
For sufficiently small \(r_0>0\), the whole inverse image
\(M=H^{-1}(y_0)\) is represented injectively inside one contractible
coordinate neighborhood in \(X\).  In its \(w\)-coordinates it is
exactly the embedded two-torus
\begin{equation}
 w(a,b)=
 \bigl(\rho_0e^{i(\phi_0-a-b)},\rho_1e^{ia},\rho_2e^{ib}\bigr),
 \qquad
 \rho_0=r_0^{1/2},\quad \rho_1=\rho_2=r_0^{1/4},
 \quad (a,b)\in\bigl(\mathbb R/(2\pi\mathbb Z)\bigr)^2.
 \label{ret:eq:M}
\end{equation}
\end{lemma}
\begin{proof}
The point \(y_0\) lies off the collapsed subset of the target, so
\eqref{ret:eq:H} and \eqref{ret:eq:q2} imply
\[
 M=\{x:t(f(x))=t_0,\ \overline\pi_f(x)=[u_0]\}.
\]
Any lift of such a point to \((\mathbb C^*)^2\times\{t_0\}\) has
\(A(t_0)^{-1}\log|x|=u_0+n\) for one \(n\in\mathbb Z^2\).
Applying the actual transformation \(\Psi_{-n}\) from
\eqref{ret:eq:deck} produces a lift with
\(\log|x|=A(t_0)u_0\).  If two such lifts represent the same point,
\(\Psi_n\) relates them and changes their lifted quotient coordinate
by \(n\); hence \(n=0\), and the lifts are identical.  This proves
both completeness and injectivity of this chosen representative.

Since
\(\log|h_j(t)|=2\pi\operatorname{Im}(C(t)u_0)_j
=-(R(t)u_0)_j\), its \(w\)-coordinates satisfy
\[
 (\log|w_1|,\log|w_2|)^{\mathsf T}
 =A(t_0)u_0-R(t_0)u_0=\ell_0p_0,
 \qquad \ell_0=\log r_0.
\]
Thus \(|w_1|=|w_2|=r_0^{1/4}\), and the unchanged product
\(w_0w_1w_2=t_0\) forces \(|w_0|=r_0^{1/2}\) and the displayed
phase.  Conversely every pair \((a,b)\) in \eqref{ret:eq:M} gives one
such lift.  The original kernel coordinates \(k_j\) had period one
in \(x_j=e^{2\pi i\zeta_j}\); the phase coordinates \(a,b\) have
period \(2\pi\), with no cover or change of lattice.

The covering \(Y_\epsilon\to N_0\) is injective on a neighborhood of
the origin in \eqref{ret:eq:z-chart}, because the original action is free
and properly discontinuous.  Choose a ball \(V\) centered at that
origin with closure inside such a neighborhood and the period disc.
Choose a ball \(V_w\) about the origin with closure contained in
\(F(V)\).  Every coordinate in \eqref{ret:eq:M} tends uniformly to zero
as \(r_0\downarrow0\).  Therefore for sufficiently small \(r_0\)
the entire torus lies in \(V_w\).  The inverse image of \(V_w\)
under \(F\), followed by the injective covering projection, is the
required contractible chart in \(X\).
\end{proof}

\subsection{The full differential and its ordered normal framing}

We calculate in a neighborhood of the compact torus \eqref{ret:eq:M}.
The quotient coordinates admit a unique local lift near \(u_0\).
Set
\[
 v=(\log|w_1|,\log|w_2|)^{\mathsf T},\quad
 \ell=\log|w_0w_1w_2|,\quad
 \phi=\arg(w_0w_1w_2)
\]
using the branch of \(\phi\) near \(\phi_0=\pi\).
Equations \eqref{ret:eq:logx} and \eqref{ret:eq:F} give the exact identity
\begin{equation}
 v=\ell B_0u+R(t)(u-u_0),\qquad
 u=u_0+A(t)^{-1}(v-\ell p_0).
 \label{ret:eq:exact-u}
\end{equation}
Differentiate this identity before restricting to the inverse image:
\begin{equation}
 du=-A(t)^{-1}(dA(t))A(t)^{-1}(v-\ell p_0)
        +A(t)^{-1}(dv-p_0\,d\ell).
 \label{ret:eq:full-derivative}
\end{equation}
On \(M\) the vector \(v-\ell p_0\) is exactly zero.  Therefore
\begin{equation}
 du|_M=A_0^{-1}(dv-p_0\,d\ell),\qquad A_0=A(t_0).
 \label{ret:eq:derivative-on-M}
\end{equation}
This explicitly accounts for all derivatives of the original bounded
period terms; none are omitted or estimated away.

Identify \(T_w\mathbb C^3\) with \(\mathbb C^3\) in the given coordinate chart,
and define real tangent vectors
\begin{equation}
 L_j=(0,\ldots,w_j,\ldots,0),\qquad
 Q_j=(0,\ldots,iw_j,\ldots,0)\qquad(0\leq j\leq2).
 \label{ret:eq:LQ}
\end{equation}
Thus \(d\log|w_i|(L_j)=\delta_{ij}\),
\(d\arg w_i(Q_j)=\delta_{ij}\), and the crossed evaluations vanish.
The ordered tangent vectors of the parametrization \eqref{ret:eq:M} are
\begin{equation}
 T_a=-Q_0+Q_1,\qquad T_b=-Q_0+Q_2.
 \label{ret:eq:tangents}
\end{equation}
They are independent because their \(w_1\) and \(w_2\) components,
respectively, are nonzero.  For the target order
\((\ell,\phi,u_1,u_2)\), define the following four lifts of its
coordinate vectors:
\begin{align}
 N_\ell&=\tfrac12L_0+\tfrac14L_1+\tfrac14L_2,
       &N_\phi&=Q_0,\label{ret:eq:normal-first}\\
 N_i&=-\bigl((A_0)_{1i}+(A_0)_{2i}\bigr)L_0
           +(A_0)_{1i}L_1+(A_0)_{2i}L_2
       &&(i=1,2).\label{ret:eq:normal-u}
\end{align}
Let \(P=(\ell,\phi,u_1,u_2)\).  Evaluation using
\eqref{ret:eq:derivative-on-M} gives the entire differential table
\begin{equation}
\begin{array}{c|rrrrrr}
 &N_\ell&N_\phi&N_1&N_2&T_a&T_b\\\hline
 d\ell&1&0&0&0&0&0\\
 d\phi&0&1&0&0&0&0\\
 du_1&0&0&1&0&0&0\\
 du_2&0&0&0&1&0&0
\end{array}.
 \label{ret:eq:differential-table}
\end{equation}
In particular \(P\), and therefore \(H\) in its open target cell,
is a submersion along \(M\).  The classes of
\((N_\ell,N_\phi,N_1,N_2)\) in \(TX|_M/TM\) give the specified
normal four-framing.  Orthogonal representatives are not required for
this quotient-bundle definition; projecting the displayed lifts along
\(TM\) to any chosen orthogonal complement preserves the quotient
frame.

If the normal frame is specified using the literal four real coordinates
in \eqref{ret:eq:four-cell}, their differential at the chosen value is
\[
 dH_{\rm cell}=K\,dP,\qquad
 K=\operatorname{diag}\!\left(
 \frac{\pi r_0}{\epsilon}\csc^2\frac{\pi r_0}{\epsilon},
 \frac12,2\pi,2\pi\right).
\]
The first factor includes \(dr=r_0\,d\ell\), the second uses
\(\phi_0=\pi\), and the final two use the actual representatives
\((3/4,1/4)\) of \(u_0\). Thus, writing
\(\mathcal N=(N_\ell,N_\phi,N_1,N_2)\), the literal dual frame is
\(\mathcal N K^{-1}\). All four entries of \(K\) are positive
constants on \(M\). The explicitly positive diagonal matrices
\((1-v)K^{-1}+vI_4\), \(0\leq v\leq1\), give a homotopy
from that literal frame to \(\mathcal N\). This retains the original
target coordinates and establishes that either coordinate description
gives the same framed bordism class.

For completeness, the original period matrix admits the following
explicit homotopy of this normal framing.  Put
\begin{equation}
 A_{0,v}=\ell_0B_0+vR(t_0),\qquad 0\leq v\leq1,
 \label{ret:eq:A-homotopy}
\end{equation}
and replace \(A_0\) by \(A_{0,v}\) only in
\eqref{ret:eq:normal-u}.  Equation \eqref{ret:eq:invertibility} proves
invertibility for every \(v\).  Their images under the actual
\((du_1,du_2)\) are the columns of \(A_0^{-1}A_{0,v}\), so these
are normal frames throughout the homotopy.  This is a homotopy of
the computed frame on the same embedded torus, not a replacement of
the original degeneration by a different one.

\begin{lemma}[The exact orientation sign]
Use the complex real-coordinate orientation, whose ordered coordinates are
\[
 (\operatorname{Re}w_0,\operatorname{Im}w_0,
  \operatorname{Re}w_1,\operatorname{Im}w_1,
  \operatorname{Re}w_2,\operatorname{Im}w_2).
\]
Then one has
\begin{equation}
 \det(N_\ell,N_\phi,N_1,N_2,T_a,T_b)
       =-\det(A_0)|w_0w_1w_2|^2<0.
 \label{ret:eq:sign}
\end{equation}
Hence the tangent orientation induced by the ordered normal frame is
the orientation of \((-T_a,T_b)\).
\end{lemma}
\begin{proof}
The ordered frame \((L_0,Q_0,L_1,Q_1,L_2,Q_2)\) has determinant
\(|w_0|^2|w_1|^2|w_2|^2>0\) in real coordinates.  Regroup its rows
as \((L_0,L_1,L_2,Q_0,Q_1,Q_2)\); this permutation has three
inversions and sign \(-1\).  Regroup the columns of the frame in
\eqref{ret:eq:sign} as
\((N_\ell,N_1,N_2,N_\phi,T_a,T_b)\); this has two inversions and
sign \(+1\).  The angular block is
\(\left(\begin{smallmatrix}1&-1&-1\\0&1&0\\0&0&1\end{smallmatrix}\right)\),
of determinant one.  The radial block has first column
\((1/2,1/4,1/4)^{\mathsf T}\) and the other columns given by
\eqref{ret:eq:normal-u}.  Adding its last two rows to its first row gives
first row \((1,0,0)\), with remaining lower-right block \(A_0\).
Its determinant is \(\det A_0\).  This proves the equality.
Equation \eqref{ret:eq:A-homotopy} joins \(A_0\) to \(\ell_0B_0\)
through invertible matrices, so
\(\operatorname{sign}\det A_0=\operatorname{sign}\det(\ell_0B_0)=+1\).
The inequality and the tangent orientation follow.
\end{proof}

\subsection{An explicit contraction of the full frame}

Define a complex-linear matrix, also regarded as a real-linear matrix,
\begin{equation}
 D(a,b)=\operatorname{diag}(e^{-i(a+b)},e^{ia},e^{ib})\in\operatorname{SU}(3).
 \label{ret:eq:D}
\end{equation}
By \eqref{ret:eq:M}, every \(L_j\) and \(Q_j\) at \((a,b)\) is
\(D(a,b)\) times its value at \((0,0)\).  The coefficients in
\eqref{ret:eq:normal-first}, \eqref{ret:eq:normal-u}, and
\eqref{ret:eq:tangents} are constant on \(M\), since \(t=t_0\).
It follows that the positively oriented full frame is exactly
\begin{equation}
 (N_\ell,N_\phi,N_1,N_2,-T_a,T_b)_{(a,b)}
       =D(a,b)G_0,
 \label{ret:eq:full-frame}
\end{equation}
where \(G_0\in\operatorname{GL}^+(6,\mathbb R)\) is its single value at \((0,0)\).
Thus neither an unspecified phase nor a position-dependent
real-linear matrix remains.

Here is a periodic contraction of \(D\), with all parameters
specified.  For \(v\in[0,1]\) put
\[
 s_v=\sin(\pi v/2),\qquad c_v=\cos(\pi v/2),\qquad
 Q_v(\theta)=
 \begin{pmatrix}s_ve^{-i\theta}&-c_v\\c_v&s_ve^{i\theta}\end{pmatrix},
 \qquad K_v(\theta)=Q_v(\theta)Q_v(0)^*.
\]
The two columns of \(Q_v\) have norm one and zero Hermitian inner
product, and \(\det Q_v=s_v^2+c_v^2=1\).  Thus
\(Q_v,K_v\in\operatorname{SU}(2)\).  Direct substitution gives
\[
 K_0(\theta)=I,\quad
 K_1(\theta)=\operatorname{diag}(e^{-i\theta},e^{i\theta}),\quad
 K_v(0)=I.
\]
Let \(\iota_{01},\iota_{02}:\operatorname{SU}(2)\to\operatorname{SU}(3)\) act in the
ordered coordinate planes \((0,1)\) and \((0,2)\), respectively,
and as the identity in the unused coordinate.  Then
\begin{equation}
 D_v(a,b)=\iota_{01}(K_v(a))\,\iota_{02}(K_v(b))
 \label{ret:eq:D-homotopy}
\end{equation}
is a well-defined continuous map on
\([0,1]\times(\mathbb R/(2\pi\mathbb Z))^2\), with
\(D_0=I\) and \(D_1=D\).  This proves the needed null-homotopy
directly, including both phase seams.  In particular
\(D_v(a,b)G_0\) contracts the comparison matrix in
\eqref{ret:eq:full-frame} to one constant positive matrix.

\begin{proposition}
The stable tangential framing induced by the actual normal four-frame
of \(H^{-1}(y_0)\) agrees, up to homotopy and the orientation sign
\eqref{ret:eq:sign}, with the Lie framing
\((\partial_a,\partial_b)\) on the native two-torus.
\end{proposition}
\begin{proof}
In the coordinate chart, the stable tangent trivialization is the
isomorphism
\[
 \mathbb R^4\oplus TM\longrightarrow T\mathbb C^3|_M,
 \qquad (v,V)\longmapsto
 v_1N_\ell+v_2N_\phi+v_3N_1+v_4N_2+V.
\]
Using the intrinsic frame \((-\partial_a,\partial_b)\) on \(TM\),
its matrix is precisely \eqref{ret:eq:full-frame}.  The homotopy
\eqref{ret:eq:D-homotopy} therefore makes this comparison constant.
A constant positive matrix changes no homotopy class of stable
framing: Gram--Schmidt deforms it to an element of \(SO(6)\), and
an orthogonal matrix of determinant one can be reduced to the identity
by finitely many oriented plane rotations, each joined to the
identity by varying its angle.  This proves the asserted comparison
within the chart.

It remains to compare the chart with the ambient sphere framing.
Pull back the usual stable tangent trivialization
\(TS^6\oplus\mathbb R\cong\mathbb R^7\) under the fixed based diffeomorphism
\(X\cong S^6\).  On the contractible \(w\)-ball containing \(M\),
it differs from the coordinate trivialization plus one real line by
a continuous map into \(\operatorname{GL}(7,\mathbb R)\).  Contracting the ball to its
center contracts this comparison map to a constant.  Thus its
restriction to \(M\) contributes no varying stable framing class.
If the chosen sphere orientation disagrees with the complex
orientation, the remaining constant has negative determinant and
reverses the resulting framed bordism class.  Both that sign and the
explicit sign \eqref{ret:eq:sign} disappear in the order-two class
computed below.  This argument compares actual ambient
trivializations on a contractible chart; it does not suppose that
\(\mathbb C^3\) globally trivializes \(TX\).
\end{proof}

\subsection{Nonzero framed bordism and the exact unstable class}

We give the nonvanishing test for the resulting torus explicitly.
A stable framing gives a spin structure by lifting its oriented frame
to the trivial \(\operatorname{Spin}\)-bundle.  On the Lie-framed
circle this is the disconnected double cover, whereas the spin
structure induced from the boundary of a framed disc is the
connected double cover: as the boundary tangent turns once, its
\(2\pi\)-rotation lifts from \(1\) to \(-1\) in
\(\operatorname{Spin}(2)\).  We call the former nonbounding and the
latter bounding.

For a closed oriented spin surface \(S\), the spin quadratic
refinement
\[
 q_S:H_1(S;\mathbb F_2)\longrightarrow\mathbb F_2
\]
takes a class represented by an embedded circle to zero or one
according to whether the induced spin structure, using its normal
line in the surface, is bounding or nonbounding.  It satisfies
\begin{equation}
 q_S(x+y)=q_S(x)+q_S(y)+(x\mathbin{\cdot}y).
 \label{ret:eq:quadratic}
\end{equation}
This is the usual spin quadratic form: resolving transverse
intersections adds one tangent turn per intersection, and the
\(2\pi\)-spin-lift calculation above gives the last term.  The
well-defined correspondence between spin structures and such
quadratic refinements is given in
\cite[Theorem 3.2 and Remark 3.3]{retbib:DebrayGunningham}; its original
surface proof is \cite{retbib:Johnson}.

For a symplectic basis \((e_i,f_i)_{i=1}^g\), put
\begin{equation}
 \operatorname{Arf}(S)=\sum_{i=1}^gq_S(e_i)q_S(f_i)\in\mathbb F_2.
 \label{ret:eq:Arf}
\end{equation}
The basis independence can be seen without selecting a preferred
basis: factor the finite sum over the symplectic planes to obtain
\[
 \sum_{x\in H_1(S;\mathbb F_2)}(-1)^{q_S(x)}
   =\prod_{i=1}^g
     \sum_{a,b\in\mathbb F_2}(-1)^{a q_S(e_i)+b q_S(f_i)+ab}
   =2^g(-1)^{\sum_i q_S(e_i)q_S(f_i)}.
\]
The left side has no choice of basis, proving the assertion.

Here is also the bordism-invariance argument needed for the present
obstruction.  A spin bordism has a handle decomposition.  An index
zero or three handle creates or removes a spin sphere, whose first
homology and Arf invariant are zero.  An index one handle either
connects two components, in which case the first-homology quadratic
forms add orthogonally, or increases the genus of one component.
In the latter case the new symplectic pair contains the meridian
of the handle.  That meridian bounds the cocore disc with its
induced spin structure, so its quadratic value is zero; the pair
contributes zero to \eqref{ret:eq:Arf}, regardless of the value on its
dual.  Existing cycles can be represented outside the attaching
discs, so their quadratic values are preserved.  An index two
handle is the reverse of an index one handle.  Therefore each
handle preserves the total Arf invariant of the boundary, and a
spin boundary has Arf invariant zero.  This also proves that a
stably framed boundary has Arf invariant zero, because forgetting
the framing gives a spin bordism.  The full spin bordism statement
is recorded in \cite[Theorem 3.6]{retbib:DebrayGunningham}.

For the intrinsic Lie framing on
\(M=(\mathbb R/(2\pi\mathbb Z))^2\), take
\[
 e=[a\mapsto(a,0)],\qquad f=[b\mapsto(0,b)].
\]
Both the tangent and normal vectors in the surface are constant in
the Lie frame along either circle.  Thus both circles have the
nonbounding spin structure, and
\begin{equation}
 e\mathbin{\cdot}f=1,\qquad q_M(e)=q_M(f)=1,
 \qquad q_M(xe+yf)=x+xy+y,
 \qquad \operatorname{Arf}(M)=1.
 \label{ret:eq:actual-Arf}
\end{equation}
Reversing the first Lie tangent vector changes the orientation but
does not change these mod-two values.  In particular the
specifically framed inverse image computed above is not a framed
boundary.

To identify the element, let \(\eta\in\pi_1^{\rm S}\) denote the
stable class of the complex Hopf map.  Its Pontryagin--Thom inverse
image is the circle with its Lie framing.  This can be verified
directly on the Hopf map \(S^3\to\mathbb{CP}^1\): near the fibre
\((z_1,z_2)=(e^{ia},0)\), the target coordinate is \(z_2/z_1\), so
the two horizontal normal vectors are \((0,e^{ia})\) and
\((0,ie^{ia})\).  Add the tangent vector \((ie^{ia},0)\) and the
outward radial vector \((e^{ia},0)\).  In \(\mathbb C^2\) the resulting
full frame makes one complete rotation in each of its two real
coordinate planes.  This loop is contractible explicitly: for the
fixed real-linear involution \(J(z_1,z_2)=(z_1,\overline z_2)\),
\[
 J\operatorname{diag}(e^{ia},e^{ia})J=\operatorname{diag}(e^{ia},e^{-ia})=K_1(-a),
\]
and \(JK_v(-a)J\) contracts the original loop to the identity using
the periodic homotopy already constructed.  Thus the induced stable
tangent frame on the circle is the Lie frame.  This is the basic
framed-circle identification in \cite{retbib:Pontryagin1955}.

The product of two normally framed inverse images represents the
product in stable homotopy: its tubular collapse is the smash of
the two tubular collapse maps, as follows immediately from the
product tubular neighborhood \(\nu_1\times\nu_2\) and its ordered
normal coordinates.  Consequently the product of the two Lie
circles represents \(\eta^2\).  Equation \eqref{ret:eq:actual-Arf}
proves \(\eta^2\ne0\) by an explicit bordism invariant.  The
stable group \(\pi_2^{\rm S}=\mathbb Z/2\) is Pontryagin's computation
\cite{retbib:Pontryagin1950,retbib:Pontryagin1955}; the stable table in
\cite[pp.~384--385]{retbib:Hatcher} gives the same group.  Thus the sign from
\eqref{ret:eq:sign} satisfies \(-\eta^2=\eta^2\).

\begin{proof}[Proof of the theorem]
The globally continuous quotient \(H\) is smooth on a neighborhood
of its compact inverse image \(M\), by
\eqref{ret:eq:exact-u} and \eqref{ret:eq:four-cell}. First make it constant
near the source basepoint, before choosing any approximation bound.
Choose a target chart \(\kappa:O\to\mathbb R^4\) about \(*\) with
\(\kappa(*)=0\), convex ball image, and \(y_0\notin O\).
Choose a source neighborhood \(U\) of \(x_{\rm base}\), disjoint
from \(M\), with \(H(U)\subset O\). Let \(\beta\) be a smooth
function supported in \(U\), equal to one near \(x_{\rm base}\),
with \(0\leq\beta\leq1\). On \(U\) use the based homotopy
\[
 H_s(x)=\kappa^{-1}\bigl((1-s\beta(x))\kappa(H(x))\bigr),
 \qquad 0\leq s\leq1,
\]
and leave \(H\) unchanged off \(U\). Its endpoint \(H_b\) is
smooth and constant near the basepoint; it agrees with \(H\)
near \(M\) and has no new inverse image of \(y_0\).

Choose one closed tubular neighborhood \(U_1\) of \(M\), with
\(M\subset\operatorname{int}U_1\), contained in the unchanged
smooth region and disjoint from the basepoint neighborhood.
Compactness gives the positive bound
\[
 d_0=\min_{x\in X\setminus\operatorname{int}U_1}
             d\bigl(H_b(x),y_0\bigr)>0.
\]
Relative smooth approximation, fixing the same closed set \(U_1\)
and a neighborhood of the basepoint, produces a smooth based map
\(H_{\rm sm}\) whose error from \(H_b\) is less than \(d_0/2\)
on its complement. Choose the approximation closer also than a
fixed target injectivity radius; the short target geodesics give a
based homotopy. Thus the exact inverse image and its differential
are unchanged on \(U_1\), and the complement acquires no new
inverse image. The relative Pontryagin construction applies as in
\cite[\S7, Theorem B and Lemma 4, pp.~43, 48--50]{retbib:Milnor1965}.

Pontryagin--Thom therefore identifies the stabilization of
\([H]\) with the bordism class of \(M\) carrying the four normal
vectors \eqref{ret:eq:normal-first}--\eqref{ret:eq:normal-u}.  The comparison
\eqref{ret:eq:full-frame} and its explicit contraction identify that
class with the Lie torus, with the harmless order-two orientation
sign already computed.  Hence
\[
 E^\infty[H]=\eta^2\ne0.
\]
Freudenthal suspension is an isomorphism
\[
 E:\pi_{n+2}(S^n)\longrightarrow\pi_{n+3}(S^{n+1})
 \qquad(n\geq4),
\]
because \(n+2\leq2n-2\).  Therefore stabilization from
\(\pi_6(S^4)\) is an isomorphism onto \(\pi_2^{\rm S}\).
The composite \(\eta_4\circ\eta_5:S^6\to S^4\) stabilizes to
\(\eta^2\), so injectivity gives the exact equality
\eqref{ret:eq:main-class}.  Finally, a based null-homotopy of \(q_2\)
would postcompose with \(p_3\) to a based null-homotopy of \(H\),
contradicting \eqref{ret:eq:main-class}.
\end{proof}

\paragraph{Exact scope of the comparison with the angle-forgetting map.}
The map here retains the angular coordinate before collapsing the
coordinate two-skeleton.  Its regular inverse image is the specific
normally four-framed \(T^2\) above.  The angle-forgetting
postcomposition has a normally three-framed \(T^3\) inverse image.
The calculation above is independent of its class and supplies an
actual obstruction to null-homotopy of the original \(q_2\), even
when that other postcomposition is null-homotopic.  The exact
morphism carrying this obstruction is
\[
 (p_3)_*:[X,\mathcal A]_*\longrightarrow[X,S^4]_*,\qquad
 [q_2]\longmapsto[H]=\eta_4\eta_5.
\]
No conclusion about an arbitrary earlier quotient follows merely
from the diffeomorphism type of one inverse image; the specified
four-framing is what proves this value.

\begin{corollary}[The exact order-two quotient detected by the original cusp map]
\label{ret:q2-detected-quotient}
Let \(d:(S^6,*)\to(X,x_{\mathrm{base}})\) be the fixed based
diffeomorphism, and put
\[
\alpha=[q_2\circ d]\in\pi_6(\mathcal A,v_0).
\]
The original top-cell postcomposition induces a surjective homomorphism
\[
\varepsilon=(p_3)_*:\pi_6(\mathcal A,v_0)
  \longrightarrow\pi_6(S^4,*)\cong\mathbb Z/2,
\qquad \varepsilon(\alpha)=\eta_4\circ\eta_5.
\]
It gives the exact quotient isomorphism
\[
\pi_6(\mathcal A,v_0)/\ker\varepsilon
 \xrightarrow{\ [\beta]\mapsto\varepsilon(\beta)\ }
\mathbb Z/2.
\]
Every odd multiple of \(\alpha\) is nonzero. Its order is either
infinite or an even positive integer; this calculation does not determine
that order or assert a splitting of the displayed quotient.
\end{corollary}
\begin{proof}
Postcomposition distributes over the pinch sum on the source sphere,
so \((p_3)_*\) is a group homomorphism. Theorem~\ref{ret:main}
gives \((p_3)_*(\alpha)=[p_3q_2d]=\eta_4\eta_5\ne0\).
This is the nonzero element of the cyclic group of order two, proving
surjectivity. The first isomorphism theorem gives the stated quotient
with exactly the displayed map. For every integer \(n\),
\[
\varepsilon(n\alpha)=(n\bmod2)\eta_4\eta_5.
\]
An odd multiple therefore has nonzero image and cannot vanish.
If \(m\alpha=0\), applying \(\varepsilon\) forces \(m\)
to be even. These arguments prove all the assertions without inferring
the exact order of \(\alpha\).
\end{proof}

\section{A nonzero Hopf--cusp morphism to the stable mod--24 group}
\label{hbridge:section}
The retained-angle input used below is the complete Theorem~\ref{ret:main} above.

\subsection{The exact source maps and conventions}

Retain the original marked source \((X,x_0)\), its fixed based
diffeomorphism \(d:(S^6,*)\to(X,x_0)\), the cusp coordinate \(t\),
the fixed cutoff \(\epsilon\), and the native quotient
\(\overline\pi_f\) of the cusp fibres by
\(\mathbb Z\langle\widehat w,\widehat\delta\rangle\).
The quotient coordinates are ordered by
\((\overline{\widehat\gamma},\overline{\widehat u})\).
Thus the fixed continuous based map is
\[
q_2:X\longrightarrow
\mathcal A=\operatorname{Susp}(S^1\times\overline T^2),
\qquad
q_2(x)=[|t|,t/|t|,\overline\pi_f(x)]
\quad(0<|t|<\epsilon).
\]
It sends the entire central fibre \(W\) to the first suspension
vertex and \(X\setminus\operatorname{int}N_0\) to the second.
The unreduced suspension \(\mathcal A\) is based at its first
vertex. In the original product torus set
\[
K^{(2)}=(\{1\}\times\overline T^2)
\cup(S^1\times\{u_1=0\})
\cup(S^1\times\{u_2=0\}),
\qquad c_3:T^3\longrightarrow T^3/K^{(2)}\cong S^3.
\]
The retained-angle sphere map and its sphere-source representative are
\begin{equation}
p_3([r,z])=[r,c_3(z)],\qquad
H=p_3q_2:X\longrightarrow\Sigma_r(S^3,*)\cong S^4,
\qquad h=H\circ d:S^6\longrightarrow S^4.
\label{hbridge:eq:retained-map}
\end{equation}
Both suspension vertices and the whole cylinder over the basepoint
are identified in \(\Sigma_r(S^3,*)\). The preceding complete
retained-angle framing theorem proves
\begin{equation}
[h]=\eta_4\circ\eta_5\ne0
\quad\hbox{in}\quad\pi_6(S^4)\cong\mathbb Z/2.
\label{hbridge:eq:proved-H}
\end{equation}
Its proof identifies the actual derivative-framed inverse torus
with the Lie-framed torus, with Arf invariant one. We use this
already proved class, retaining the particular map
\eqref{hbridge:eq:retained-map}; see \cite{hbridgebib:HBRetained}.

Fix the complex Hopf map
\[
\eta_2:(S^3,(1,0))\longrightarrow
(\mathbb{CP}^1,[1:0])\cong(S^2,*),\qquad
(z_1,z_2)\longmapsto[z_1:z_2].
\]
Use ordered smash coordinates on spheres. In this note
\(E^k f=f\wedge\operatorname{id}_{S^k}\) denotes right reduced
suspension; the usual ordered sphere homeomorphisms are understood.
Put \(\eta_n=E^{n-2}\eta_2\), and write \(\eta\) for the
stable Hopf class. The order of each factor is retained below.

\subsection{The Hopf--cusp homomorphism and its nonzero value}

\begin{theorem}
\label{hbridge:mu-theorem}
The assignment
\begin{equation}
\mu_\eta:\pi_6(S^4)\longrightarrow\pi_9(S^6),\qquad
\mu_\eta([x])=[\eta_2\wedge x]
\label{hbridge:eq:mu}
\end{equation}
is an injective group homomorphism. Its literal representative has
the complete domain and codomain
\[
\eta_2\wedge x:S^3\wedge S^6\cong S^9
\longrightarrow S^2\wedge S^4\cong S^6.
\]
Choose any generator \(\nu_6\) of
\(\pi_9(S^6)\cong\mathbb Z/24\). Then the particular original
cusp map supplies
\begin{equation}
\kappa_H=[\eta_2\wedge h]
=\mu_\eta([h])=12\nu_6\ne0.
\label{hbridge:eq:kappa}
\end{equation}
It has exact order two. Its value is independent of orientation
reversals in the sphere identifications and of the choice of the
generator \(\nu_6\).
\end{theorem}

\begin{proof}
A based homotopy of \(x\) smashes with the fixed based Hopf map,
so \eqref{hbridge:eq:mu} is well-defined. The exact ordered
point-set factorization is
\begin{equation}
\eta_2\wedge x
=(\eta_2\wedge\operatorname{id}_{S^4})
 \circ(\operatorname{id}_{S^3}\wedge x):
S^3\wedge S^6\longrightarrow S^3\wedge S^4
\longrightarrow S^2\wedge S^4.
\label{hbridge:eq:factorization}
\end{equation}
The first applied map has source dimension nine and target dimension
seven; the second has source dimension seven and target dimension six.
The second map is \(E^4\eta_2\). To compare the first with
\(E^3x=x\wedge\operatorname{id}_{S^3}\), interchange its sphere
factors at source and target. The interchange degrees are respectively
\((-1)^{3\cdot6}=+1\) and \((-1)^{3\cdot4}=+1\). A based map
of a sphere of degree one is based homotopic to the identity, so
\begin{equation}
\mu_\eta([x])=(E^4\eta_2)_*(E^3[x]).
\label{hbridge:eq:hom-formula}
\end{equation}
Suspension is a group homomorphism. Postcomposition is also a group
homomorphism on homotopy groups, since it distributes over the
source pinch map defining the sum. Thus
\eqref{hbridge:eq:hom-formula} proves additivity, including zero
and inverses, before passage to stable homotopy.

Freudenthal makes
\[
E:\pi_{n+3}(S^n)\longrightarrow\pi_{n+4}(S^{n+1})
\]
an isomorphism for \(n\ge5\), because \(n+3<2n-1\).
In particular stabilization
\(\pi_9(S^6)\to\pi_3^{\mathrm S}\) is an isomorphism.
Stabilization from \(\pi_6(S^4)\) is likewise an isomorphism
onto \(\pi_2^{\mathrm S}\), because \(n+2<2n-1\) for
\(n\ge4\). The exact factorization
\eqref{hbridge:eq:factorization} shows that its stable class is
the stable composition product of \(\eta\) with \([x]\).
Using the proved original class \eqref{hbridge:eq:proved-H},
\begin{equation}
E^\infty\kappa_H
=\eta\,E^\infty[h]
=\eta\,\eta^2
=\eta^3
=12\nu\ne0\quad\hbox{in}\quad
\pi_3^{\mathrm S}\cong\mathbb Z/24.
\label{hbridge:eq:stable-product}
\end{equation}
The integral relation here is the one explicitly stated in
\cite[p.~385]{hbridgebib:HBHatcher}, not the coefficient \(4\) in the
two-primary quotient. More precisely, the two-primary relation
is \(\eta^3=4\nu\) modulo eight. Bilinearity and \(2\eta=0\)
give \(2\eta^3=0\) integrally. In \(\mathbb Z/24\) the only
elements killed by two are \(0\) and \(12\nu\), and only the
second has nonzero image \(4\nu\) modulo eight. This recovers
the exact integral coefficient and its nonvanishing.

The inverse of the stabilization isomorphism therefore gives
\eqref{hbridge:eq:kappa}. Since the source of \(\mu_\eta\) is
cyclic of order two, generated by \([h]\), its nonzero value
proves injectivity and identifies its image as
\(\{0,12\nu_6\}\).

An orientation reversal at a sphere source or target multiplies
the corresponding stable homotopy class by \(-1\); here
\(-12\nu_6=12\nu_6\). A sign change in the chosen Hopf
representative also leaves the stable product unchanged since
\(-\eta=\eta\). Finally, replacing the generator by
\(u\nu_6\), with \(u\) a unit modulo twenty-four, does not
change the element: every such \(u\) is odd, and
\(12u\equiv12\pmod{24}\). All orientation and generator
assertions consequently hold for the actual order-two element.
\end{proof}

\begin{corollary}[The same nonzero class in the twelve-dimensional receiver]
Put \(\nu_9=E^3\nu_6\). The particular based map
\[
E^3(\eta_2\wedge h):S^{12}\longrightarrow S^9
\]
has class \(12\nu_9\ne0\) in
\(\pi_{12}(S^9)\cong\mathbb Z/24\).
\end{corollary}
\begin{proof}
Each of these three suspensions is an isomorphism in the preceding
stable range. Apply their composite to
\eqref{hbridge:eq:kappa}; it preserves the exact order two.
\end{proof}

\subsection{The actual marked peripheral character}

Retain the original three elliptic-kernel matrices, with their
original signs and ordering:
\begin{equation}
C_0=\begin{pmatrix}1&1\\0&1\end{pmatrix},\qquad
C_1=\begin{pmatrix}-1&-1\\1&0\end{pmatrix},\qquad
C_2'=\begin{pmatrix}0&1\\-1&0\end{pmatrix}.
\label{hbridge:eq:matrices}
\end{equation}
The existing peripheral theorem \cite{hbridgebib:HBHigher} defines the
surjective marked abelianization character
\begin{equation}
\chi:SL_2(\mathbb Z)\longrightarrow\mathbb Z/12,
\qquad
(\chi(C_0),\chi(C_1),\chi(C_2'))=(-1,4,-3).
\label{hbridge:eq:chi}
\end{equation}
For an explicit check of this normalization, let
\[
S_0=\begin{pmatrix}0&-1\\1&0\end{pmatrix},\qquad T=C_0.
\]
Then \(C_2'=S_0^3\), \(C_1=T^{-1}S_0\), and direct matrix
multiplication gives \(C_1C_2'=C_0^{-1}\). In the modular
presentation \(S_0^4=1\), \(S_0^2=(S_0T)^3\), the
abelianized generators satisfy \(4s=0\) and \(s+3t=0\).
Thus \(s=-3t\) and \(12t=0\). The assignment
\(t\mapsto-1\), \(s\mapsto3\) satisfies both relations and
maps onto \(\mathbb Z/12\); consequently the order of the
abelianized generator is exactly twelve. It gives, without a
change of peripheral matrices, precisely
\(\chi(C_0)=-1\), \(\chi(C_1)=1+3=4\), and
\(\chi(C_2')=3\cdot3=-3\pmod{12}\).

\begin{theorem}[An exact nonzero peripheral homotopy morphism]
\label{hbridge:peripheral-theorem}
Define
\begin{equation}
\Theta_H:\mathbb Z/12\longrightarrow\pi_6(S^4),
\qquad \Theta_H([n])=n[h],
\qquad
\Psi_H=\mu_\eta\circ\Theta_H\circ\chi:
SL_2(\mathbb Z)\longrightarrow\pi_9(S^6).
\label{hbridge:eq:Psi}
\end{equation}
These are group homomorphisms. The complete coefficient formula is
\begin{equation}
\Psi_H(M)=12\,\widetilde{\chi(M)}\,\nu_6,
\label{hbridge:eq:coefficient}
\end{equation}
where \(\widetilde{\chi(M)}\) is any integer lift of the
marked residue. In the original matrix order of
\eqref{hbridge:eq:matrices}, its exact values are
\begin{equation}
(\Psi_H(C_0),\Psi_H(C_1),\Psi_H(C_2'))
=(12\nu_6,0,12\nu_6).
\label{hbridge:eq:values}
\end{equation}
Its image and kernel are exactly
\begin{equation}
\operatorname{im}\Psi_H=\{0,12\nu_6\},\qquad
\ker\Psi_H=\chi^{-1}(2\mathbb Z/12).
\label{hbridge:eq:kernel}
\end{equation}
In particular its kernel has index two in \(SL_2(\mathbb Z)\).
\end{theorem}

\begin{proof}
The element \([h]\) has exact order two, so replacing \(n\)
by \(n+12k\) leaves \(n[h]\) unchanged. Addition in its
abelian homotopy group proves that \(\Theta_H\) is a group
homomorphism. Moreover
\[
\ker\Theta_H=2\mathbb Z/12,
\qquad \operatorname{im}\Theta_H=\pi_6(S^4).
\]
The composite in \eqref{hbridge:eq:Psi} is consequently a
group homomorphism, and \eqref{hbridge:eq:kappa} gives
\(\mu_\eta(n[h])=12n\nu_6\). Changing the integer lift by
\(12k\) changes its coefficient by \(144k\), which is a
multiple of twenty-four. This proves
\eqref{hbridge:eq:coefficient} with its independence of lifts.
The three coefficients are respectively
\(-12\equiv12\), \(48\equiv0\), and
\(-36\equiv12\pmod{24}\), proving
\eqref{hbridge:eq:values}. Since \(12n\equiv0\pmod{24}\)
holds exactly when \(n\) is even, the asserted kernel is exact.
Surjectivity of \(\chi\) and nonvanishing at \(C_0\) prove
that the image is the full order-two subgroup, so the first
isomorphism theorem gives the index assertion.
\end{proof}

Here is the complete parity factorization. Let
\(\operatorname{red}_2:\mathbb Z/12\to\mathbb Z/2\) be
ordinary parity reduction, let
\(\iota_2([e])=12e\pmod{24}\), and let
\(c_{\nu_6}:\pi_9(S^6)\to\mathbb Z/24\) send \(\nu_6\)
to one. Every arrow and its codomain in the diagram are fixed:
\begin{equation}
\begin{CD}
SL_2(\mathbb Z) @>{\Psi_H}>> \pi_9(S^6)\\
@V{\operatorname{red}_2\circ\chi}VV @VV{c_{\nu_6}}V\\
\mathbb Z/2 @>{\iota_2}>> \mathbb Z/24.
\end{CD}
\label{hbridge:eq:parity-square}
\end{equation}
Commutativity is exactly the coefficient calculation above; the
bottom arrow is injective because \(12e\) is zero modulo
twenty-four exactly when \(e\) is zero modulo two.

\subsection{Comparison with the earlier maps, with their exact types}

Write \(\psi_H=c_{\nu_6}\circ\Psi_H\). Ordinary reduction
\(r:\mathbb Z/24\to\mathbb Z/12\) gives
\begin{equation}
r\circ\psi_H=0,
\qquad (r\circ\psi_H)(C_0)=0\ne-1=\chi(C_0).
\label{hbridge:eq:not-lift}
\end{equation}
Thus this nonzero character is not a reduction-compatible lift of
\(\chi\). Its natural commuting square is
\eqref{hbridge:eq:parity-square}. The existing nonexistence
theorem for a reduction-compatible lift remains valid: every
homomorphism from \(SL_2(\mathbb Z)\) to an abelian group
factors through its abelianization \(\mathbb Z/12\), whereas
a lift of \(\chi(C_0)=-1\) would require \(C_0\) to map
to \(11\) or \(23\) modulo twenty-four, each of order
twenty-four. The abelianized generator has order twelve and
cannot have either image.

The universal even-subgroup injection is the separate map
\[
j:\mathbb Z/12\hookrightarrow\mathbb Z/24,
\qquad j([a])=[2a].
\]
The new coefficient morphism factors through it by the exact
identity
\begin{equation}
\psi_H=j\circ[6]\circ\chi,
\qquad [6]:\mathbb Z/12\longrightarrow\mathbb Z/12,
\quad [a]\longmapsto[6a].
\label{hbridge:eq:j-factor}
\end{equation}
Indeed \(2(6a)=12a\pmod{24}\). The map \(j\) is
injective with image of order twelve; the composite
\eqref{hbridge:eq:j-factor} has image of order two because
\([6]\) retains precisely parity. The map \(j\) is realized
by the already proved universal suspension
\(E^6:\pi_6(S^3)\to\pi_{12}(S^9)\) after compatible
generator choices. Equation \eqref{hbridge:eq:j-factor} is a
coefficient identity; the actual new geometric representative is
\(\eta_2\wedge h\), followed, when desired, by three
suspensions as above.

For the original angle-forgetting map
\(g=\widehat p_2q_2:X\to S^3\), the separate framed
calculation has proved \([g\circ d]=0\). Therefore the
previous morphism
\[
\Theta_g:\mathbb Z/12\longrightarrow\pi_6(S^3),
\qquad [n]\longmapsto n[g\circ d]
\]
and \(E^6\Theta_g\chi\) are zero. In contrast,
\(E^3\Psi_H(C_0)=12\nu_9\ne0\) in this same stable
receiver \(\pi_{12}(S^9)\). These statements agree because
the exact postcompositions from \(\mathcal A\) are different:
\(p_3\) retains the angular coordinate in its top-cell map,
while \(\widehat p_2\) forgets it. The new homomorphism
uses the already proved value of \(p_3q_2\) and the specified
Hopf map, with all maps displayed in
\eqref{hbridge:eq:retained-map}--\eqref{hbridge:eq:Psi}.

The constructed representative has source \(S^3\wedge X\cong
S^9\); its threefold suspension has source \(S^{12}\).
The number twenty-four is the order of the receiving stable
homotopy group. This calculation does not identify either source
with a twenty-four-dimensional carrier, does not supply new
periods or analytic overlap maps, and does not complete the
global higher torus family. It supplies the exact nonzero
homotopy and peripheral morphisms just proved.

\paragraph{Source checks.}
The three matrices and the marked values were checked directly
against equations \emph{elliptic-peripheral-matrices} and
\emph{peripheral-mod12-values} in the existing higher-rung draft
\cite{hbridgebib:HBHigher}. The calculation above reproduces their signs
using the same \(S_0,T\). The retained-angle input is the
complete derivative-framing theorem in \cite{hbridgebib:HBRetained}.
For the stable facts, the author-hosted text
\cite{hbridgebib:HBHatcher} was read at Corollary~4.24, printed p.~360,
and at printed pp.~384--385. Page~385 explicitly distinguishes
the two-primary coefficient four from the integral coefficient
twelve. Proposition~4.56 on that page, with its proof on
p.~388, gives the stable composition product and its smash
description. The fixed-factor proof
\eqref{hbridge:eq:factorization}--\eqref{hbridge:eq:hom-formula}
establishes the needed homomorphism here directly.



% END RETAINED_ANGLE_AND_HOPF_BRIDGE_20260905

% BEGIN RETAINED_ANGLE_SUPPLEMENTARY_REFERENCES_20260905
\begingroup
\sloppy
\raggedright
\renewcommand{\refname}{Supplementary references for the retained-angle and Hopf--cusp proofs}
\begin{thebibliography}{99}

\bibitem[Alp\"oge(2026)]{retbib:Canonical}
L.~Alp\"oge, with the disclosed Fable/Claude contribution,
\emph{The \((3,4,\infty)\) modular family of \(2\)-tori, completed at
its three special points, is a complex structure on \(S^6\)},
canonical audit and reconstruction source bundle, Zenodo record 22238779
(2026), original manuscript \S3.7 and \S\S4.1--4.5.
\url{https://doi.org/10.5281/zenodo.22238779}.

\bibitem[Cusp sphere-route companion(2026)]{retbib:ConeRoutes}
\emph{Real \(S^{12}\) and \(S^{24}\) Routes from Alp\"oge's
Candidate \((3,4,\infty)\) Two-Torus Fibration},
companion research draft in the preceding canonical source bundle;
native quotient, localized cusp map, and framed-cusp calculation.
Current proof-bearing edition:
\url{https://www.overleaf.com/read/rtmyqxyrzprn}.

\bibitem[Milnor(1965)]{retbib:Milnor1965}
J.~W.~Milnor, \emph{Topology from the Differentiable Viewpoint},
University Press of Virginia, 1965, \S7, Theorem B and Lemma 4.
\url{https://www.maths.ed.ac.uk/~aar/papers/milnortop.pdf}.

\bibitem[Pontryagin(1950)]{retbib:Pontryagin1950}
L.~S.~Pontryagin,
\emph{Homotopy classification of mappings of an \((n+2)\)-dimensional
sphere on an \(n\)-dimensional one},
Dokl. Akad. Nauk SSSR \textbf{70} (1950), 957--959.
\url{https://www.sas.rochester.edu/mth/sites/doug-ravenel/otherpapers/pont3.pdf}.

\bibitem[Pontryagin(1955)]{retbib:Pontryagin1955}
L.~S.~Pontryagin,
\emph{Smooth manifolds and their applications in homotopy theory},
Trudy Mat. Inst. Steklov \textbf{45} (1955);
English translation, Amer. Math. Soc. Transl. (2) \textbf{11} (1959).
\url{https://www.sas.rochester.edu/mth/sites/doug-ravenel/otherpapers/pont4.pdf}.

\bibitem[Johnson(1980)]{retbib:Johnson}
D.~Johnson, \emph{Spin structures and quadratic forms on surfaces},
J. London Math. Soc. (2) \textbf{22} (1980), 365--373.

\bibitem[Debray and Gunningham(2018)]{retbib:DebrayGunningham}
A.~Debray and S.~Gunningham,
\emph{The Arf--Brown TQFT of pin-minus surfaces},
arXiv:1803.11183, especially \S3.1.1 and Proposition 5.1.
\url{https://adebray.github.io/papers/arf_brown.pdf}.

\bibitem[Hatcher(2002)]{retbib:Hatcher}
A.~Hatcher, \emph{Algebraic Topology}, Cambridge University Press,
2002; group table p.~339, Corollary 4.24, p.~360, and pp.~384--385.
\url{https://pi.math.cornell.edu/~hatcher/AT/AT.pdf}.

\bibitem[Retained-angle cusp theorem(2026)]{hbridgebib:HBRetained}
\emph{The retained-angle cusp quotient has nonzero top-cell class},
complete derivative-framing theorem and its based
Pontryagin--Thom proof, 5 September 2026;
companion to the canonical cusp/sphere-route draft.
Current cumulative source collection:
\url{https://www.overleaf.com/read/rtmyqxyrzprn}.

\bibitem[Higher torus wrap(2026)]{hbridgebib:HBHigher}
\emph{The higher torus wrap: obstruction, exact local lift, and the
Jordan--triality rung},
canonical higher-rung research draft, section on peripheral
modulus and the mod--24 boundary, 5 September 2026.
Current cumulative source collection:
\url{https://www.overleaf.com/read/rtmyqxyrzprn}.

\bibitem[Hatcher(2002)]{hbridgebib:HBHatcher}
A.~Hatcher, \emph{Algebraic Topology}, Cambridge University Press,
2002, Corollary~4.24, p.~360, stable homotopy groups and the
integral relation \(\eta^3=12\nu\), pp.~384--385,
and Proposition~4.56 with proof on p.~388.
\url{https://pi.math.cornell.edu/~hatcher/AT/AT.pdf}.

\end{thebibliography}
\endgroup
% END RETAINED_ANGLE_SUPPLEMENTARY_REFERENCES_20260905

\bibliographystyle{abbrvnat}
\bibliography{s6_cone_sphere_routes_companion}

\end{document}
